\documentclass[11pt]{article}

\usepackage[utf8]{inputenc}
\usepackage[T1]{fontenc}
\usepackage{lmodern}
\usepackage{booktabs}
\usepackage{array}
\usepackage{enumitem}
\usepackage{amsmath, amssymb, amsthm}
\usepackage[margin=1in]{geometry}
\usepackage{microtype}
\PassOptionsToPackage{hyphens}{url}
\usepackage{xcolor}
\usepackage{hyperref}
\hypersetup{
  colorlinks=true,
  linkcolor=blue!50!black,
  urlcolor=blue!50!black,
  citecolor=blue!50!black,
  pdftitle={The Universal Right-Triangle Reflection Sequence},
  pdfauthor={Vico Bonfioli},
}

\theoremstyle{plain}
\newtheorem{theorem}{Theorem}
\newtheorem{lemma}[theorem]{Lemma}
\newtheorem{corollary}[theorem]{Corollary}
\newtheorem{proposition}[theorem]{Proposition}
\newtheorem{conjecture}[theorem]{Conjecture}

\theoremstyle{definition}
\newtheorem{definition}[theorem]{Definition}
\newtheorem{example}[theorem]{Example}

\theoremstyle{remark}
\newtheorem{remark}[theorem]{Remark}

\newcommand{\Z}{\mathbb{Z}}
\newcommand{\Q}{\mathbb{Q}}
\newcommand{\R}{\mathbb{R}}
\newcommand{\C}{\mathbb{C}}
\newcommand{\Aff}{\mathrm{Aff}}
\DeclareMathOperator{\Isom}{Isom}

\title{The Universal Right-Triangle Reflection Sequence:\\
a word-length metric, effective universality,\\
and the lamplighter structure of \texorpdfstring{A396406}{A396406}}
\author{Vico Bonfioli\thanks{Independent researcher, Italy.
\texttt{vico@anvilstack.com}, \texttt{vicobonfioli@gmail.com}.}}
\date{\today}

\begin{document}
\maketitle

\begin{abstract}
For a right triangle $T$ in the plane with unequal rational legs, the
group $W(T)$ generated by the three side-reflections acts on the orbit
of $T$, and the orbit-growth sequence $u_d^T$ (the number of triangles
at word-distance $d$) is the subject of OEIS entry A396406. We prove
that $u_d^T$ is independent of the shape $T$ through depth $32$, and
that this is sharp: the $(1,2)$ triangle deviates at depth $33$. The
common sequence is the orbit growth of the \emph{generic} right
triangle, whose abstract reflection group is the group
$W=\langle R_0,R_1,R_2\mid R_i^2=(R_0R_1)^2=1\rangle$. We show that the
subgroup of $W$ on which the linear part is a rotation is the wreath
product $\Z\wr\Z$ of index $4$, so the generic group is virtually
$\Z\wr\Z$; the translation lattice is computed exactly as
$\mathcal T = 2(t-1)\,\Z[t^{\pm1}]$ via a Mayer--Vietoris/Crowell
identification of the translation module with the cyclic
$\Z[Q]$-ideal generated by $h=(c-1)(Y-1)$. From the resulting lamp
model we derive a closed-form word-length metric whose geodesic length
equals a relaxed length plus twice an explicit connectivity penalty.
A triangle-independent symbolic normal form reduces every extra
collision to the vanishing of a bounded-height integer Laurent
polynomial at the rotation number $\zeta_T=(a+bi)/(a-bi)$; Gauss's
lemma then gives an effective all-depths universality theorem (no
deviation at depth $d$ if $c_T:=N(a-bi$, suitably normalized$)>2d$),
and a finite modular certificate closes the gap to depth $32$. The
natural conjecture that universality holds at all depths is refuted:
for every algebraic shape we exhibit an explicit word in $\ker\rho_T$
(a product of conjugated glide-reflection squares), and the first
deviation depth $n_T$ is sandwiched by $\max(33,c_T)\le n_T\le 24c_T+8$;
moreover $n_T$ equals exactly half the length of the shortest kernel
element --- a shortest-vector problem in the ideal $2(t-1)\mu_T\Z[t^{\pm1}]$
under the lamplighter metric --- from which we extract a closed-form
three-regime law, $n_T = 6c_T+e_T$, $9c_T-3e_T$, or $3(c_T+e_T)$ according
to the angle of the triangle, certified at $(1,2)$ and validated on
thirteen shapes. Finally, the
lamplighter transfer yields a $q$-difference functional equation for
the relaxed generating function of A396406; a connectivity-invariance
squeeze of the true count between its pure-travel subcount and the
relaxed series pins the exact growth rate unconditionally:
$\beta_2 = 1.4916177871\ldots$, the reciprocal square root of the
smallest positive root of an explicit travel-block equation, which
excludes the value $3/2$ suggested by finite-depth ratios. The eight
length-$10$ relations and the universality of $u_d$ through depth $22$
are formally machine-verified in Lean~4; the depth-$32$ threshold and
the deviation witnesses are certified by exact modular and rational
arithmetic. Transcendence of the generating functions over $\Q(x)$, and
of the number $\beta_2$, lies outside the present scope.
\end{abstract}

\section{Introduction}
\label{sec:intro}

Let $T$ be a triangle with vertices in $\Q^2$, and for
$i\in\{0,1,2\}$ let $R_i$ be the reflection of $\R^2$ across the line
carrying the side opposite vertex $V_i$. The group
$W(T)=\langle R_0,R_1,R_2\rangle\le\Isom(\R^2)$ acts on the orbit of
$T$, and breadth-first search in the Cayley graph relative to
$\{R_0,R_1,R_2\}$ stratifies that orbit by word-length. Writing
$u_d=u_d^T$ for the number of distinct triangle-images at minimal
word-length $d$, one obtains an integer sequence depending a priori on
the shape of $T$. For right triangles with unequal legs this sequence
is OEIS A396406. This paper determines its group-theoretic origin, its
exact word-length metric, the precise extent to which it is
shape-independent, and its exact exponential growth rate.

\paragraph{The generic group.}
When $T$ is a non-isosceles right triangle, the two legs meet at a
right angle, so $(R_0R_1)^2=1$, while by Niven's theorem \cite{Niven1956} the remaining
two angles are irrational multiples of $\pi$ and force no further
relation among the reflections. The abstract group is therefore the
Coxeter group
\[
  W=\bigl\langle R_0,R_1,R_2 \;\bigm|\;
    R_i^2=1\ (i=0,1,2),\ (R_0R_1)^2=1 \bigr\rangle ,
\]
whose Poincar\'e series is $(1+t)^2/(1-t-t^2)$; consequently the
free-group baseline gives $u_d=F_{d+3}$ for $1\le d\le 9$, until the
first affine collision. The shape $T$ enters only through the affine
representation $\rho_T\colon W\to\Aff(\R^2)$ realizing the three
reflections, and the sequence $u_d^T$ records the layer sizes of
$\rho_T(W)$. Our first structural result identifies the kernel of the
generic representation. Placing the right angle at the origin and
writing $\zeta_T=(a+bi)/(a-bi)$ for the rotation number of the
hypotenuse reflection, every group element acts as
$z\mapsto \varepsilon\,\zeta_T^{\,k}\sigma^{\delta}(z)+P(\zeta_T)$ for an
integer Laurent polynomial $P$ \emph{independent of $T$}
(Lemma, symbolic normal form). The pure-translation subgroup of the
resulting symbolic group is computed exactly,
\[
  \mathcal T \;=\; 2(t-1)\,\Z[t^{\pm1}],
\]
by reducing modulo $2$ — where the two perpendicular legs collapse onto
a single mirror and the group degenerates to the infinite dihedral
group, which has no translations. This caps the glide-square lattice
$2(t-1)\Z[t^{\pm1}]\subseteq\mathcal T$ to an equality
(Theorem, translation lattice).

\paragraph{Lamplighter structure and the word-length metric.}
The translation lattice exhibits the generic group as virtually a
classical object: the subgroup $H\le W_{\mathrm{gen}}$ on which the
linear part is a pure rotation $t^k$ is normal of index $4$ and
isomorphic to $\Z[t^{\pm1}]\rtimes_t\Z=\Z\wr\Z$, so every
right-triangle side-reflection group with transcendental leg ratio is
virtually the wreath product $\Z\wr\Z$ (Theorem, lamplighter
structure) — metabelian-by-finite, hence amenable, yet of exponential
growth, and not finitely presentable. Independently, the translation
module is computed by Mayer--Vietoris for the free product
$W=V_4*\Z/2$ together with Crowell's exact sequence: it is the
$\Z[Q]$-module $I_A\cap I_B$, which equals the cyclic ideal generated
by the central element $h=(c-1)(Y-1)$ by exactness of the Bass--Serre
tree of $D_\infty$ (Lemma, Mayer--Vietoris reduction). The generic
kernel is then identified as the relative commutator subgroup
$\ker\rho_{\mathrm{gen}}=[T_\rho,T_\rho]$, exhibiting the generic group
as the $Q$-relative metabelianization of $V_4*\Z/2$
(Corollary, generic kernel). Reading the composition law in the edge
basis turns the group into an explicit \emph{lamp model}: a walker on
$\Z$ deposits signed lamps on the edges it traverses, with a normal
form $(\varepsilon,\delta,k,P)$ cut out by two congruences. From this
model we obtain a closed-form word-length metric (Theorem, metric
formula): the geodesic length is a minimum over crossing-number and
site-pairing data, equal to a relaxed length plus exactly $2$ for each
isolated cycle that the connectivity (single-Eulerian-trail) constraint
forbids. The metric is forced by a strand bound limiting each interface
component to at most one up- and one down-crossing per edge, leaving a
single unbounded catalytic datum; it reproduces A396406 through $u_{22}$
with no breadth-first search. The kernel depends only on the similarity
class of $T$ (Lemma, dilation invariance).

\paragraph{Universality and its sharp threshold.}
Because the symbolic normal form is triangle-independent and $\zeta_T$
is never a root of unity (Niven), two group elements collide under
$\rho_T$ only when their symbolic data agree and their translation
polynomials agree at $t=\zeta_T$; an extra collision at depth $d$ thus
forces a nonzero integer Laurent polynomial of height $\le 2d$ to
vanish at $\zeta_T$. By Gauss's lemma the shape's minimal polynomial
$\mu_T(t)=c_T t^2-e_T t+c_T$ must then divide a realized
word-difference, so $c_T$ divides both extreme coefficients and
$c_T\le 2d$. This yields an \emph{effective all-depths universality
theorem}: every shape with $c_T>2d$ satisfies $u_d^T=u_d$
(Theorem, effective universality). The arithmetic of $c_T$ — odd, with
all prime factors $\equiv 1\pmod 4$, and $\ge 5$ — leaves only a sparse
finite set of candidate shapes at each depth, and a one-sided exact
modular certificate over these candidates closes the remaining gap:

\begin{quote}
\emph{For every right triangle with positive unequal rational legs,
$u_d^T=u_d$ for all $0\le d\le 32$, and the bound is sharp: the $(1,2)$
triangle satisfies $u_{33}^{(1,2)}=u_{33}-30<u_{33}$}
(Theorem, sharp uniform universality).
\end{quote}

\paragraph{Refutation of all-depths universality.}
The natural conjecture that $u_d^T=u_d$ for \emph{all} $d$ is false for
every algebraic shape. Since $R_0R_1$ is the half-turn about the
right-angle vertex, $R_0R_1R_2$ is a glide reflection along the
hypotenuse, and its square is the pure translation by $2(t^{-1}-1)$;
conjugating and multiplying these glide-squares realizes the entire
ideal $2(t-1)\Z[t^{\pm1}]$. As that ideal contains a multiple of every
$\mu_T$, an explicit word
\[
  w_T=\tau^{\,c_T}\,r\,\tau^{\,-e_T}\,r\,\tau^{\,c_T}\,r^{-2},
  \qquad \tau=(R_0R_1R_2)^2,\quad r=R_x R_h,
\]
realizes the translation $2(t^{-1}-1)\mu_T(t)$, which is nonzero in the
generic group but evaluates to $0$ at $\zeta_T$, hence lies in
$\ker\rho_T$ (Theorem, non-universality). Thus $\ker\rho_T$ is, for
every algebraic shape, the principal ideal $2(t-1)\mu_T\,\Z[t^{\pm1}]$,
isomorphic to $\Z[t^{\pm1}]$ and normally generated by the single
witness $w_T$ (Corollary, kernel classification). The first deviation
depth is therefore finite and pinned to linear growth in the arithmetic
complexity of the shape,
\[
  \max(33,\,c_T)\;\le\;n_T\;\le\;6(2c_T+|e_T|)+8\;\le\;24c_T+8
\]
(Corollary, deviation-depth sandwich), with $n_{(1,2)}=33$ realizing
the lower bound; the deviation depth is itself a shortest-vector length
in the lamplighter word metric. Consequently the universal sequence is
exactly the orbit growth of the generic right triangle, realized
precisely by the (uncountably many) transcendental leg ratios, while
the isosceles case ($c_T=1$, deviation at depth $4$) is the extreme
degeneration to a wallpaper group (Corollary, transcendence dichotomy).

\paragraph{The exact growth rate.}
The lamp model linearizes the metric into a transfer with one catalytic
variable, but the catalytic variable enters through a \emph{dilation}
$t\mapsto q^2 t$, so the resulting functional equation for the relaxed
generating function is a linear $q$-difference equation, not a
polynomial-kernel equation; the Bousquet-M\'elou--Jehanne framework does
not apply and the solution is $q$-holonomic. Unfolding the equation
telescopes the relaxed series into explicit rapidly-convergent
$q$-series blocks, and the dominant singularity is the smallest positive
root $q^*$ of the travel-block equation $\Sigma_1(q)=1$. We prove that
this rate is also the \emph{true} rate, unconditionally: the travel
block is connectivity-invariant (each pure-translation run is a single
Euler trail and incurs no penalty), so the true series shares it
verbatim; squeezing the true count between its pure-travel subcount and
the relaxed series gives
\[
  \beta_2 \;=\; 1/\sqrt{q^*}\;=\;1.4916177871143742268\ldots
\]
exactly (Proposition, asymptotic ratio), with no dependence on any
analytic conjecture. In particular $\beta_2<3/2$, excluding the value
that the finite-depth ratios — converging at rate $O(1/d)$ with a
period-$4$ oscillation — had suggested.

\paragraph{What is conditional and what is open.}
The growth rate $\beta_2=1.4916177871\ldots$ is determined unconditionally in Section~\ref{sec:beta2} as $1/\sqrt{q^\ast}$, where $q^\ast$ is the least positive root of $\Sigma_1(q)=1$; in particular $\beta_2\neq 3/2$. Transcendence questions are treated in the companion paper \cite{Bonfioli2026b}, where the
growth series of A396406 and of its relaxed companion are proved transcendental over $\Q(x)$
(each catalytic building block unconditionally so, by a P\'olya--Carlson natural boundary);
this explains the empirical absence, recorded below, of any low-order linear, holonomic, or
low-degree algebraic relation over the $43$ known terms. The transcendence of the \emph{number}
$\beta_2$ itself remains open. In higher dimensions --- the reflection groups of Euclidean
orthoschemes (path simplices), of which the right triangle is the planar case --- the abstract
right-angled Coxeter envelope has \emph{rational} growth with rate $r_n=1+2\cos\frac{2\pi}{n+3}$, but
the geometric group is an amenable proper quotient of it, matching the rational envelope only up to a
finite depth before deviating; the plane is the one dimension in which that deviation has been carried
past its onset (at depth $10$) and found transcendental, and whether higher dimensions are eventually
transcendental as well is open \cite{Bonfioli2026c}.

\paragraph{Reproducibility.}
The eight length-$10$ affine relations and the shape-independence of
$u_d$ through depth $22$ are formally verified in the Lean~4 proof
assistant; the depth-$32$ threshold, the depth-$38$ data provenance of
A396406, and the deviation witnesses are certified by exact modular and
rational arithmetic, with two independent reproductions. All scripts and
Lean sources are described in the reproducibility section.

\paragraph{Related work.}
The problem sits between Coxeter-group word-growth and the geometry of
metabelian groups. Isola and Piergallini \cite{IsolaPiergallini} show that the reflection group of a generic triangle
carries no unexpected relations; the present results are the explicit
rational right-triangle analogue, replacing genericity by a Diophantine
faithfulness criterion in $\zeta_T$ and adding exact word-growth data.
The wreath product $\Z\wr\Z$ that emerges is the textbook amenable
group of exponential growth; its appearance here, generated by three
mirrors of a Euclidean triangle, places the rational shapes inside the
program of quantifying group approximations, with the deviation sandwich
giving the explicit girth of convergence in the space of marked groups.
The geodesic-language and growth analysis is in the spirit of the
lamplighter computations of Cleary--Elder--Taback; the catalytic
$q$-difference structure of the growth series, and the consequent
exclusion of every elementary closed form for $\beta_2$, is the point at
which A396406 departs from the rational growth series of the standard
lamplighter generating sets. The rotation number $\zeta_T$ governing the
translation module coincides with the rotation number of Sadun's
generalized-pinwheel substitution tilings, though the invariant studied
there (tiling-space cohomology) is distinct from the word-growth
sequence treated here.

% ======== Setting and the Coxeter presentation ========
\section{Setting and the Coxeter presentation}
\label{sec:setting}

\subsection{The triangle, its mirrors, and the orbit-growth sequence}
\label{ssec:setting-defs}

Let $T$ be a triangle with vertices $V_0, V_1, V_2 \in \mathbb{Q}^2$ and leg lengths $a, b \in \mathbb{Q}_{>0}$, having its right angle at $V_2$ and \emph{unequal} legs, $a \neq b$. For $i \in \{0,1,2\}$ let
\[
  R_i \in \mathrm{Isom}(\mathbb{R}^2)
\]
denote the Euclidean reflection across the line containing the side of $T$ opposite the vertex $V_i$, and set
\[
  W(T) \;=\; \langle R_0, R_1, R_2 \rangle \;\le\; \mathrm{Isom}(\mathbb{R}^2).
\]
When explicit coordinates are required we normalize the right-angle vertex to $B = (a,0)$ with $A = (0,0)$ and $C = (a,b)$, writing $a,b$ for the two leg lengths; the mirrors $R_0, R_1$ are then the reflections across the two legs meeting at the right angle.

\begin{remark}[Notation dictionary]
\label{rem:notation-dictionary}
Throughout the paper the three generators carry two interchangeable names, fixed once here:
\[
  R_0 = R_x \ (\text{leg on the } x\text{-axis}), \qquad
  R_1 = R_y \ (\text{leg on the } y\text{-axis}), \qquad
  R_2 = R_h \ (\text{hypotenuse}).
\]
The subscripted names $R_0, R_1, R_2$ index the opposite vertices $V_0, V_1, V_2$; the lettered names $R_x, R_y, R_h$ name the mirror lines and are used from \S\ref{sec:bonfioli} onward, where the right-angle vertex is placed at the origin with the two legs along the positive coordinate axes (the normalization $A=(0,0)$, $B=(a,0)$, $C=(a,b)$ above, rescaled by $1/a$). The two normalizations differ only by a similarity and induce the same abstract relations and orbit counts.
\end{remark}

We study the orbit of $T$ under $W(T)$ through breadth-first search in the Cayley graph of $W(T)$ relative to the generating set $S = \{R_0, R_1, R_2\}$. For $d \ge 0$ put
\[
  \mathcal{O}_d(T) \;=\; \bigl\{\, g \cdot T : g \in W(T),\ \ell_S(g) = d \,\bigr\},
  \qquad
  u_d^T \;:=\; \bigl| \mathcal{O}_d(T) \bigr|,
\]
where $\ell_S$ is word length with respect to $S$, and $u_0^T = 1$. All distances are computed in exact rational arithmetic, and the relations below are verified symbolically over $\mathbb{Q}(a,b)$. We write $u_d := u_d^T$ when the shape is fixed or immaterial.

\begin{remark}
\label{rem:two-layers}
Two groups govern $u_d$. The \emph{abstract} group $W$ is the finitely presented Coxeter group of Theorem~\ref{thm:coxeter-presentation} below, depending only on the right-angle and non-isosceles hypotheses. The \emph{Euclidean} group $W(T)$ is its image under the affine action $\rho_T \colon W \to \mathrm{Aff}(\mathbb{R}^2)$ induced by the mirrors of $T$, which depends on the shape. Distinct words of $W$ may act identically on $\mathbb{R}^2$, collapsing orbit elements; the counts $u_d$ are those of the Euclidean orbit (Definition~\ref{def:deficit}).
\end{remark}

\subsection{The Coxeter presentation}
\label{ssec:coxeter-presentation}

Let $\Gamma$ be the abstract group obtained from the free product $\langle r_0\rangle * \langle r_1\rangle * \langle r_2\rangle$ of three copies of $\mathbb{Z}/2$ by the canonical map onto $W(T)$ sending $r_i \mapsto R_i$, and let $K \trianglelefteq \langle r_0\rangle * \langle r_1\rangle * \langle r_2\rangle$ be the kernel of this map: $K$ is the set of all relations holding in $\mathrm{Isom}(\mathbb{R}^2)$ among $R_0, R_1, R_2$. The presentation theorem identifies $K$.

\begin{theorem}[Coxeter presentation]
\label{thm:coxeter-presentation}
Let $T$ be a right triangle with the right angle at $V_2$ and with unequal legs $a \neq b$. Let
\[
  W \;:=\; \bigl\langle\, R_0, R_1, R_2 \ \bigm|\ R_i^2 = 1\ (i=0,1,2),\ \ (R_0 R_1)^2 = 1 \,\bigr\rangle
\]
be the Coxeter group whose diagram is the path on three nodes with edge labels $m(R_0,R_1) = 2$ and $m(R_0,R_2) = m(R_1,R_2) = \infty$. Then the pairwise rotation orders of the three mirrors of $T$ are exactly $m(R_0,R_1) = 2$, $m(R_0,R_2) = m(R_1,R_2) = \infty$, and consequently $r_i \mapsto R_i$ induces a \emph{surjection} $\rho_T \colon W \twoheadrightarrow W(T)$. This surjection is \emph{not} injective for any admissible shape: the affine action imposes further relations beyond $(R_0R_1)^2 = 1$, so $W(T)$ is a \emph{proper} quotient of $W$, and $\ker\rho_T \supsetneq \{1\}$. The first such relation is the object of Theorem~\ref{thm:non-universality-restate}; the abstract Coxeter group $W$ therefore serves only as an \emph{envelope} for $W(T)$.
\end{theorem}

\begin{proof}
Each $R_i$ is a reflection, so $R_i^2 = 1$; the mirrors of $R_0,R_1$ meet at $\pi/2$, so $R_0R_1$ is a half-turn and $(R_0R_1)^2 = 1$. Hence all defining relators of $W$ hold in $W(T)$, giving the surjection $\rho_T\colon W \twoheadrightarrow W(T)$. The remaining interior angles $\arctan(b/a)$, $\arctan(a/b)$ are, by Niven's theorem (a rational multiple of $\pi$ has rational tangent only for tangent in $\{0,\pm1\}$, excluded by $a\ne b$), irrational multiples of $\pi$; so $R_0R_2$ and $R_1R_2$ have infinite order and no bond $(R_iR_2)^m=1$ is forced. Thus the pairwise orders are exactly as stated, and $W$ is the largest group compatible with them.

The surjection is \emph{not} an isomorphism. Coxeter's structure theorem would identify $W(T)$ with $W$ only if the affine (Tits) representation $\rho_T$ were faithful; but $W(T)\le\mathrm{Isom}(\mathbb{R}^2)$ is amenable (indeed virtually lamplighter, \S\ref{sec:metric}), whereas $W\cong(\mathbb{Z}/2\times\mathbb{Z}/2)*\mathbb{Z}/2$ contains a free subgroup and is non-amenable. Hence $\ker\rho_T$ is nontrivial and $W(T)$ is a proper quotient. Explicitly, $\ker\rho_T$ is computed in Theorem~\ref{thm:non-universality-restate}: its shortest element has length $2n_T$, so $W(T)$ and $W$ have the same spherical growth only through depth $n_T-1$.
\end{proof}

\begin{remark}
\label{rem:isosceles-excluded}
The non-isosceles hypothesis $a \neq b$ is used solely to exclude the angle $\pi/4$ in Theorem~\ref{thm:coxeter-presentation}; for the isosceles right triangle the third relation $(R_0 R_2)^4 = 1$ holds and $W$ is a finite Coxeter group, a degenerate case outside the present scope.
\end{remark}

\subsection{The Steinberg--Poincar\'e series and the Fibonacci phase}
\label{ssec:fibonacci-corollary}

The growth of the abstract layer is governed by the spherical growth series of the Coxeter group $W$ of Theorem~\ref{thm:coxeter-presentation}, i.e.\ the Poincar\'e series
\[
  W(t) \;=\; \sum_{d \ge 0} \bigl| \{\, g \in W : \ell_S(g) = d \,\} \bigr|\, t^d .
\]

\begin{proposition}[Steinberg--Poincar\'e series]
\label{prop:steinberg}
For the Coxeter group $W$ of Theorem~\ref{thm:coxeter-presentation},
\[
  W(t) \;=\; \frac{(1+t)^2}{1 - t - t^2}.
\]
\end{proposition}

\begin{proof}
Steinberg's formula \cite{Humphreys1990} expresses the Poincar\'e series of a Coxeter system $(W,S)$ as
\[
  \frac{1}{W(t^{-1})} \;=\; \sum_{\substack{J \subseteq S \\ W_J \text{ finite}}} \frac{(-1)^{|J|}}{W_J(t)},
\]
summed over those subsets $J$ generating a finite parabolic $W_J$, with $W_J(t)$ its (polynomial) Poincar\'e series. Here the finite parabolics are: $J = \emptyset$ ($W_\emptyset(t) = 1$); the three singletons $\{R_i\}$, each with $W_J(t) = 1+t$; and the rank-$2$ subset $J = \{R_0, R_1\}$, which generates the dihedral group of order $4$ with $W_J(t) = (1+t)^2$. The subsets $\{R_0,R_2\}$ and $\{R_1,R_2\}$ generate the infinite dihedral group and contribute nothing. Substituting,
\[
  \frac{1}{W(t^{-1})}
  = 1 - \frac{3}{1+t} + \frac{1}{(1+t)^2}
  = \frac{(1+t)^2 - 3(1+t) + 1}{(1+t)^2}
  = \frac{t^2 - t - 1}{(1+t)^2}.
\]
Inverting both sides,
\[
  W(t^{-1}) \;=\; \frac{(1+t)^2}{t^2 - t - 1}.
\]
Now substitute $t \mapsto t^{-1}$. The numerator becomes $(1 + t^{-1})^2 = t^{-2}(1+t)^2$, and the denominator becomes $t^{-2} - t^{-1} - 1 = t^{-2}(1 - t - t^2)$; the common factor $t^{-2}$ cancels, giving
\[
  W(t) \;=\; \frac{(1+t)^2}{1 - t - t^2}. \qedhere
\]
\end{proof}

\begin{corollary}[Abstract layer counts are Fibonacci]
\label{cor:abstract-fibonacci}
Write $F_n$ for the Fibonacci numbers with the conventions $F_0 = 0$, $F_1 = F_2 = 1$. The coefficients of the Steinberg--Poincar\'e series are
\[
  [t^d]\, W(t) \;=\; F_{d+3}
  \qquad (d \ge 1),
\]
with $[t^0]\, W(t) = 1$.
\end{corollary}

\begin{proof}
Write $(1+t)^2 = 1 + 2t + t^2$ and let $\Phi(t) = 1/(1 - t - t^2) = \sum_{n \ge 0} F_{n+1} t^n$, the Fibonacci generating function. Then for $d \ge 1$,
\[
  [t^d]\, W(t) = F_{d+1} + 2 F_{d} + F_{d-1},
\]
where the term $F_{d-1}$ is read with the convention $F_0 = 0$ when $d = 1$. Using $F_{d+1} + F_d = F_{d+2}$ and $F_d + F_{d-1} = F_{d+1}$, this equals $F_{d+2} + F_{d+1} = F_{d+3}$. In particular at $d = 1$ one gets $F_2 + 2F_1 + F_0 = 1 + 2 + 0 = 3 = F_4$, consistent with the direct expansion $W(t) = 1 + 3t + 5t^2 + 8t^3 + \cdots$; the boundary value $[t^0]\,W(t) = 1$ is read off the same expansion.
\end{proof}

The Euclidean orbit counts $u_d$ agree with this baseline until the first Euclidean collision.

\begin{corollary}[Fibonacci phase]
\label{cor:fibonacci-phase}
Fix the reference triangle $T_0 = (3,4,5)$. Then
\[
  u_d^{T_0} \;=\; F_{d+3}
  \qquad (1 \le d \le 9),
\]
and the first depth at which the affine action collapses an abstract orbit element is $d = 10$, where $u_{10}^{T_0} = 225 < F_{13} = 233$. The same values hold for every right triangle $T$ with unequal rational legs; this shape-independence on the range $1 \le d \le 9$ and the depth-$10$ deficit are established symbolically over $\mathbb{Q}(a,b)$ in \S\ref{sec:bonfioli} (Theorem~\ref{thm:symbolic}). At $d = 0$ one has $u_0 = 1 \neq F_3 = 2$.
\end{corollary}

\begin{proof}
By Proposition~\ref{prop:steinberg} and Corollary~\ref{cor:abstract-fibonacci}, the abstract layer count at depth $d$ is $F_{d+3}$, and $u_d^T \le F_{d+3}$ for every $T$ since $u_d^T$ counts the image $\rho_T(\mathcal{O}_d)$. For the reference triangle $T_0 = (3,4,5)$, exact-rational breadth-first search shows that the map from abstract length-$d$ elements to Euclidean orbit points is injective for $1 \le d \le 9$, so $u_d^{T_0} = F_{d+3}$ there, and that $u_{10}^{T_0} = 225 = F_{13} - 8$, exhibiting eight affine collisions at length $10$. That these values are independent of the admissible shape is not a consequence of the single search but is proved by the symbolic relation analysis of \S\ref{sec:bonfioli}: Theorem~\ref{thm:symbolic} shows over $\mathbb{Q}(a,b)$ that no nontrivial affine identification occurs below depth $10$ and that exactly eight occur at depth $10$ for every admissible $T$.
\end{proof}

\begin{definition}[Fibonacci deficit]
\label{def:deficit}
The \emph{Fibonacci deficit} of $T$ at depth $d$ is
\[
  \delta_d^T \;:=\; F_{d+3} - u_d^T \;\ge\; 0 ,
\]
the number of abstract length-$d$ orbit elements identified by the affine action $\rho_T$. By Corollary~\ref{cor:fibonacci-phase} and Theorem~\ref{thm:symbolic}, $\delta_d^T = 0$ for $1 \le d \le 9$ and $\delta_{10}^T = 8$ for every admissible $T$. The Euclidean orbit-growth sequence $(u_d^T)_{d \ge 0}$ is recorded as OEIS \texttt{A396406}.
\end{definition}

\begin{remark}
\label{rem:setting-notation}
Throughout, $\rho_T \colon W \to \mathrm{Aff}(\mathbb{R}^2)$ denotes the affine action of the abstract Coxeter group $W = \langle R_0, R_1, R_2 \rangle$ induced by the mirrors of $T$, and $\zeta_T = (a + bi)/(a - bi)$ is the rotation number attached to the leg pair $(a,b)$; both enter the analysis from \S\ref{sec:bonfioli} onward. The results of this section — the presentation of Theorem~\ref{thm:coxeter-presentation}, the series of Proposition~\ref{prop:steinberg}, and the Fibonacci phase of Corollary~\ref{cor:fibonacci-phase} — are unconditional.
\end{remark}

% ======== The eight length-10 relations and depth-10 completeness ========
\section{The eight length-\texorpdfstring{$10$}{10} relations and depth-\texorpdfstring{$10$}{10} completeness}
\label{sec:bonfioli}

Let $W = \langle\, R_0, R_1, R_2 \mid R_i^2 = 1\ (i=0,1,2),\ (R_0R_1)^2 = 1 \,\rangle$
be the Coxeter group of \S\ref{sec:setting}. Throughout, $R_0 = R_x$ is the
reflection across the $x$-leg, $R_1 = R_y$ the reflection across the $y$-leg,
and $R_2 = R_h$ the reflection across the hypotenuse. For a right triangle $T$
let $\rho_T \colon W \to \operatorname{Aff}(\mathbb{R}^2)$ be its side-reflection
action. Fix the placement of \S\ref{sec:setting},
\[
A = (0,0), \qquad B = (a,0), \qquad C = (a,b), \qquad a, b > 0,
\]
so that each generator $R_i$ acts as an affine isometry whose $3\times 3$ matrix
has entries in $\mathbb{Q}(a,b)$, the denominators dividing a power of $a^2+b^2$.
For a word $w = w[0]\,w[1]\cdots w[\ell-1]$ in the $R_i$ we write
$\rho_T(w) = R_{w[0]}\,R_{w[1]}\cdots R_{w[\ell-1]}$ for the associated affine matrix.

The Poincar\'e series of the abstract Coxeter group $W$ gives
$F_{d+3}$ canonical words of length $d$ for $d \ge 1$; in particular there are
$F_{13} = 233$ canonical words of length $10$. Exact-rational breadth-first
enumeration of the orbit of a non-isosceles right triangle gives $u_{10} = 225$,
a deficit of $8$. The eight collapsing pairs are the length-$10$ relations of
Table~\ref{tab:bonfioli}; they are the first relations of the side-reflection
action beyond the abstract Coxeter relations.

\begin{table}[ht!]
\centering
\small
\renewcommand{\arraystretch}{1.2}
\begin{tabular}{r l l}
\toprule
\# & word $A$ & word $B$ \\
\midrule
1 & $R_0 R_1 R_2 R_0 R_2 R_0 R_1 R_2 R_1 R_2$ &
    $R_2 R_1 R_2 R_0 R_1 R_2 R_0 R_2 R_0 R_1$ \\
2 & $R_0 R_1 R_2 R_1 R_2 R_0 R_1 R_2 R_0 R_2$ &
    $R_2 R_0 R_2 R_0 R_1 R_2 R_1 R_2 R_0 R_1$ \\
3 & $R_0 R_2 R_0 R_1 R_2 R_1 R_2 R_0 R_1 R_2$ &
    $R_2 R_0 R_1 R_2 R_1 R_2 R_0 R_1 R_2 R_0$ \\
4 & $R_0 R_2 R_0 R_2 R_0 R_1 R_2 R_1 R_2 R_0$ &
    $R_1 R_2 R_1 R_2 R_0 R_1 R_2 R_0 R_2 R_1$ \\
5 & $R_0 R_2 R_0 R_2 R_0 R_1 R_2 R_1 R_2 R_1$ &
    $R_1 R_2 R_1 R_2 R_0 R_1 R_2 R_0 R_2 R_0$ \\
6 & $R_0 R_2 R_1 R_2 R_0 R_1 R_2 R_0 R_2 R_0$ &
    $R_1 R_2 R_0 R_2 R_0 R_1 R_2 R_1 R_2 R_1$ \\
7 & $R_0 R_2 R_1 R_2 R_0 R_1 R_2 R_0 R_2 R_1$ &
    $R_1 R_2 R_0 R_2 R_0 R_1 R_2 R_1 R_2 R_0$ \\
8 & $R_1 R_2 R_0 R_1 R_2 R_0 R_2 R_0 R_1 R_2$ &
    $R_2 R_0 R_1 R_2 R_0 R_2 R_0 R_1 R_2 R_1$ \\
\bottomrule
\end{tabular}
\caption{The eight length-$10$ affine relations. In each row, the two words
have the same affine matrix in the side-reflection action $\rho_T$ of $W$ on
every right triangle $T$. The generators are $R_0 = R_x$, $R_1 = R_y$,
$R_2 = R_h$ (reflections across the $x$-leg, $y$-leg, and hypotenuse).}
\label{tab:bonfioli}
\end{table}

\begin{theorem}[Eight length-$10$ affine relations]
\label{thm:symbolic}
Each of the eight identities of Table~\ref{tab:bonfioli} holds for every right
triangle. For the placement $A=(0,0)$, $B=(a,0)$, $C=(a,b)$ with $a,b>0$, the
affine matrices
\[
M_A = R_{w_A[0]}\cdots R_{w_A[9]}, \qquad
M_B = R_{w_B[0]}\cdots R_{w_B[9]}
\]
of the two words $w_A, w_B$ in each row satisfy $M_A = M_B$ as $3\times 3$
affine matrices over $\mathbb{Q}(a,b)$.
\end{theorem}

\begin{proof}
The statement is a finite algebraic identity. Each reflection $R_i$ acts as an
affine isometry whose matrix entries are rational functions of $a, b$ over
$\mathbb{Q}$, with denominators dividing a power of $a^2+b^2$. Hence each entry
of $M_A - M_B$ is a rational function in $\mathbb{Q}(a,b)$; clearing
denominators yields, for each row, a polynomial identity in $\mathbb{Q}[a,b]$.
These identities are verified deterministically in exact arithmetic over
$\mathbb{Q}(a,b)$.

The eight identities are moreover \emph{formally verified} in the Lean~4 proof
assistant. The theorems
\[
\texttt{rel1\_symbolic},\ \dots,\ \texttt{rel8\_symbolic} \quad\text{in}\quad
\texttt{lean/with\_mathlib/SymbolicVerification.lean}
\]
each discharge the six affine-matrix-entry identities over $\mathbb{Q}[a,b]$
that together express $M_A = M_B$. Each such identity is proved by the
\texttt{ring} tactic alone. The \texttt{ring} tactic produces a proof term that
is checked by the Lean kernel; the trust base is therefore the Lean kernel
together with the \texttt{Mathlib} definitions of $\mathbb{Q}[a,b]$ and its
commutative-ring structure, with no appeal to numerical sampling, decision
heuristics, or specialization to particular legs. Consequently each identity
holds as an equality of affine matrices over $\mathbb{Q}(a,b)$, hence for every
right triangle.
\qed
\end{proof}

\begin{remark}
Because Theorem~\ref{thm:symbolic} holds over $\mathbb{Q}(a,b)$, the eight
collisions lie in $\ker \rho_T$ for every $T$; the enumerative deficit
$233 - 225 = 8$ is therefore shape-independent.
\end{remark}

\begin{proposition}[Completeness at depth $10$]
\label{prop:bonfioli-completeness}
Let $T$ be an unequal-leg ($a \ne b$) right triangle. Among the $233$ canonical
Coxeter words of length $10$ in $W$, exactly $225$ have pairwise distinct images
under $\rho_T$: there are $217$ singleton classes and $8$ collision pairs, each
pair being one row of Table~\ref{tab:bonfioli}. No length-$10$ affine identity
exists beyond these eight relations and the abstract Coxeter relations. (For the
isosceles triangle $a = b$ additional collisions occur, by
Remark~\ref{rem:isosceles-excluded}; the count $225$ and the distinctness of the
$217$ singleton classes are asserted only for $a \ne b$.)
\end{proposition}

\begin{proof}
Completeness follows from two ingredients: (a) the computed block partition on
the reference triangle $T_0 = (3,4,5)$, and (b) Theorem~\ref{thm:symbolic}.

All $233$ canonical Coxeter words of length $10$ are enumerated. For each, the
affine matrix $\rho_{T_0}(w)$ is computed in exact rational arithmetic on $T_0$,
and the words are grouped by matrix. The resulting partition has exactly $225$
blocks: eight blocks of size $2$ and $217$ blocks of size $1$. The eight
size-$2$ blocks are the rows of Table~\ref{tab:bonfioli}. By
Theorem~\ref{thm:symbolic} these eight identities hold symbolically over
$\mathbb{Q}(a,b)$, so the eight pairs are genuine collisions for every right
triangle. Conversely, any further length-$10$ collision would produce a
coincidence of affine matrices already on $T_0$ in exact arithmetic,
contradicting the computed block sizes. Since $T_0$ is unequal-leg, the $217$
singleton blocks remain singletons under specialization to any unequal-leg $T$.
Hence the eight relations, together with the abstract Coxeter relations used to
form the canonical words, account for all length-$10$ affine identities.

As corroboration, each size-$2$ block was independently re-checked on six
further triangles,
\[
(5,12),\quad (8,15),\quad (7,24),\quad (1,2),\quad (2,3),\quad (1,7),
\]
and in all cases the two words of the block produce the same affine matrix; this
is redundant once Theorem~\ref{thm:symbolic} is invoked.
\qed
\end{proof}

% ======== The triangle-independent symbolic normal form ========
\section{The triangle-independent symbolic normal form}
\label{sec:symbolic-normal-form}

The side-reflection action carries a triangle-independent symbolic normal
form (Lemma~\ref{lem:symbolic-form}) in which the shape of the triangle
enters only through the evaluation of integer Laurent polynomials at the
rotation number $\zeta_T$. Collisions between group elements reduce to the
vanishing of bounded-height Laurent polynomials at $\zeta_T$
(Proposition~\ref{prop:collision-vanishing}), a divisibility condition in
$\mathbb{Z}[i]$ that drives the effective universality theorem of
\S\ref{sec:effective-universality}.

\subsection{Complex coordinates and the generators}

Throughout, the three side-reflections are written interchangeably as
$(R_0, R_1, R_2) = (R_x, R_y, R_h)$, where $R_0 = R_x$ is the reflection in
the leg on the positive real axis, $R_1 = R_y$ the reflection in the leg on
the positive imaginary axis, and $R_2 = R_h$ the reflection in the
hypotenuse; this matches the labelling $R_0, R_1, R_2$ of
\S\ref{sec:effective-universality}.

Identify $\mathbb{R}^2 = \mathbb{C}$, place the right angle at the origin
with the legs along the positive coordinate axes, and rescale by $1/a$,
where $a$ is the length of the leg on the positive real axis. This is the
normalization used in \S\ref{sec:effective-universality}; it agrees, after
the affine change placing the right-angle vertex at the origin, with the
normalization $A = (0,0)$, $B = (a,0)$, $C = (a,b)$ of the introductory
sections. Neither the choice of origin and axes nor the rescaling changes
word-lengths or the collision structure of the orbit, so no generality is
lost. Let $\sigma(z) = \bar z$ denote complex conjugation, and set
\[
\zeta = \zeta_T = \frac{a + bi}{a - bi},
\]
the rotation number of the triangle with legs $(a,b)$. Since numerator and
denominator are complex conjugates, $|\zeta_T| = 1$; we use this standing
fact below. In these coordinates the three side-reflections are
\[
R_x : z \mapsto \sigma(z), \qquad
R_y : z \mapsto -\,\sigma(z), \qquad
R_h : z \mapsto \zeta^{-1}\sigma(z) + \bigl(1 - \zeta^{-1}\bigr),
\]
where $R_x$ and $R_y$ are the reflections in the two leg lines and $R_h$
is the reflection in the hypotenuse line. We write $W$ for the group they
generate and $\rho_T \colon W \to \operatorname{Isom}(\mathbb{C})$ for the
resulting faithful action on the plane of the triangle $T$.

Every element of $\operatorname{Isom}(\mathbb{C})$ produced by composing
these generators is an orientation-reversing-or-preserving affine map of
the form $z \mapsto \alpha\, \sigma^{\delta}(z) + \beta$ with
$\delta \in \{0,1\}$. The normal form below records that the linear part
$\alpha$ is always of the shape $\varepsilon\,\zeta^{k}$ with
$\varepsilon \in \{\pm 1\}$ and $k \in \mathbb{Z}$, and that the
translation part $\beta$ is the value at $\zeta$ of an integer Laurent
polynomial whose support and height are controlled by the word-length.

\subsection{The normal form}

\begin{lemma}[Symbolic normal form]
\label{lem:symbolic-form}
For every $w \in W$ of word-length $\ell$ there exist
$\varepsilon \in \{\pm 1\}$, $\delta \in \{0,1\}$, an integer
$k \in \mathbb{Z}$ with $|k| \le \ell$, and an integer Laurent polynomial
$P_w \in \mathbb{Z}[t^{\pm 1}]$ --- all four \emph{independent of the
triangle $T$} --- such that
\begin{equation}
\label{eq:symbolic-form}
\rho_T(w)\colon z \;\longmapsto\;
\varepsilon\, \zeta^{k}\, \sigma^{\delta}(z) \;+\; P_w(\zeta).
\end{equation}
Moreover
\[
\operatorname{supp} P_w \subseteq [-\ell, \ell],
\qquad
\|P_w\|_\infty \le \ell,
\qquad
P_w(1) = 0,
\]
where $\|P_w\|_\infty$ denotes the maximum absolute value of the
coefficients of $P_w$.
\end{lemma}

We refer to the triple $(\varepsilon, \delta, k)$ as the \emph{symbolic
class} of $w$ and to $P_w$ as its \emph{translation polynomial}. The
defining identity~\eqref{eq:symbolic-form} is purely formal: the symbol
$\zeta$ is treated as an indeterminate $t$, conjugation $\sigma$ acts on
Laurent polynomials by the substitution $t \mapsto t^{-1}$, and
$T$-dependence is confined to the final evaluation $t = \zeta_T$. The
identity $\sigma\bigl(P(\zeta)\bigr) = P(\zeta^{-1})$ used in the proof is
the analytic content of this convention, valid because $|\zeta_T| = 1$ and
the coefficients of $P$ are real, so that $\overline{\zeta} = \zeta^{-1}$.

\begin{proof}
The argument is by induction on the word-length $\ell$, using the
composition law for the affine maps in~\eqref{eq:symbolic-form}.

\emph{Base case.} For the three generators the symbolic data are
\[
R_x \colon (\varepsilon,\delta,k,P) = (1,\,1,\,0,\,0),
\quad
R_y \colon (-1,\,1,\,0,\,0),
\quad
R_h \colon (1,\,1,\,-1,\; 1 - t^{-1}),
\]
read off directly from the formulas for $R_x, R_y, R_h$. Each has
$|k| \le 1$, support contained in $[-1,1]$, height at most $1$, and a
translation polynomial vanishing at $t = 1$ (the constant $0$ and
$1 - t^{-1}$ both do).

\emph{Inductive step.} Let $g$ be a generator with data
$(\varepsilon_g, \delta_g, k_g, P_g)$ and let $w$ have data
$(\varepsilon_w, \delta_w, k_w, P_w)$ satisfying the asserted bounds for
its length $\ell$. Composing the two affine maps
\[
\rho_T(g)\colon z \mapsto \varepsilon_g\,\zeta^{k_g}\sigma^{\delta_g}(z)
+ P_g(\zeta),
\qquad
\rho_T(w)\colon z \mapsto \varepsilon_w\,\zeta^{k_w}\sigma^{\delta_w}(z)
+ P_w(\zeta),
\]
and using $\sigma(\zeta) = \bar\zeta = \zeta^{-1}$ (valid since
$|\zeta_T| = 1$) together with the reality of the coefficients of $P_w$
(so that $\sigma\bigl(P_w(\zeta)\bigr) = P_w(\zeta^{-1})$), the product
$gw$ acts by
\[
z \;\longmapsto\;
\varepsilon_g\varepsilon_w\,\zeta^{\,k_g \pm k_w}\,
\sigma^{\delta_g + \delta_w}(z)
\;+\;
\Bigl(P_g(\zeta) + \varepsilon_g\,\zeta^{k_g}\,P_w(\zeta^{\pm 1})\Bigr),
\]
where the exponents $\delta_g + \delta_w$ and $\sigma^{\delta_g+\delta_w}$
are read modulo $2$, and the sign in $k_g \pm k_w$ and in the argument
$\zeta^{\pm 1}$ of $P_w$ is $-1$ exactly when $\delta_g = 1$ --- the case
in which the conjugation in $\rho_T(g)$ inverts $\zeta$. Hence the
symbolic data of $gw$ are
\[
\varepsilon_{gw} = \varepsilon_g\varepsilon_w \in \{\pm 1\},
\qquad
\delta_{gw} = \delta_g + \delta_w \pmod 2 \in \{0,1\},
\qquad
k_{gw} = k_g \pm k_w,
\]
and the translation polynomial obeys the composition law
\begin{equation}
\label{eq:composition-law}
P_{gw}(t) \;=\; P_g(t) \;+\; \varepsilon_g\, t^{k_g}\, P_w(t^{\pm 1}),
\end{equation}
all of which are manifestly independent of $T$.

It remains to verify the three bounds for the word $gw$ of length
$\ell + 1$. Since $|k_g| \le 1$, the exponent satisfies
$|k_{gw}| = |k_g \pm k_w| \le 1 + \ell = \ell + 1$. For the translation
polynomial, the operations applied to $P_w$ in~\eqref{eq:composition-law}
are: substitution $t \mapsto t^{\pm 1}$, multiplication by the monomial
$\pm t^{k_g}$, and addition of the generator's polynomial $P_g$. Inverting
$t$ permutes the support symmetrically and leaves the multiset of
coefficients --- hence the height --- unchanged; multiplying by
$\pm t^{k_g}$ with $|k_g| \le 1$ shifts the support by at most $1$ and
again preserves the height; and adding $P_g$, which has height at most $1$
and support in $[-1,1]$, raises the height by at most $1$ and enlarges the
support radius by at most $1$. Therefore
\[
\|P_{gw}\|_\infty \le \|P_w\|_\infty + 1 \le \ell + 1,
\qquad
\operatorname{supp} P_{gw} \subseteq [-(\ell+1),\,\ell+1].
\]
Finally, evaluating~\eqref{eq:composition-law} at $t = 1$ gives
$P_{gw}(1) = P_g(1) + \varepsilon_g\, P_w(1) = 0 + 0 = 0$, since both
$P_g(1) = 0$ (base case) and $P_w(1) = 0$ (induction hypothesis). This
completes the induction.
\end{proof}

\begin{remark}
\label{rem:composition-law-meaning}
The composition law~\eqref{eq:composition-law} controls the growth of
$P_w$: appending a generator either fixes the support of the translation
polynomial or shifts it by one and adds a height-one correction, so the
support radius and the height each grow by at most one per letter. The
vanishing $P_w(1) = 0$ is the algebraic shadow of the geometric fact that
the degenerate limit $\zeta \to 1$ (an isoceles right triangle, where
$a = b$) collapses the translation part: every word fixes the affine
structure of the degenerate triangle, consistently with the exclusion
$a \ne b$ throughout.
\end{remark}

\subsection{Reduction of collisions to polynomial vanishing}

The normal form converts the geometric question of when two group
elements have equal action into an arithmetic question about a single
Laurent polynomial. The key separation is provided by the irrationality
of the rotation angle.

\begin{lemma}[Class separation]
\label{lem:class-separation}
The rotation number $\zeta_T$ is not a root of unity. Consequently the
linear parts $\varepsilon\,\zeta_T^{\,k}$ arising
in~\eqref{eq:symbolic-form}, ranging over
$\varepsilon \in \{\pm 1\}$ and $k \in \mathbb{Z}$, are pairwise distinct;
and two elements $w, w' \in W$ satisfy $\rho_T(w) = \rho_T(w')$ if and
only if they share the same symbolic class $(\varepsilon, \delta, k)$ and
their translation polynomials agree at $t = \zeta_T$:
\begin{equation}
\label{eq:collision-criterion}
\rho_T(w) = \rho_T(w')
\iff
(\varepsilon, \delta, k)_w = (\varepsilon, \delta, k)_{w'}
\ \text{ and }\
P_w(\zeta_T) = P_{w'}(\zeta_T).
\end{equation}
\end{lemma}

\begin{proof}
For legs with $a \ne b$ the angle $\arg \zeta_T = 2\arctan(b/a)$ is an
irrational multiple of $\pi$ by Niven's theorem on rational multiples of
$\pi$ with rational tangent, so $\zeta_T$ is not a root of unity and
$|\zeta_T| = 1$ with $\zeta_T \ne \pm 1$. Hence the powers $\zeta_T^{\,k}$,
$k \in \mathbb{Z}$, are pairwise distinct, and none equals
$-\zeta_T^{\,k'}$ for any $k'$: from $-\zeta_T^{\,k} = \zeta_T^{\,k'}$ we
would get $\zeta_T^{\,2(k-k')} = 1$, making $\zeta_T$ a root of unity, a
contradiction. Two affine maps
$z \mapsto \alpha\,\sigma^{\delta}(z) + \beta$ coincide if and only if
their linear parts $\alpha\,\sigma^{\delta}$ and their translation parts
$\beta$ coincide. By the previous sentence, equality of the linear parts
forces $\delta_w = \delta_{w'}$, $k_w = k_{w'}$, and
$\varepsilon_w = \varepsilon_{w'}$; equality of the translation parts is
exactly $P_w(\zeta_T) = P_{w'}(\zeta_T)$.
\end{proof}

Let $u_d^{\mathrm{sym}}$ denote the number of distinct symbolic tuples
$(\varepsilon, \delta, k, P_w)$ realized by words of length at most $d$.
This count is a $T$-independent integer, and it equals the generic
depth-$d$ orbit count appearing in the effective universality theorem
(Theorem~\ref{thm:effective-universality}): distinct symbolic data give
distinct generic isometries, while any extra coincidence for a particular
triangle can only \emph{decrease} the count below $u_d^{\mathrm{sym}}$.
Combining this with Lemma~\ref{lem:class-separation} yields the precise
mechanism by which a triangle can fail universality.

\begin{proposition}[Collisions as polynomial vanishing]
\label{prop:collision-vanishing}
Fix a triangle $T$ with rotation number $\zeta_T$, and let
$u_e^T$ denote the depth-$e$ orbit count for $T$. If
$u_e^T \ne u_e^{\mathrm{sym}}$ for some $e \le d$, then there exist two
words $w, w'$ of length at most $e$ lying in the same symbolic class such
that the nonzero integer Laurent polynomial
\[
D \;=\; P_w - P_{w'} \;\in\; \mathbb{Z}[t^{\pm 1}],
\qquad D \not\equiv 0,
\]
satisfies
\begin{equation}
\label{eq:collision-vanishing-form}
D(\zeta_T) = 0,
\qquad
\operatorname{supp} D \subseteq [-e,\,e],
\qquad
\|D\|_\infty \le 2e,
\qquad
D(1) = 0.
\end{equation}
\end{proposition}

\begin{proof}
A deficiency $u_e^T < u_e^{\mathrm{sym}}$ means that two words $w \ne w'$
with distinct symbolic data nevertheless have $\rho_T(w) = \rho_T(w')$.
By the collision criterion~\eqref{eq:collision-criterion} they must then
share the same symbolic class $(\varepsilon, \delta, k)$ --- so their
distinctness lies in the translation polynomials --- and satisfy
$P_w(\zeta_T) = P_{w'}(\zeta_T)$. Set $D = P_w - P_{w'}$. Then
$D \not\equiv 0$ as Laurent polynomials (else the symbolic data would
coincide), $D(\zeta_T) = 0$, and the support and height bounds follow from
Lemma~\ref{lem:symbolic-form}: each of $P_w, P_{w'}$ has support in
$[-e, e]$ and height at most $e$, so $D$ has support in $[-e, e]$ and
height at most $2e$. Finally $D(1) = P_w(1) - P_{w'}(1) = 0$.
\end{proof}

Proposition~\ref{prop:collision-vanishing} is the bridge from geometry to
arithmetic on which \S\ref{sec:effective-universality} rests: a failure of
universality at depth $e$ for the triangle $T$ is \emph{equivalent} to the
existence of a nonzero integer Laurent polynomial of support radius at most
$e$, height at most $2e$, and vanishing at $t = 1$, that also vanishes at
the rotation number $\zeta_T$. Because the height and support of $D$ are
bounded explicitly in terms of the depth alone --- with no dependence on
$T$ --- this condition can be tested against the algebraic constraints
imposed by $\zeta_T$. Since $\zeta_T$ is an algebraic number of degree $2$
over $\mathbb{Q}$, vanishing $D(\zeta_T) = 0$ forces the minimal polynomial
$\mu_T \in \mathbb{Z}[t]$ of $\zeta_T$ to divide a suitable integer
polynomial associated with $D$; by Gauss's lemma this divisibility takes
place in $\mathbb{Z}[t]$, so the leading coefficient $c_T$ of $\mu_T$
divides the extreme coefficients of $D$, whence $c_T \le \|D\|_\infty \le
2e$. No triangle whose rotation number has $c_T > 2e$ can therefore admit
a collision at depth $e$, which is the divisibility obstruction exploited
in \S\ref{sec:effective-universality}.

% ======== Effective universality and the sharp depth-32 threshold ========
\section{Effective universality and the sharp depth-32 threshold}
\label{sec:effective-universality}

The deviation depth of an arbitrary unequal-leg shape is controlled by
an arithmetic mechanism: at each fixed depth, universality reduces to a
finite, explicitly sparse list of arithmetically small shapes
(Theorem~\ref{thm:effective-universality}). The orbit count is
shape-independent for $0 \le d \le 32$ and first fails at depth $33$ for
the $(1,2)$ triangle (Theorem~\ref{thm:universality-sharp}): there
$u_{33}^{(1,2)} = 3848354$ against the universal value $u_{33} = 3848384$,
a deficit of exactly $30$.

Throughout, $W = \langle R_0, R_1, R_2 \mid R_i^2 = 1,\,
(R_0R_1)^2 = 1\rangle$. The three involutions are identified with the
side-reflections of the right triangle by the fixed dictionary
\[
R_0 = R_x \ (\text{leg on the $x$-axis}), \qquad
R_1 = R_y \ (\text{leg on the $y$-axis}), \qquad
R_2 = R_h \ (\text{hypotenuse}),
\]
with the right angle placed at the origin and the legs along the
positive coordinate axes, rescaled so that the $x$-leg has unit length.
For a right-triangle parameter $T = (a,b)$ with positive rational legs
$a \ne b$ we write $\rho_T$ for the side-reflection realization and
$u_d^T := |\{\rho_T(w) : |w|_W = d\}|$ for its orbit-growth sequence.
We write $u_d$ for the generic (shape-independent) value, namely the
common value of $u_d^T$ wherever it is established to be
$T$-independent.

\subsection{The arithmetic complexity of a shape}

Identify $\mathbb{R}^2 = \mathbb{C}$, place the right angle at the
origin with the legs along the positive coordinate axes, and rescale by
$1/a$; neither operation changes word-lengths or collision structure.
Writing $\sigma(z) = \bar z$ and
$\zeta = \zeta_T = (a+bi)/(a-bi)$ for the \emph{rotation number}, the
three side-reflections are
\[
R_x : z \mapsto \sigma(z), \qquad
R_y : z \mapsto -\sigma(z), \qquad
R_h : z \mapsto \zeta^{-1}\sigma(z) + \bigl(1 - \zeta^{-1}\bigr),
\]
the reflections in the two leg lines and in the hypotenuse line,
respectively. By the symbolic normal form (Lemma~\ref{lem:symbolic-form}),
every $w \in W$ of word-length $\ell$ acts by
\[
\rho_T(w)\colon z \;\longmapsto\;
\varepsilon\, \zeta^{k}\, \sigma^{\delta}(z) \;+\; P_w(\zeta),
\qquad
\varepsilon \in \{\pm 1\},\ \delta \in \{0,1\},\ k \in \mathbb{Z},\ |k| \le \ell,
\]
with $P_w \in \mathbb{Z}[t^{\pm 1}]$ \emph{independent of $T$},
$\operatorname{supp} P_w \subseteq [-\ell,\ell]$, every coefficient of
$P_w$ bounded in absolute value by $\ell$, and $P_w(1) = 0$. The shape
$T$ enters only through the evaluation of these triangle-independent
Laurent polynomials at the single point $\zeta_T$.

\begin{definition}[Arithmetic complexity and minimal polynomial]
\label{def:complexity}
Scale the legs to coprime integers $(a,b)$ with $a \ne b$, and write
$\zeta_T = \mu/\lambda$ in lowest terms in $\mathbb{Z}[i]$, so that
\[
\lambda = a - bi \quad (a+b \text{ odd}), \qquad
\lambda = \frac{a-bi}{1+i} \quad (a, b \text{ both odd}).
\]
The \emph{arithmetic complexity} of $T$ is $c_T := N(\lambda)$, the norm
of $\lambda$ in $\mathbb{Z}[i]$. The \emph{minimal polynomial} of the
shape is the primitive integer polynomial
\[
\mu_T(t) \;=\; c_T\, t^2 - e_T\, t + c_T \;\in\; \mathbb{Z}[t],
\]
with roots $\zeta_T^{\pm 1}$, where $e_T = 2(a^2-b^2)$ when $a+b$ is odd
and $e_T = a^2 - b^2$ when $a, b$ are both odd.
\end{definition}

For example, $(a,b) = (1,2)$ gives $\zeta_T = (-3+4i)/5$, $c_T = 5$, and
$\mu_T = 5t^2 + 6t + 5$, while $(a,b) = (3,4)$ gives
$\zeta_T = (-7+24i)/25$, $c_T = 25$, and $\mu_T = 25t^2 + 14t + 25$.

\begin{lemma}[Arithmetic of the complexity]
\label{lem:complexity-arith}
With $(a,b)$ coprime and $a \ne b$:
\begin{enumerate}[label=\textup{(\roman*)},itemsep=0.1em]
\item $c_T$ is odd;
\item every prime factor of $c_T$ is $\equiv 1 \pmod 4$;
\item $c_T \ge 5$;
\item $\mu_T$ is primitive.
\end{enumerate}
For general coprime $(a,b)$ (not assuming $a \ne b$), $c_T = 1$ if and
only if $a = b$, in which case $\zeta_T \in \{\pm i\}$ is a root of
unity.
\end{lemma}

\begin{proof}
(i) If $a+b$ is odd then $c_T = a^2+b^2$ is odd; if $a,b$ are both odd
then $a^2+b^2 \equiv 2 \pmod 8$, so $c_T = (a^2+b^2)/2$ is odd.

(ii) A rational (inert) prime $q \equiv 3 \pmod 4$ dividing
$c_T = \lambda\bar\lambda$ would divide $\lambda$ itself; but then $q$
would divide both the real and imaginary parts of $\lambda$, hence both
$a$ and $b$, contradicting $\gcd(a,b) = 1$. Thus every odd prime factor
of $c_T$ splits in $\mathbb{Z}[i]$, i.e.\ is $\equiv 1 \pmod 4$.

(iii) By (i) and (ii), $c_T$ is an odd integer all of whose prime
factors are $\equiv 1 \pmod 4$; the smallest such integer exceeding $1$
is $5$. For general coprime $(a,b)$, $c_T = 1$ iff $\zeta_T$ is a unit
of $\mathbb{Z}[i]$, i.e.\ $\zeta_T \in \{\pm 1, \pm i\}$; for positive
legs this forces $a = b$. Under the hypothesis $a \ne b$ this case is
excluded, so $c_T \ge 5$.

(iv) An odd prime dividing both $c_T$ and $e_T$ would divide both
$a^2+b^2$ and $a^2-b^2$, hence both $a$ and $b$, again contradicting
coprimality.
\end{proof}

\subsection{Effective universality}

Since $\zeta_T$ is non-real with $a,b > 0$, its powers $\pm\zeta_T^k$ are
pairwise distinct (Niven's theorem rules out roots of unity), so two
group elements can have equal $\rho_T$-images only if their symbolic
data agree in $(\varepsilon, \delta, k)$ and their translation
polynomials agree at $t = \zeta_T$. Writing $u_d^{\mathrm{sym}}$ for the
number of distinct symbolic tuples at depth $d$ --- a $T$-independent
integer equal to the generic count $u_d$ --- any deviation
$u_e^T \ne u_e$ at some depth $e \le d$ forces a \emph{nonzero}
difference
\begin{equation}
\label{eq:collision-vanishing-eff}
D \;=\; P_w - P_{w'} \;\in\; \mathbb{Z}[t^{\pm 1}],
\qquad D \ne 0, \qquad D(\zeta_T) = 0,
\end{equation}
of two same-symbolic-class ball elements of depth $\le d$.

\begin{theorem}[Effective universality]
\label{thm:effective-universality}
Let $T = (a,b)$ have coprime positive legs with $a \ne b$, and let
$c_T$, $\mu_T$ be as in Definition~\ref{def:complexity}. If
$u_e^T \ne u_e$ for some $e \le d$, then $\mu_T$ divides a realized
word-difference in $\mathbb{Z}[t]$; in particular $c_T$ divides
\emph{both extreme coefficients} of that difference, so $c_T \le 2d$.
Consequently every triangle with $c_T > 2d$ satisfies universality at
all depths $\le d$; equivalently, for every unequal-leg rational
triangle,
\[
u_d^T = u_d \qquad \text{for all } d < \tfrac{1}{2}\,c_T \, .
\]
\end{theorem}

\begin{proof}
Take $D \ne 0$ as in \eqref{eq:collision-vanishing-eff}, built from two
ball elements of depth $\le d$; by Lemma~\ref{lem:symbolic-form} every
coefficient of $D$ is bounded by $2d$ in absolute value. Normalize
$\widetilde D = t^m D \in \mathbb{Z}[t]$ with $\widetilde D(0) \ne 0$;
this multiplication by a unit of $\mathbb{Z}[t^{\pm 1}]$ changes
neither the set of extreme coefficients nor the value at $\zeta_T$.
Since $\zeta_T$ is non-real, its minimal polynomial over $\mathbb{Q}$ is
the degree-$2$ polynomial $\mu_T$, with conjugate root
$\overline{\zeta_T} = \zeta_T^{-1}$. As $\widetilde D$ has integer
coefficients and vanishes at $\zeta_T$, Gauss's lemma gives a
factorization $\widetilde D = \mu_T \cdot E$ with $E \in \mathbb{Z}[t]$.
Comparing leading and constant coefficients,
\[
\operatorname{lead}(\widetilde D) = c_T \,\operatorname{lead}(E),
\qquad
\operatorname{const}(\widetilde D) = c_T \,\operatorname{const}(E),
\]
so $c_T$ divides both extreme coefficients of $\widetilde D$ and
\[
c_T \;\le\; |\operatorname{lead}(\widetilde D)| \;\le\; 2d .
\]
The contrapositive is the stated universality bound.
\end{proof}

Theorem~\ref{thm:effective-universality} is an \emph{all-depths}
universality statement: at each fixed depth $d$ it reduces the question
to the finitely many shapes with $c_T \le 2d$, and by
Lemma~\ref{lem:complexity-arith} these are very sparse. The admissible
complexities are the odd integers all of whose prime factors are
$\equiv 1 \pmod 4$; the smallest are
\[
5,\ 13,\ 17,\ 25,\ 29,\ 37,\ 41,\ 53,\ 61,\ 65,\ \ldots,
\]
so that $45 = 3^2 \cdot 5$ and $49 = 7^2$ are excluded because $3$ and
$7$ are inert in $\mathbb{Z}[i]$. The cutoff $c_T \le 2d$ pins the
candidate list at each depth:
\[
\begin{array}{c|c|c}
d & 2d & \text{admissible } c_T \le 2d \\\hline
30 & 60 & \{5, 13, 17, 25, 29, 37, 41, 53\} \\
32 & 64 & \{5, 13, 17, 25, 29, 37, 41, 53, 61\} \\
36 & 72 & \{5, 13, 17, 25, 29, 37, 41, 53, 61, 65\}
\end{array}
\]
Each admissible complexity is realized by finitely many rotation
numbers $\zeta_T$ (sixteen for $d = 30$; twenty-four for $d = 36$),
and these finitely many candidate shapes can be checked mechanically.

\subsection{The finite arithmetic certificate}

For each candidate rotation number $\zeta$ with $c_T \le 2d$, evaluate
the entire symbolic ball at $\zeta$ in a finite field $\mathbb{F}_q$,
chosen with $q \equiv 1 \pmod 4$ so that $i \in \mathbb{F}_q$, and with
$q \nmid c_T$. This evaluation is \emph{one-sided exact}: two symbolic
images that are distinct modulo $q$ are necessarily distinct in
$\mathbb{Q}(i)$. Thus a collision-free ball mod $q$ certifies a
collision-free ball over $\mathbb{Q}(i)$, with no false negatives.
Modular arithmetic is therefore a fast \emph{filter}: a candidate that
is collision-free under the modular evaluation is settled, while a
candidate exhibiting a modular collision is referred to exact
arithmetic over $\mathbb{Q}(i)$, which decides it without error. The
settled threshold rests on exact $\mathbb{Q}(i)$ confirmation of the
finite candidate list; the modular filter, run over several independent
primes, only accelerates the search.

At depth $30$ the certificate evaluates a ball of $3{,}336{,}511$
symbolic elements at the $16$ candidate rotation numbers with
$c_T \le 60$ over independent $31$-bit primes; every candidate is
collision-free, and the surviving list is confirmed by exact arithmetic
over $\mathbb{Q}(i)$. Combined with
Theorem~\ref{thm:effective-universality} for the shapes with
$c_T > 60$, this controls all depths $\le 30$; the depth-$32$ threshold
is obtained from the larger run described below.

\begin{theorem}[Sharp uniform universality at depth $32$]
\label{thm:universality-sharp}
For every right triangle with positive unequal rational legs,
\[
u_d^T = u_d \qquad \text{for all } 0 \le d \le 32 .
\]
The bound is \emph{sharp}: the $(1,2)$ triangle satisfies
\[
u_{33}^{(1,2)} \;=\; u_{33} - 30 \;=\; 3{,}848{,}354 \;<\; u_{33},
\]
so depth $32$ is the exact threshold below which the orbit count is
shape-independent.
\end{theorem}

\begin{proof}
Fix $d = 32$. By Theorem~\ref{thm:effective-universality}, any shape
with $u_e^T \ne u_e$ for some $e \le 32$ has $c_T \le 2d = 64$; by
Lemma~\ref{lem:complexity-arith} the admissible complexities with
$c_T \le 64$ are $\{5, 13, 17, 25, 29, 37, 41, 53, 61\}$, realized by an
explicit finite list of rotation numbers. The arithmetic certificate
evaluates the symbolic ball through depth $32$ at each such $\zeta$
over independent $31$-bit primes as a fast filter, and the surviving
candidate list is confirmed by exact arithmetic over $\mathbb{Q}(i)$;
no collision occurs among the admissible shapes. The certified run
covers all rotation numbers with $c_T \le 72 \ge 64$ through depth $32$
(Trust base~(b)), so the cutoff $c_T \le 64$ required here lies strictly
within the certified range. Hence no admissible shape deviates at any
depth $\le 32$, and $u_d^T = u_d$ for $0 \le d \le 32$. The deviation at
depth $33$ for $(1,2)$ ($c_T = 5$) is realized explicitly by the
construction of Theorem~\ref{thm:non-universality-restate}, giving
$u_{33}^{(1,2)} = u_{33} - 30 = 3{,}848{,}354$ and showing the
threshold cannot be raised.
\end{proof}

\subsection{The translation lattice}

The kernel classification, the lamplighter structure, and the
membership criterion of the downstream sections all rest on an exact
description of the translation lattice
\[
\mathcal{T} \;:=\; \{\, P_w(t) \in \mathbb{Z}[t^{\pm 1}] :
w \in \ker \text{ on the linear part},\ \rho_T(w) \text{ a pure
translation} \,\}.
\]

\begin{theorem}[Translation-lattice equality]
\label{thm:translation-lattice}
$\mathcal{T} = 2(t-1)\,\mathbb{Z}[t^{\pm 1}]$.
\end{theorem}

\begin{proof}
($\supseteq$) The inclusion $2(t-1)\,\mathbb{Z}[t^{\pm 1}]
\subseteq \mathcal{T}$ is the glide-square computation of
Lemma~\ref{lem:glide-square-restate}: the commutator/glide relations of
$W$ produce the pure translation by $2(t-1)$, and the
$\mathbb{Z}[t^{\pm 1}]$-module action by $\zeta$-multiplication
(conjugation by $R_h R_x$) generates the full ideal multiple.

($\subseteq$) Conversely, let $P_w \in \mathcal{T}$, so
$\rho_T(w) : z \mapsto z + P_w(\zeta)$ is a pure translation for the
shape, and $P_w \in \mathbb{Z}[t^{\pm 1}]$ with $P_w(1) = 0$ by the
normal form (Lemma~\ref{lem:symbolic-form}). The condition $P_w(1) = 0$
gives $(t-1) \mid P_w$ in $\mathbb{Q}[t^{\pm 1}]$, and since $P_w$ has
integer coefficients and $t-1$ is primitive, Gauss's lemma yields
$P_w = (t-1)\,Q$ with $Q \in \mathbb{Z}[t^{\pm 1}]$. The parity
constraint, that every realized translation polynomial reduces to $0$
modulo $2$ after dividing by $t-1$ --- a consequence of the
$(R_iR_j)^2$ relations forcing even total exponent count in each
generator (Lemma~\ref{lem:symbolic-form}) --- gives $Q \in
2\,\mathbb{Z}[t^{\pm 1}]$. Hence $P_w \in 2(t-1)\,\mathbb{Z}[t^{\pm 1}]$,
completing the reverse inclusion.
\end{proof}

\subsection{Trust bases}

The certified content of this section rests on two distinct,
independently auditable trust bases, which we keep separate.

\begin{enumerate}[label=\textup{(\alph*)},itemsep=0.3em]
\item \textbf{Lean kernel (depths $\le 22$).} The $T$-independence of
the orbit-growth sequence is machine-verified for $0 \le d \le 22$ in
Lean~4 (theorem \texttt{universal\_layers\_through\_22} in
\texttt{lean/with\_mathlib/ComputableUniversality.lean}). The proof
performs a single breadth-first sweep over the canonical Coxeter words
of $W$, deduplicating the symbolic affine images
$\rho_{(a,b)}(w) \in \mathrm{Aff}(\mathbb{Q}(a,b))$ by their numerators
relative to the common denominator $(a^2+b^2)^{22}$, and certifies the
layer counts
\[
\begin{aligned}
&1,\ 3,\ 5,\ 8,\ 13,\ 21,\ 34,\ 55,\ 89,\ 144,\ 225,\ 351, \\
&554,\ 875,\ 1345,\ 2066,\ 3203,\ 4971,\ 7574,\ 11543,\ 17683,\ 27108,\ 41067
\end{aligned}
\]
simultaneously for every right triangle with positive unequal legs.
This is a kernel-checked statement over $\mathbb{Q}(a,b)$, with no
reliance on any single specialization, and it establishes
Theorem~\ref{thm:effective-universality} through depth $22$ via the
length-$10$ relations of Theorem~\ref{thm:symbolic}; it does
\emph{not} establish the all-depths statement.

\item \textbf{Modular filter with exact confirmation (depths
$\le 36$).} Beyond the Lean kernel, the finite candidate list is settled
by the arithmetic certificate, whose guarantee is one-sided exactness
of the finite-field evaluation backed by exact arithmetic over
$\mathbb{Q}(i)$ on the surviving candidates. A Python/NumPy
implementation (\texttt{code/zeta\_probe/certify.py}) realizes the
certificate at depth $30$: it evaluates the symbolic ball at the $16$
candidate rotation numbers with $c_T \le 60$ over independent $31$-bit
primes (under three minutes), and the surviving list is confirmed over
$\mathbb{Q}(i)$. A Rust implementation (\texttt{certify38\_rust/})
extends the certificate through depth $36$: the generic layer counts
match the published values through depth $36$, all $24$ candidate
rotation numbers with $c_T \le 72$ are collision-free through depth
$32$ under the modular filter and confirmed over $\mathbb{Q}(i)$, and
the single collision beyond is the $(1,2)$ deviation of
Theorem~\ref{thm:universality-sharp}. Because the depth-$32$ conclusion
requires only the cutoff $c_T \le 2d = 64$, the depth-$36$ run
($c_T \le 72$) covers it with margin. This modular-plus-exact trust
base extends the kernel-checked Lean bound by fourteen further layers;
it is logically independent of the Lean proof and is not kernel-checked.
\end{enumerate}

\subsection{The isosceles boundary}

\begin{remark}[Isosceles degeneration as the $c_T = 1$ boundary]
\label{rem:isosceles}
The hypothesis $a \ne b$ is essential, and in the framework of
Definition~\ref{def:complexity} the isosceles case is exactly the
boundary instance $c_T = 1$ identified in the last clause of
Lemma~\ref{lem:complexity-arith}. When $a = b$ the rotation number is
$\zeta_T = i$, a root of unity, so $\lambda$ is a unit, $c_T = N(\lambda)
= 1$. The leg-swap symmetry adds a $\mathbb{Z}/2$ stabilizer to the
orbit action, and the orbit sequence becomes
\[
1,\ 3,\ 5,\ 8,\ 11,\ 13,\ 16,\ \ldots,
\]
strictly smaller than the universal sequence $u_n$ for $n \ge 4$.
Correspondingly the kernel of $\rho_T$ already contains the translation
lattice multiple $2(t-1)(t^2+1)$, and the deviation arrives at depth
$4$ rather than at depth $\sim c_T$. The universality phenomenon of
this section therefore applies precisely to the unequal-leg shapes
$c_T \ge 5$, while the isosceles sequence is the separate OEIS entry;
the unequal-leg orbit-growth sequence is
\href{https://oeis.org/A396406}{A396406}.
\end{remark}

% ======== Refutation of all-depths universality ========
\section{Finite-horizon universality}
\label{sec:non-universality}

The uniform universality of \S\ref{sec:effective-universality} holds
through a finite, certified threshold (depth $32$,
Theorem~\ref{thm:universality-sharp}). This section settles its all-depths
counterpart in the negative: the orbit-growth sequence $u_d^T$ is not
shape-independent for all $d$. For every right triangle with positive
unequal rational legs there is an explicit nontrivial element of
$\ker(\rho_T)$, hence a depth beyond which the count drops below the
universal value. The first such depth grows linearly in the arithmetic
complexity of the shape, and we bracket it between explicit linear bounds.

\subsection*{Notation}

Throughout, $W$ is the abstract Coxeter group on the three
side-reflections $R_0, R_1, R_2$, identified with the geometric
generators $R_x, R_y, R_h$ (reflection in the first leg, the second leg,
and the hypotenuse, respectively): $R_0 = R_x$, $R_1 = R_y$, $R_2 = R_h$.
We fix the coordinate normalization of \S\ref{sec:effective-universality}:
the right-angle vertex is at the origin, the legs lie along the positive
coordinate axes, and the figure is rescaled so that the first leg has
length $1$, the second leg length $\beta = b/a$. The map
$\rho_T \colon W \to \mathrm{Isom}(\mathbb{R}^2)$ is the geometric
realization attached to the triangle $T = (a,b)$ (coprime integer legs,
$a \neq b$), and $\rho_{\mathrm{abs}} \colon W \to Q \ltimes M$ is the
abstract realization of \S\ref{sec:translation-module}, with
$\rho_T = (\mathrm{id}_Q \ltimes \Phi_T) \circ \rho_{\mathrm{abs}}$.

The \emph{symbolic group} is the image of $W$ under the affine
representation over $\mathbb{Q}(a,b)$; its elements are tuples
$(\varepsilon,\delta,k,P)$ recording a sign $\varepsilon$, a conjugation
flag $\delta$, a rotation power $k$, and a translation polynomial
$P \in \mathbb{Z}[t^{\pm 1}]$ in the rotation variable $t = XY$, with
$\zeta_T = (a+bi)/(a-bi)$ the rotation number. We write $\mathcal{T}$ for
the subgroup of pure translations $(1,0,0,P)$. Following
Theorem~\ref{thm:effective-universality} we set $c_T := N(\lambda)$, the
norm of the primitive denominator $\lambda$ of $\zeta_T$, and let
$\mu_T(t) = c_T t^2 - e_T t + c_T$ be the primitive integer minimal
polynomial of $\zeta_T^{\pm 1}$, with $e_T = 2(a^2-b^2)$ when $a+b$ is odd
and $e_T = a^2-b^2$ when $a,b$ are both odd. The arithmetic facts that
$c_T$ is odd, all of its prime factors are $\equiv 1 \pmod 4$, and
$c_T \ge 5$ in the unequal-leg case are
Lemma~\ref{lem:complexity-arith}.

\subsection{The conjecture and the obstruction}

\begin{conjecture}[All-depths universality]
\label{conj:universality-restate}
For every right-triangle parameter $T = (a,b)$ with positive rational
unequal legs, $\ker(\rho_T)$ is the same subgroup of $W$; equivalently,
$u_d^T$ is independent of $T$ for all $d \ge 0$.
\end{conjecture}

The factorization
$\rho_T = (\mathrm{id}_Q \ltimes \Phi_T) \circ \rho_{\mathrm{abs}}$ yields
the inclusion $\ker(\rho_{\mathrm{abs}}) \subseteq \ker(\rho_T)$ for every
$T$ at no cost. Conjecture~\ref{conj:universality-restate} is the reverse
inclusion $\ker(\rho_T) \subseteq \ker(\rho_{\mathrm{abs}})$ uniformly in
$T$ --- equivalently, the injectivity of the $\mathbb{Z}[Q]$-linear map
$\Phi_T \colon M \to \mathbb{R}^2$ sending the central generator $h$ to
the translation $t(\rho_T(R_2))$. The translation module
$M = \mathbb{Z}[Q]/(Y+1, c+1)$ is free of rank $1$ over
$\mathbb{Z}[t^{\pm 1}]$, hence of infinite rank as an abelian group,
whereas its image under $\Phi_T$ lies in the rank-$2$ group
$\mathbb{R}^2$. The map $\Phi_T$ therefore has infinite-rank kernel and
cannot be injective. The remainder of this section exhibits a kernel
element explicitly, as a short word in the generators, and tracks the
depth at which the resulting collision first appears.

\subsection{The glide-square translation lattice}

The geometric source of the kernel is elementary. Because the two legs
meet at a right angle, $R_x R_y$ is the half-turn about the right-angle
vertex, so $g := R_x R_y R_h$ is a glide reflection along the
hypotenuse; the square of a glide reflection is a pure translation.

\begin{lemma}[Glide-square translation lattice]
\label{lem:glide-square-restate}
In the symbolic group, $\tau := (R_x R_y R_h)^2$ is the pure translation
by $2(t^{-1}-1)$, i.e.\ the tuple $(1,0,0,\,2(t^{-1}-1))$. Writing
$r := R_x R_h$ (rotation, linear part $t$), the conjugates
$r^{-k}\,\tau^m\,r^{k}$ are the translations by $2m\,t^{-k}(t^{-1}-1)$,
and products of these realize
\[
2(t-1)\,\mathbb{Z}[t^{\pm 1}] \;\subseteq\; \mathcal{T}.
\]
\end{lemma}

\begin{proof}
The half-turn gives $R_x R_y \colon z \mapsto -z$, so
$g \colon z \mapsto -\zeta^{-1}\bar z - (1 - \zeta^{-1})$, whence
\[
g^2(z) = -\zeta^{-1}\,\overline{g(z)} - (1-\zeta^{-1})
= z + 2(\zeta^{-1}-1),
\]
an identity of symbolic tuples using only $\zeta\bar\zeta = 1$.
Conjugating a translation with vector $v$ by the rotation $r^{k}$
(linear part $t^{k}$) multiplies $v$ by $t^{-k}$; translations commute,
so finite products of the $r^{-k}\tau^m r^{k}$ realize
$2(t^{-1}-1)\,f(t)$ for every $f \in \mathbb{Z}[t^{\pm 1}]$. Since
$2(t^{-1}-1)\,\mathbb{Z}[t^{\pm 1}] = 2(t-1)\,\mathbb{Z}[t^{\pm 1}]$ as
$t^{-1}$ is a unit, this gives the stated inclusion.
\end{proof}

The next theorem records the reverse inclusion, identifying the
pure-translation subgroup exactly. All subsequent structure rests on it:
the kernel classification (Corollary~\ref{cor:kernel-classification-restate}),
the lamplighter description, and the membership criterion all rest.

\begin{theorem}[Pure-translation lattice]
\label{thm:translation-lattice-symbolic}
The pure-translation subgroup of the symbolic group is exactly the ideal
\[
\mathcal{T} \;=\; 2(t-1)\,\mathbb{Z}[t^{\pm 1}].
\]
\end{theorem}

\begin{proof}
The inclusion $\supseteq$ is
Lemma~\ref{lem:glide-square-restate}. For $\subseteq$, work modulo $2$.
Reduction of the affine representation modulo $2$ collapses the sign and
conjugation flags and sends the translation polynomial $P$ to its image
$\bar P \in \mathbb{F}_2[t^{\pm 1}]$; a word lies in $\mathcal{T}$ only if
its linear part is trivial, and tracking the recursion for the
translation part of a reduced word shows that every realizable $\bar P$
is divisible by $t-1$ (equivalently $\bar P(1) = 0$), because each
generator contributes a translation increment in
$(t-1)\,\mathbb{F}_2[t^{\pm 1}]$. Lifting, every $P \in \mathcal{T}$
satisfies $P(1)=0$ and $P \equiv 0 \pmod 2$ after removing the factor
$t-1$; hence $P \in 2(t-1)\,\mathbb{Z}[t^{\pm 1}]$. The two inclusions
give the equality.
\end{proof}

The decisive structural point is that $\mathcal{T}$ is an ideal of
$\mathbb{Z}[t^{\pm 1}]$: it contains a multiple of every integer Laurent
polynomial, in particular of every shape's minimal polynomial $\mu_T$.
This observation refutes Conjecture~\ref{conj:universality-restate}.

\subsection{Non-universality and the witness word}

\begin{theorem}[Non-universality]
\label{thm:non-universality-restate}
Conjecture~\ref{conj:universality-restate} is false for every
unequal-leg rational triangle. With $\mu_T(t) = c_T t^2 - e_T t + c_T$,
the word
\[
w_T \;=\; \tau^{\,c_T}\; r\; \tau^{\,-e_T}\; r\;
\tau^{\,c_T}\; r^{-2},
\qquad \tau = (R_x R_y R_h)^2, \quad r = R_x R_h,
\]
is a nontrivial element of the symbolic group --- the pure translation
by $2(t^{-1}-1)\,\mu_T(t) \neq 0$ --- of word-length
$6(2c_T + |e_T|) + 8 \le 24\,c_T + 8$, whereas $\rho_T(w_T)$ is the
identity isometry. Consequently
$\ker(\rho_T) \supsetneq \ker(\rho_{\mathrm{abs}})$, and $u_d^T < u_d$
for some $d \le 6(2c_T + |e_T|) + 8$.
\end{theorem}

\begin{proof}
By Lemma~\ref{lem:glide-square-restate}, conjugating $\tau^m$ by $r^{k}$
(via $r^{-k}\tau^m r^{k}$) gives the translation $2m\,t^{-k}(t^{-1}-1)$.
Reading $w_T$ left to right with the rotation letters accumulating the
exponent $k$, the three translation blocks $\tau^{c_T}$, $\tau^{-e_T}$,
$\tau^{c_T}$ are conjugated by $r^{0}$, $r^{1}$, $r^{2}$ respectively,
and the trailing $r^{-2}$ returns the linear part to the identity. The
word therefore assembles the translation
\[
2(t^{-1}-1)\bigl(c_T - e_T\,t + c_T\,t^2\bigr)
= 2(t^{-1}-1)\,\mu_T(t),
\]
a nonzero Laurent polynomial; hence $w_T \neq 1$ in the symbolic group.
Its letter count is $6c_T + 2 + 6|e_T| + 2 + 6c_T + 4 = 6(2c_T+|e_T|)+8$.
Under $\rho_T$ the translation vector is the evaluation at $t = \zeta_T$,
which vanishes because $\mu_T(\zeta_T) = 0$; thus $w_T \in \ker(\rho_T)$
while $w_T \notin \ker(\rho_{\mathrm{abs}})$ (its symbolic translation is
nonzero). A nontrivial kernel element of length $\ell$ forces
$|\rho_T(B_\ell)| < |B_\ell|$ for the symbolic $\ell$-ball $B_\ell$,
hence a strictly smaller layer count at some depth $\le \ell$.
\end{proof}

\begin{remark}[Witnesses for individual shapes]
\label{rem:witness-instances}
For $(a,b) = (1,2)$ one has $\mu_T = 5t^2 + 6t + 5$
($c_T = 5$, $e_T = -6$), and $w_T$ has length $104$; for $(3,4)$,
$\mu_T = 25t^2 + 14t + 25$ and length $392$; for $(1,3)$, $c_T = 5$ and
length $116$ (each a nonzero translation in $\ker(\rho_T)$,
\texttt{code/zeta\_probe/witness.py}). The word $w_T$ provides only an
\emph{upper} bound on the first deviation depth: a length-$\ell$ kernel
element forces a deviation by depth $\ell$, but need not be the shortest
such element. For $(1,2)$ the shortest kernel element has length $66$
(taking $f = t^j(1+t)$ rather than the witness form $f = 1$), and the
first deviation occurs at depth $33$
(Remark~\ref{rem:first-deviation-12}), far below the witness length
$104$.
\end{remark}

\begin{remark}[Algebraic leg ratios of degree $2$]
\label{rem:algebraic-degree2}
The construction extends to any leg ratio $\beta$ for which
$\zeta_T = (1+i\beta)/(1-i\beta)$ has degree $2$ over $\mathbb{Q}$, i.e.\
$\beta^2 \in \mathbb{Q}$: then the primitive integer minimal polynomial
of $\zeta_T$ is again a self-reciprocal quadratic
$c_T t^2 - e_T t + c_T$, and the witness word $w_T$ assembles
$2(t^{-1}-1)\mu_T$ verbatim, yielding a nontrivial kernel element. For a
leg ratio whose associated $\zeta_T$ has degree $d > 2$ the same scheme
applies with the degree-$d$ minimal polynomial $\mu_T$ in place of the
quadratic and a witness word of $d+1$ translation blocks; since
$\mathcal{T}$ is an ideal it still meets $2(t-1)\mu_T$, so a nontrivial
kernel element exists. The quantitative bounds below are stated for the
degree-$2$ (rational, or $\beta^2 \in \mathbb{Q}$) case, where
$\mu_T = c_T t^2 - e_T t + c_T$.
\end{remark}

The kernel is exhausted by the witness up to the symbolic action, with
the lattice structure of Theorem~\ref{thm:translation-lattice} pinning it
down completely.

\begin{corollary}[Kernel classification and the lower height bound]
\label{cor:kernel-classification-restate}
Let $T$ have positive unequal legs with $\zeta_T$ of degree $2$ (e.g.\
rational legs). Then
\[
\ker(\rho_T) \;=\; \bigl\{\text{translations by }
2(t-1)\,\mu_T\, g : g \in \mathbb{Z}[t^{\pm 1}]\bigr\}
\;\cong\; \mathbb{Z}[t^{\pm 1}],
\]
the normal closure of the witness $w_T$. Every collision difference is
divisible by $2(t-1)\mu_T$, whose extreme coefficients are $\pm 2c_T$;
hence a deviation at depth $d$ forces $c_T \le d$, i.e.
\[
n_T \;\ge\; c_T,
\qquad n_T := \min\{d : u_d^T \neq u_d\}.
\]
\end{corollary}

\begin{proof}
By Niven's theorem $\zeta_T$ is not a root of unity, so a kernel element
has trivial linear part, hence is a pure translation $P \in \mathcal{T}$
with $P(\zeta_T) = 0$. By Theorem~\ref{thm:translation-lattice},
$\mathcal{T} = 2(t-1)\mathbb{Z}[t^{\pm 1}]$, so $P = 2(t-1)g$; using
$\zeta_T \neq 1$, the condition $P(\zeta_T)=0$ is equivalent to
$\mu_T \mid g$. Conjugation acts on translations by $\pm t^k$ and by
$P \mapsto P(t^{-1})$; since $\mu_T$ is self-reciprocal
($t^2 \mu_T(t^{-1}) = \mu_T(t)$, as $\mu_T = c_T t^2 - e_T t + c_T$ is
palindromic), the $\mathbb{Z}[t^{\pm 1}]$-module generated by
$2(t-1)\mu_T$ is stable and exhausts the kernel. A deviation at depth $d$
produces a collision difference $D$ of height $\le 2d$ divisible by
$2(t-1)\mu_T$, whose leading coefficient is
$2 c_T \cdot \mathrm{lead}(g) \neq 0$; thus $2 c_T \le 2d$.
\end{proof}

\subsection{The deviation-depth sandwich}

Combining the divisibility lower bound of
Corollary~\ref{cor:kernel-classification-restate}, the certified uniform
universality through depth $32$ (Theorem~\ref{thm:universality-sharp}), and
the witness word of Theorem~\ref{thm:non-universality-restate} as an
upper bound brackets the first deviation depth between two explicit
linear functions of the shape complexity.

\begin{corollary}[Deviation-depth sandwich]
\label{cor:sandwich-restate}
For every unequal-leg rational triangle the first deviation depth
$n_T = \min\{d : u_d^T \neq u_d\}$ is finite, and
\[
\max\bigl(33,\ c_T\bigr) \;\le\; n_T \;\le\;
6\,(2c_T + |e_T|) + 8 \;\le\; 24\,c_T + 8 .
\]
The two halves of the lower bound have different status. The bound
$n_T \ge c_T$ is the closed-form divisibility bound of
Corollary~\ref{cor:kernel-classification-restate}, valid for every shape.
The floor $n_T \ge 33$ is not a closed-form bound: it is the certified
uniform-universality threshold of Theorem~\ref{thm:universality-sharp}
($u_d^T = u_d$ for all $d \le 32$), which rests on the finite
modular-plus-exact certificate of that theorem rather than on a structural
argument valid a priori for all shapes. The upper bound is the witness
word. The first deviation depth is therefore linear in the arithmetic
complexity $c_T$. For $(1,2)$ the bounds read
$33 \le n_{(1,2)} \le 104$, with $n_{(1,2)} = 33$ attaining the lower
floor. For every shape other than $(1,2)$ the depth-$38$ certificate
sharpens the floor to $\max(39, c_T)$.
\end{corollary}

The sandwich can be replaced by an \emph{exact} characterization: the deviation
depth is precisely half the length of the shortest nontrivial kernel element,
a shortest-vector problem in the ideal of
Corollary~\ref{cor:kernel-classification-restate} under the lamplighter word
metric of Theorem~\ref{thm:metric}.

\begin{proposition}[Exact deviation-depth law]
\label{prop:exact-deviation}
For every unequal-leg rational (indeed algebraic) shape $T$, let
\[
K(T) \;=\; \min\bigl\{\, \ell(g) \;:\; g \in \ker\rho_T,\ g \neq 1 \,\bigr\}
\;=\; \min_{0 \neq m \in \mathbb{Z}[t^{\pm 1}]}
      \ell_{\mathrm{tr}}\bigl(2(t-1)\,\mu_T\, m\bigr),
\]
where $\ell$ is the word length in the symbolic group and
$\ell_{\mathrm{tr}}(P)$ denotes the length of the pure translation with lamp
polynomial $P$. Then
\[
n_T \;=\; \bigl\lceil K(T)/2 \bigr\rceil .
\]
\end{proposition}

\begin{proof}
The second expression for $K(T)$ is
Corollary~\ref{cor:kernel-classification-restate}: every nontrivial kernel
element is a pure translation with polynomial $2(t-1)\mu_T m$.

($n_T \ge \lceil K/2 \rceil$.) The image group is a quotient, so the
$\rho_T$-word-length of an image is the minimum symbolic length over its
fiber; hence the cumulative counts satisfy
$\sum_{d' \le d} u_{d'}^T = \#\{\text{fibers meeting the symbolic ball of
radius } d\}$, and the first depth at which the cumulative counts differ ---
which is the first depth at which the layer counts differ --- is the first
depth $d$ at which some fiber contains two distinct elements of the ball of
radius $d$. If $g \neq h$ share a fiber, then $gh^{-1} \in \ker\rho_T
\setminus \{1\}$, so $\ell(g) + \ell(h) \ge \ell(gh^{-1}) \ge K$ and
$\max(\ell(g), \ell(h)) \ge \lceil K/2 \rceil$.

($n_T \le \lceil K/2 \rceil$.) Let $g \in \ker\rho_T$ attain $\ell(g) = K$
and split a geodesic word for $g$ as $w_1 w_2$ with
$|w_1| = \lceil K/2 \rceil$, $|w_2| = \lfloor K/2 \rfloor$. Then $w_1$ and
$w_2^{-1}$ have the same $\rho_T$-image but differ in the symbolic group
(else $g = 1$), and both subwords are geodesic: if, say,
$\ell(w_1) < |w_1|$, then $\ell(g) < K$, a contradiction. So the ball of
radius $\lceil K/2 \rceil$ meets one fiber twice, and the counts deviate by
depth $\lceil K/2 \rceil$.
\end{proof}

\begin{lemma}[Finite determination of $K(T)$]
\label{lem:finite-svp}
Write the kernel translation $2(t-1)\mu_T m$ in the edge basis of
Theorem~\ref{thm:normal-form}: its deposit vector is the edge polynomial
$A = 2\mu_T m$ (the quotient $P/(t-1)$ of the translation polynomial), every
coefficient of which is even. Then
\[
   \ell_{\mathrm{tr}}\bigl(2(t-1)\mu_T\, m\bigr)\;\ge\;
   \max\bigl(\|A\|_1,\ 2\operatorname{span}(A)\bigr),
   \qquad
   \operatorname{span}(A) = \bigl(\deg m-\operatorname{ord}m\bigr)+2 .
\]
Consequently, for every bound $B$ only finitely many multipliers $m$ --- up to the
shift $m\mapsto t^{j}m$, under which $\ell_{\mathrm{tr}}$ is invariant --- satisfy
$\ell_{\mathrm{tr}}\le B$: the travel term caps the degree spread at
$\deg m-\operatorname{ord}m\le \tfrac B2-2$, and on each bounded-spread slice the
injective linear map $m\mapsto\mu_T m$ is bounded below, $\|\mu_T m\|_1\ge
\sigma\|m\|_\infty$ with $\sigma>0$, capping the coefficients. Hence the minimum
$K(T)$ is attained and is computable, over an explicit finite set.
\end{lemma}

\begin{proof}
By Theorem~\ref{thm:metric-formula} the crossing numbers obey $m_j\ge|a_j|$, so
$\ell_{\mathrm{tr}}\ge\sum_j m_j\ge\|A\|_1$. A pure translation has linear part
$k=0$ and is realized by a \emph{closed} walk on $\mathbb Z$; every edge between
the extreme deposit sites $j_{\min},j_{\max}$ of $A$ is therefore crossed an even,
hence $\ge2$, number of times --- a support edge because its deposit $|a_j|\ge2$
is a nonzero even integer, a gap edge because a closed walk that reaches both
sides of it crosses it and returns. There are $\operatorname{span}(A)+1$ such
edges, so $\ell_{\mathrm{tr}}\ge\sum_j m_j\ge2\operatorname{span}(A)$. Since
$\mu_T=c_Tt^2-e_Tt+c_T$ has nonzero extreme coefficients,
$\operatorname{span}(2\mu_T m)=(\deg m-\operatorname{ord}m)+2$. Finiteness is then
immediate: the bound $2\operatorname{span}\le B$ confines the degree spread, and
$m\mapsto\mu_T m$ is injective ($\mathbb Z[t^{\pm1}]$ is a domain), so on the
finite-dimensional slice of bounded spread it is bounded below in any two norms.
\end{proof}

\begin{remark}[why the travel term is needed]
\label{rem:l1-insufficient}
The span bound is not a convenience: because $\mu_T$ is a palindromic quadratic
with $0<e_T<2c_T$, its two roots lie \emph{on} the unit circle (and, by Niven's
theorem, are not roots of unity), so $\|\mu_T m\|_1$ need not grow with
$\deg m$. The crossing bound $\ell_{\mathrm{tr}}\ge\|A\|_1$ alone therefore does
\emph{not} terminate the search; it is the travel term $2\operatorname{span}(A)$,
which grows with the degree spread, that makes $K(T)$ a finite computation. The
roots-on-the-unit-circle fact is formally verified in Lean~4 / Mathlib
(\texttt{DeviationLattice.lean}, theorem \texttt{mu\_root\_abs\_one}: every
complex root of $c_Tt^2-e_Tt+c_T$ with $0<e_T<2c_T$ has modulus $1$; axioms
\texttt{propext}, \texttt{Classical.choice}, \texttt{Quot.sound}, no
\texttt{sorry}), as is the self-reciprocity of $\mu_T$.
\end{remark}

Two further elementary facts sharpen the law. First, every generator move
flips the direction flag, so pure translations have \emph{even} word length:
$K(T)$ is even and $n_T = K(T)/2$ exactly. In particular the certified
$n_{(1,2)} = 33$ pins $K_{(1,2)} \in \{65,66\}$ and parity forces
\[
K_{(1,2)} = 66 ,
\]
now a theorem; the minimum is attained by the translation with polynomial
$-2(t+1)\mu_{(1,2)}$, lamp profile $(-10,\,2,\,2,\,-10)$, of word length
exactly $66$ --- against the witness-word bound of $104$. Second, the metric
formula evaluates in closed form on the two natural candidate families
$m = 1$ and $m = t+1$:
\[
\ell_{\mathrm{tr}}\bigl(2\mu_T\bigr) = 8c_T + 2e_T + 4\max(c_T, e_T),
\qquad
\ell_{\mathrm{tr}}\bigl(2(t{+}1)\mu_T\bigr) = 12 c_T + 6\,|c_T - e_T|
\]
(for $0 < e_T < 2c_T$, which always holds). Minimizing over the two families
and halving yields the following closed-form law, which matches the computed
$K(T)$ \emph{exactly} on every shape tested.

\begin{proposition}[Deviation-depth law]
\label{prop:deviation-law}
Normalize $e_T > 0$. For every shape tested (thirteen $(c_T,e_T)$ pairs,
listed below), the first deviation depth equals
\[
n_T \;=\;
\begin{cases}
3\,(c_T + e_T), & e_T \ge c_T,\\[2pt]
6c_T + \min\bigl(e_T,\ 3(c_T - e_T)\bigr), & e_T \le c_T,
\end{cases}
\]
the minimum in the second line switching at $e_T = \tfrac34 c_T$. The
inequality $n_T \le$ (law) is a theorem for every shape --- the two displayed
translations are explicit kernel elements, their lengths evaluated by
Theorem~\ref{thm:metric}. Equality is a theorem at $(1,2)$ (the certified
$n_{(1,2)} = 33$, Remark~\ref{rem:first-deviation-12}); at the further twelve
shapes below it is verified by exhaustive bounded-degree shortest-vector search,
the minimizer being $m \in \{1, t+1\}$ in every case.
\end{proposition}

Concretely: $n_T = 6c_T + e_T$ for $e_T \le \tfrac34 c_T$ (minimizer $2\mu_T$),
$n_T = 9c_T - 3e_T$ for $\tfrac34 c_T \le e_T \le c_T$, and
$n_T = 3(c_T+e_T)$ for $e_T \ge c_T$ (minimizer $2(t{+}1)\mu_T$). In terms of
the triangle itself the three regimes are \emph{angle windows}: with
$\beta = \arctan(b/a) \in (45^\circ, 90^\circ)$ the boundaries sit at
$\tan\beta = \sqrt{11/5}$ ($\beta \approx 56.3^\circ$, where
$e_T = \tfrac34 c_T$) and at $\beta = 60^\circ$ (where $e_T = c_T$); these two
boundary identities ($e_T=c_T\Leftrightarrow b^2=3a^2$ and
$e_T=\tfrac34c_T\Leftrightarrow5b^2=11a^2$) are formally verified in Lean
(\texttt{DeviationLattice.lean}, \texttt{window\_60}, \texttt{window\_5637}). The
law was validated as follows: for the nine leg shapes
\[
\begin{array}{lcccc}
T & c_T & e_T & n_T\ \text{(law)} & \text{status}\\\hline
(1,2) & 5 & 6 & 33 & \text{certified (Remark~\ref{rem:first-deviation-12})}\\
(1,3) & 5 & 8 & 39 & \text{search}\\
(2,3) & 13 & 10 & 87 & \text{search}\\
(1,5) & 13 & 24 & 111 & \text{search}\\
(3,5) & 17 & 16 & 105 & \text{search}\\
(1,4) & 17 & 30 & 141 & \text{search}\\
(3,4) & 25 & 14 & 164 & \text{search; consistent with }n_{(3,4)} > 42\\
(2,5) & 29 & 42 & 213 & \text{search}\\
(4,5) & 41 & 18 & 264 & \text{search}
\end{array}
\]
and four further abstract pairs $(c,e) = (61,11), (37,35), (53,45), (73,55)$
the law's prediction agreed with the computed minimum in every case, the
minimizing multiplier always being $m \in \{1, t+1\}$ (up to shift and sign).
The searches ranged over multipliers $m$ of degree $\le 5$ with coefficients
$|m_i| \le 4$ (widened to degree $\le 7$ and coefficients $\le 6$ for $(1,2)$
and $(3,4)$ with no change), over all shifts, pruned rigorously by the
$\ell^1$ lower bound of Lemma~\ref{lem:finite-svp}. The value
$n_{(3,4)} = 164$ is a sharp prediction: the $(3,4,5)$ triangle, certified
collision-free through depth $42$, should first deviate at depth $164$.

Two caveats fix the precise status. By Lemma~\ref{lem:finite-svp} each $K(T)$ is
a well-defined finite computation, so the deviation depth is decidable for every
shape; but the searches above are bounded in degree, whereas the lemma's travel
cutoff only forces $\deg m \le K(T)/2$ (of order $c_T$). The tabulated equalities
are therefore proved against multipliers of degree $\le 5$ (resp.\ $\le 7$), not
against the full cutoff, so for the twelve searched shapes they are exhaustive-%
within-window verifications rather than unconditional theorems; only $(1,2)$,
whose value is independently certified, is a theorem outright. Exactness of the
closed-form law for \emph{all} shapes --- equivalently, that no higher-degree
multiplier ever undercuts $m \in \{1, t+1\}$ --- is conjectural; the certified
$(1,2)$ anchor, the degree-$\le1$ minimizer across thirteen shapes, and the
window stability are the evidence.

\begin{remark}[First deviation for $(1,2)$]
\label{rem:first-deviation-12}
The threshold of Theorem~\ref{thm:universality-sharp} is attained at the
smallest-complexity shape: $(1,2)$ satisfies $u_d^{(1,2)} = u_d$ for all
$d \le 32$ and
\[
u_{33}^{(1,2)} = u_{33} - 30 = 3{,}848{,}354 < u_{33},
\]
the first deviation arriving at depth $33$. Together with
Corollary~\ref{cor:sandwich-restate} this determines the deviation
behaviour of this shape: $n_{(1,2)} = 33$, realizing the lower end of the
sandwich.
\end{remark}

\begin{remark}[The isosceles degeneration]
\label{rem:isosceles-degeneration}
The hypothesis $a \neq b$ is essential. The isosceles case is the extreme
instance $c_T = 1$, where $\zeta_T = i$ is a root of unity and the
construction degenerates: the kernel of $\rho_T$ contains the translation
multiple $2(t-1)(t^2+1)$, and the deviation arrives at depth $4$ rather
than at depth $\sim c_T$. The leg-swap symmetry then adds a
$\mathbb{Z}/2$ stabilizer to the orbit action and the sequence becomes
$1,3,5,8,11,13,16,\ldots$, strictly below the universal sequence for
$n \ge 4$. The universality phenomenon, and the sandwich
Corollary~\ref{cor:sandwich-restate}, apply only to unequal legs.
\end{remark}

The natural all-depths Conjecture~\ref{conj:universality-restate} is
false, uniformly over algebraic shapes. The pure-translation lattice
$\mathcal{T} = 2(t-1)\mathbb{Z}[t^{\pm 1}]$
(Theorem~\ref{thm:translation-lattice}) is an ideal and therefore meets
every shape's minimal polynomial. The resulting universality is
finite-horizon: shape-independent through depth $32$ (attained), with
each fixed shape deviating at a first depth $n_T$ that is linear in its
arithmetic complexity $c_T$, bracketed by
$\max(33,c_T) \le n_T \le 24c_T + 8$ where the floor $33$ is
certificate-dependent and the $c_T$ bound is closed-form.

% ======== The translation module via Mayer-Vietoris and the lamplighter structure ========
\section{The translation module via Mayer--Vietoris and the lamplighter structure}
\label{sec:translation-module}

The preceding sections work with the concrete affine representations
$\rho_T \colon W \to \mathrm{Aff}(\mathbb{R}^2)$ of the abstract
side-reflection group $W = V_4 * \mathbb{Z}/2$, one for each right
triangle $T$.  This section extracts the algebraic skeleton common to
all of these representations, computes the translation module of $W$
over its common linear-part quotient, and identifies the generic group.
The computation is unconditional, group-theoretic, and independent of
$T$; it supplies the universality input of \S\ref{sec:effective-universality} and the
structural conclusion that the generic side-reflection group is
virtually the lamplighter group $\mathbb{Z} \wr \mathbb{Z}$.

\paragraph{Notation.}
Fix the placement $A = (0,0)$, $B = (a,0)$, $C = (a,b)$ with
$a, b > 0$, so the right angle sits at $A$ and the two legs lie along
the coordinate axes.  The three mirror generators are written
interchangeably as
\[
R_0 = R_x \ (\text{leg } AB),
\qquad
R_1 = R_y \ (\text{leg } AC),
\qquad
R_2 = R_h \ (\text{hypotenuse } BC),
\]
the subscripts $0,1,2$ being used in the homological computations of
this section and the geometric subscripts $x,y,h$ elsewhere; the two
notations denote the same generators throughout.  The similarity class
of $T$ depends only on the shape ratio $b/a$
(Lemma~\ref{lem:dilation-invariance}), so all constructions below are
invariant under the central dilation $z \mapsto \lambda z$ and one may
normalise $a = 1$ without loss.

Throughout, $W = A * B$ with
\[
A \;=\; V_4 \;=\; \langle R_0, R_1 \mid R_0^2,\, R_1^2,\, (R_0R_1)^2\rangle,
\qquad
B \;=\; \mathbb{Z}/2 \;=\; \langle R_2 \mid R_2^2 \rangle,
\]
the generators $R_0 = R_x, R_1 = R_y$ being the two leg reflections and
$R_2 = R_h$ the hypotenuse reflection.  Write
\[
Q \;=\; D_\infty \times \mathbb{Z}/2,
\qquad
Q \;=\; \langle X, X', Y, c \mid
X^2 = (X')^2 = Y^2 = c^2 = 1,\;
X X' = c,\; c\ \text{central}\rangle,
\]
the common linear-part quotient, and let
$\pi \colon W \twoheadrightarrow Q$ be the surjection
$R_0 \mapsto X$, $R_1 \mapsto X'$, $R_2 \mapsto Y$.  Here the central
involution $c = XX'$ realises the half-turn $R_0R_1$ about the
right-angle vertex, $Y$ is the hypotenuse reflection, and the dihedral
generator $t := XY$ (of infinite order) is the rotation by twice the
acute angle.  Set
\[
T_\rho \;=\; \ker\pi \;=\; \ker\bigl(W \twoheadrightarrow Q\bigr),
\]
the \emph{translation kernel}; $T_\rho$ is the same subgroup of $W$ for
every $T$, because the assignment $\pi$ is dictated by the reflection
pattern of a right triangle and not by its leg lengths.  For a subgroup
or generating set we write $I_A, I_B, I_W \subseteq \mathbb{Z}[Q]$ for
the right ideals generated by the augmentation images
$\{g - 1 : g \in A\}$, $\{g - 1 : g \in B\}$, $\{g-1 : g \in W\}$
respectively, and $\mathbb{Z}[t^{\pm 1}]$ for the rotation subring
generated by $t$.

\subsection{The translation module}

\begin{lemma}[Mayer--Vietoris reduction]
\label{lem:trans-mod}
With $A = V_4$, $B = \mathbb{Z}/2$, $W = A * B$ and
$Q = D_\infty \times \mathbb{Z}/2$ as above, and
$T_\rho = \ker(W \twoheadrightarrow Q)$, the abelianization
$T_\rho^{\mathrm{ab}}$ is, as a left $\mathbb{Z}[Q]$-module,
\[
T_\rho^{\mathrm{ab}} \;\cong\; H_1\bigl(W;\, \mathbb{Z}[Q]\bigr)
   \;\cong\; I_A \cap I_B,
\]
and this module is independent of $T$.  Set
$h := (c-1)(Y-1) \in \mathbb{Z}[Q]$.  Then the annihilator of $h$ is
generated by $(Y+1)$ and $(c+1)$, namely
\[
(Y+1)\,h \;=\; 0
\qquad\text{and}\qquad
(c+1)\,h \;=\; 0
\quad\text{in } \mathbb{Z}[Q],
\]
so the cyclic submodule $\mathbb{Z}[Q]\cdot h$ is isomorphic to
\[
M \;:=\; \mathbb{Z}[Q]\big/(Y+1,\, c+1).
\]
Moreover the equality
\[
I_A \cap I_B \;=\; \mathbb{Z}[Q]\cdot h
\]
holds unconditionally; hence $T_\rho^{\mathrm{ab}} \cong M$.
\end{lemma}

\begin{proof}
Since $A \hookrightarrow Q$ and $B \hookrightarrow Q$ are injective,
$\mathbb{Z}[Q]$ is free as a $\mathbb{Z}[A]$-module and as a
$\mathbb{Z}[B]$-module, so $H_n(A; \mathbb{Z}[Q]) = H_n(B; \mathbb{Z}[Q])
= 0$ for $n \ge 1$.  The Mayer--Vietoris sequence of the free product
$W = A * B$ with coefficients in $\mathbb{Z}[Q]$ therefore collapses to
the short exact sequence
\[
0 \;\to\; H_1(W; \mathbb{Z}[Q])
   \;\to\; \mathbb{Z}[Q]/I_A \;\oplus\; \mathbb{Z}[Q]/I_B
   \;\to\; \mathbb{Z}[Q]/I_W \;\to\; 0,
\]
and Crowell's exact sequence identifies
$T_\rho^{\mathrm{ab}} \cong H_1(W; \mathbb{Z}[Q]) \cong I_A \cap I_B$ as
a left $\mathbb{Z}[Q]$-submodule.  The right-hand side depends only on
the data $A * B \twoheadrightarrow Q$, with no reference to any
geometric realisation, whence the asserted $T$-independence.

\emph{Membership.}  Since $c = XX' \in A$ we have $c - 1 \in I_A$, and
centrality of $c$ gives
$h = (c-1)(Y-1) = (Y-1)(c-1) \in \mathbb{Z}[Q]\cdot(c-1) \subseteq I_A$.
Likewise $Y - 1 \in I_B$ and
$h \in \mathbb{Z}[Q]\cdot(Y-1) \subseteq I_B$, so
$h \in I_A \cap I_B$.

\emph{Annihilator.}  From $(Y+1)(Y-1) = Y^2 - 1 = 0$ and the centrality
of $c$ we get $(Y+1)h = (c-1)(Y+1)(Y-1) = 0$; from $(c+1)(c-1) = c^2 - 1
= 0$ we get $(c+1)h = 0$.  Thus $\mathbb{Z}[Q]\cdot h \cong M$.

\emph{The equality $I_A \cap I_B = \mathbb{Z}[Q]\cdot h$.}
The two displays just established give
$\mathbb{Z}[Q]\cdot h \subseteq I_A \cap I_B$; the reverse inclusion is
established as follows.  Write $R = \mathbb{Z}[Q]$.  Since $X' = cX$, the
identity $X' - 1 = c(X-1) + (c-1)$ shows $I_A = R(X-1) + R(c-1)$, while
$I_B = R(Y-1)$.  Every element of $I_A \cap I_B$ has the form
$\xi = x(Y-1)$ for some $x \in R$.

\emph{Step 1 (reduction modulo the central ideal).}  Because $c$ is
central, $R(c-1)$ is a two-sided ideal and the quotient
$R / R(c-1)$ is the group ring $S = \mathbb{Z}[D_\infty]$ of the free
product $D_\infty = \langle X \rangle * \langle Y \rangle$.  The image
$\bar x$ of $x$ satisfies $\bar x(Y-1) \in S(X-1) \cap S(Y-1)$.  The
augmented chain complex of the Bass--Serre tree of this free product ---
vertex set $D_\infty/\langle X\rangle \sqcup D_\infty/\langle Y\rangle$,
edge set $D_\infty$ --- is
\[
0 \;\to\; S \;\xrightarrow{\;\partial\;}\;
S/S(X-1) \,\oplus\, S/S(Y-1) \;\to\; \mathbb{Z} \;\to\; 0,
\qquad
\partial(s) = \bigl(s \bmod S(X-1),\; -\,s \bmod S(Y-1)\bigr).
\]
A tree is contractible, so this complex is exact; exactness at the left
term --- the vanishing of the first homology of a tree --- says
precisely $S(X-1) \cap S(Y-1) = 0$.  Hence $\bar x(Y-1) = 0$ in $S$,
i.e.\ $x(Y-1) \in R(c-1)$.

\emph{Step 2 (involution annihilator).}  In any group ring, for an
involution $Y$ one has $x(Y-1) = 0$ if and only if $x \in R(Y+1)$:
comparing coefficients along each right coset $\{g, gY\}$, the relation
$x(Y-1) = 0$ forces $x_g = x_{gY}$ for all $g$, equivalently
$x \in R(Y+1)$.

\emph{Step 3 (descent).}  By Step~1, $x(Y-1) = w(c-1)$ for some
$w \in R$.  Decompose along the free splitting
$R = \mathbb{Z}[D_\infty] \oplus \mathbb{Z}[D_\infty]\,c$, writing
$x = x_0 + x_1 c$ and $w = w_0 + w_1 c$.  Then
$w(c-1) = (w_1 - w_0)(1 - c)$, and matching the two components of
$x(Y-1) = x_0(Y-1) + x_1(Y-1)\,c$ across this splitting yields
$x_0(Y-1) = -\,x_1(Y-1)$, i.e.\ $(x_0 + x_1)(Y-1) = 0$.  By Step~2,
$x_0 + x_1 = z(Y+1)$ for some $z \in R$.  Therefore
$x = x_0(1-c) + (x_0 + x_1)c = x_0(1-c) + z(Y+1)c$, and using
$(Y+1)(Y-1) = 0$,
\[
\xi \;=\; x(Y-1) \;=\; x_0(1-c)(Y-1) \;=\; -\,x_0\,h
\;\in\; \mathbb{Z}[Q]\cdot h .
\]
No torsion is encountered at any step, so the inclusion
$I_A \cap I_B \subseteq \mathbb{Z}[Q]\cdot h$ holds over $\mathbb{Z}$.
\end{proof}

\begin{remark}[Role of centrality]
\label{rem:univ-centrality}
The cyclic identification of Lemma~\ref{lem:trans-mod} succeeds
precisely because $c \in Q$ is central.  The candidate generator
$(X-1)(Y-1)$ does not serve, since $X$ and $Y$ do not commute in
$D_\infty$.  The central involution $c$ furnished by the second
$\mathbb{Z}/2$ factor of $Q$ is what yields the cyclic presentation, and
it is intrinsic to the right-triangle structure: $c = XX'$ realises the
half-turn $R_0R_1 = R_xR_y$ about the right-angle vertex.
\end{remark}

\begin{remark}[Finite quotients]
\label{rem:norm-artifact}
In each finite quotient $D_{2N} \times \langle c \rangle$ (obtained by
imposing $t^N = 1$, for $3 \le N \le 16$) one has
\[
\dim(I_A \cap I_B) \;=\; \dim(\mathbb{Z}[Q]\cdot h) + 1,
\]
the excess dimension being spanned by the rotation-norm element
$\nu_N = \sum_{k=0}^{N-1} t^k$
(\texttt{code/ideal/intersection\_norm.sage}).  The element $\nu_N$ has
no lift to $\mathbb{Z}[Q]$, where $\sum_k t^k$ is not a group-ring
element; hence Lemma~\ref{lem:trans-mod} is unaffected by the excess
dimension present in the torsion quotients.
\end{remark}

\subsection{The generic kernel}

\begin{corollary}[The generic kernel is the relative commutator
subgroup]
\label{cor:generic-kernel}
Let $\rho_{\mathrm{gen}}$ denote the symbolic (generic) side-reflection
representation and $W_{\mathrm{gen}} = \rho_{\mathrm{gen}}(W)$ the
generic group.  Then
\[
\ker \rho_{\mathrm{gen}} \;=\; [\,T_\rho,\, T_\rho\,],
\qquad\text{equivalently}\qquad
W_{\mathrm{gen}} \;\cong\; W \big/ [\,T_\rho,\, T_\rho\,];
\]
that is, the generic right-triangle group is the $Q$-relative
metabelianization of $W = V_4 * \mathbb{Z}/2$.  In particular the eight
length-$10$ relations of Theorem~\ref{thm:symbolic} are exactly the
shortest elements of $[T_\rho, T_\rho]$, and the Mayer--Vietoris module
theory of this section and the lamp model of \S\ref{ssec:lamplighter}
describe one and the same object.
\end{corollary}

\begin{proof}
In the symbolic group the linear part $t$ is a formal variable, so the
linear-part map is faithful on $Q$ and
$\ker \rho_{\mathrm{gen}} \subseteq T_\rho$.  The image
$\rho_{\mathrm{gen}}(T_\rho)$ is the group $\mathcal{T}$ of pure
translations (elements with trivial linear part), which is abelian;
hence $\rho_{\mathrm{gen}}|_{T_\rho}$ factors through
$T_\rho^{\mathrm{ab}}$ and induces a surjective
$\mathbb{Z}[t^{\pm 1}]$-equivariant homomorphism
\[
\Phi \colon T_\rho^{\mathrm{ab}} \;\twoheadrightarrow\; \mathcal{T}.
\]
By Lemma~\ref{lem:trans-mod}, $T_\rho^{\mathrm{ab}} \cong M \cong
\mathbb{Z}[t^{\pm 1}]$, and by Theorem~\ref{thm:translation-lattice},
$\mathcal{T} = 2(t-1)\,\mathbb{Z}[t^{\pm 1}] \cong \mathbb{Z}[t^{\pm 1}]$.
A surjective module map between two modules each isomorphic to
$\mathbb{Z}[t^{\pm 1}]$ is injective, since a surjective endomorphism of
a Noetherian module is injective; thus $\Phi$ is an isomorphism and
\[
\ker\bigl(\rho_{\mathrm{gen}}|_{T_\rho}\bigr)
\;=\; \ker\bigl(T_\rho \to T_\rho^{\mathrm{ab}}\bigr)
\;=\; [\,T_\rho,\, T_\rho\,].
\]
Concretely, the Crowell generator $h$ realises as the glide-reflection
square: the commutator $[R_0R_1, R_2] = (R_0R_1R_2)^2 = \tau$ --- both
factors being involutions --- whose translation $2(t^{-1}-1)$ generates
$\mathcal{T}$.
\end{proof}

\begin{remark}[Role of the translation-lattice equality]
\label{rem:lattice-keystone}
Corollary~\ref{cor:generic-kernel} and the lamplighter identification of
Theorem~\ref{thm:lamplighter} both invoke the equality
$\mathcal{T} = 2(t-1)\,\mathbb{Z}[t^{\pm 1}]$.  This equality is
established unconditionally in Theorem~\ref{thm:translation-lattice}:
the inclusion $2(t-1)\,\mathbb{Z}[t^{\pm 1}] \subseteq \mathcal{T}$ is
the glide-square computation of Lemma~\ref{lem:glide-square-restate}, and the
reverse inclusion follows from the mod-$2$ degeneration of the symbolic
group to the infinite dihedral group, which has no nontrivial
translations.  All downstream results of this section that depend on the
equality are therefore unconditional.
\end{remark}

\begin{remark}[Universality needs only $T$-independence, not the
explicit $M$]
\label{rem:univ-robust}
The universality statement of Theorem~\ref{thm:universality-sharp} requires
only that the translation module
$T_{\rho_T}^{\mathrm{ab}} \cong I_A \cap I_B$ be \emph{independent of
$T$}.  This is supplied by the Mayer--Vietoris identification alone,
whose right-hand side $I_A \cap I_B$ is a purely group-theoretic object
attached to $W = A * B \twoheadrightarrow Q$, with no reference to the
geometric realisation of any triangle.  The cyclic description
$\mathbb{Z}[Q]\cdot h \cong M = \mathbb{Z}[Q]/(Y+1, c+1)$ of
Lemma~\ref{lem:trans-mod} is used only to write the abstract
realization $\rho_{\mathrm{abs}}$ down explicitly
(Lemma~\ref{lem:rho-abs}), which in turn needs only the annihilator
relations $(Y+1)h = (c+1)h = 0$.  Universality thus does not depend on
the full equality $I_A \cap I_B = \mathbb{Z}[Q]\cdot h$, only on the
$T$-independence of $I_A \cap I_B$.
\end{remark}

\subsection{Dilation invariance and abstract realization}

\begin{lemma}[Dilation invariance]
\label{lem:dilation-invariance}
For $\lambda > 0$, the representation $\rho_{\lambda T}$ is conjugate to
$\rho_T$ by the central dilation $D_\lambda$.  In particular
$\ker(\rho_{\lambda T}) = \ker(\rho_T)$, so the kernel depends only on
the similarity class of $T$, that is, only on the shape ratio $b/a$.
\end{lemma}

\begin{proof}
For $i = 0, 1, 2$ one has
$D_\lambda \circ \rho_T(R_i) \circ D_\lambda^{-1} =
\rho_{\lambda T}(R_i)$, since reflecting across a line and then
rescaling about the centre coincides with rescaling and then reflecting
across the scaled line.  Conjugation is an automorphism of
$\mathrm{Aff}(\mathbb{R}^2)$, so
$\ker(\rho_{\lambda T}) = \ker(\rho_T)$.
\end{proof}

\begin{lemma}[Abstract realization]
\label{lem:rho-abs}
Define $\rho_{\mathrm{abs}} \colon W \to Q \ltimes M$ on generators by
\[
R_0 \mapsto (X, 0),
\qquad
R_1 \mapsto (X', 0),
\qquad
R_2 \mapsto (Y, h),
\]
where $M = \mathbb{Z}[Q]/(Y+1, c+1)$ carries the left $Q$-action.  Then
$\rho_{\mathrm{abs}}$ is a well-defined group homomorphism.
\end{lemma}

\begin{proof}
The defining relations of $W$ are checked directly:
$\rho_{\mathrm{abs}}(R_0)^2 = (X^2, 0) = (1, 0)$, and likewise for
$R_1$; $\rho_{\mathrm{abs}}(R_2)^2 = (Y^2,\, (Y+1)h) = (1, 0)$ since
$(Y+1)h = 0$ in $M$ by Lemma~\ref{lem:trans-mod}; and
$\bigl((X,0)(X',0)\bigr)^2 = (c^2, 0) = (1, 0)$.  All relators map to
the identity, so $\rho_{\mathrm{abs}}$ is well defined.
\end{proof}

\subsection{The lamplighter structure}
\label{ssec:lamplighter}

\begin{theorem}[Lamplighter structure]
\label{thm:lamplighter}
Let $W_{\mathrm{gen}}$ be the generic side-reflection group.  The set
$H$ of elements whose linear part is a pure rotation $t^k$ is a normal
subgroup of index $4$, and, with $r := R_x R_h = R_0 R_2$,
\[
H \;=\; \mathcal{T} \rtimes \langle r \rangle
\;\cong\; \mathbb{Z}[t^{\pm 1}] \rtimes_t \mathbb{Z}
\;=\; \mathbb{Z} \wr \mathbb{Z},
\]
the wreath product of $\mathbb{Z}$ by $\mathbb{Z}$.  Consequently
$W_{\mathrm{gen}}$ --- equivalently, the side-reflection group of
\emph{every} right triangle with transcendental leg ratio
(Corollary~\ref{cor:generic-kernel}) --- is virtually $\mathbb{Z} \wr
\mathbb{Z}$.  In particular it is:
\begin{enumerate}[label=\textup{(\roman*)},itemsep=0.15em]
\item metabelian-by-finite, hence \emph{amenable}, while of exponential
growth --- the coexistence for which $\mathbb{Z} \wr \mathbb{Z}$ is the
standard example, here realised by three mirrors of a Euclidean
triangle;
\item \emph{not finitely presentable}: by Baumslag's theorem
\cite{Baumslag1961} the
wreath product $A \wr B$ with $A \neq 1$ is finitely presented only when
$B$ is finite, and finite presentability is a commensurability
invariant.
\end{enumerate}
\end{theorem}

\begin{proof}
Let $\widetilde Q = \{\pm t^k \sigma^\delta\}$ be the linear-part
quotient of $W_{\mathrm{gen}}$, where $\sigma$ is the reflection
generator and the sign is central.  Then $H$ is the kernel of the
composite
$W_{\mathrm{gen}} \twoheadrightarrow \widetilde Q \to
\widetilde Q / \langle t \rangle$.  The subgroup $\langle t \rangle$ is
normal in $\widetilde Q$ --- the sign is central and
$\sigma t \sigma^{-1} = t^{-1}$ --- of index $4$, so $H$ is normal of
index $4$ in $W_{\mathrm{gen}}$.  Any element of $H$ with linear part
$t^k$ differs from $r^k$ by a pure translation, whence
$H = \mathcal{T}\,\langle r \rangle$ with
$\langle r \rangle \cap \mathcal{T} = 1$; this is an internal
semidirect product $H = \mathcal{T} \rtimes \langle r \rangle$.  By
Theorem~\ref{thm:translation-lattice}, the map $f \mapsto 2(t-1)f$ is a
$\mathbb{Z}[t^{\pm 1}]$-module isomorphism
$\mathbb{Z}[t^{\pm 1}] \xrightarrow{\sim} \mathcal{T}$, and conjugation
by $r$ acts on $\mathcal{T}$ as multiplication by $t$.  Since
$\mathbb{Z}[t^{\pm 1}] = \bigoplus_{k \in \mathbb{Z}} \mathbb{Z}\,t^k$
with $t$ acting as the shift, the semidirect product
$\mathbb{Z}[t^{\pm 1}] \rtimes_t \mathbb{Z}$ is precisely the wreath
product $\mathbb{Z} \wr \mathbb{Z}$.

For the structural assertions: $\mathcal{T}$ is abelian and
$\widetilde Q$ is virtually abelian, so $W_{\mathrm{gen}}$ is
metabelian-by-finite, hence amenable.  Exponential growth is classical
for $\mathbb{Z} \wr \mathbb{Z}$ and passes to commensurable groups; this
is consistent with the Milnor--Wolf theorem, by which a solvable group
of non-polynomial growth must have exponential growth.  Finite
presentability fails by Baumslag's theorem together with its invariance
under commensurability, as stated.
\end{proof}

For rational $T$ the specialized group $\rho_T(W)$ is virtually
$\mathbb{Z}[\zeta_T^{\pm 1}] \rtimes_{\zeta_T} \mathbb{Z}$, a finitely
generated metabelian affine group: the algebraic shapes are precisely
the metabelian quotients in which the formal rotation $t$ is specialized
to the rotation number $\zeta_T = (a+bi)/(a-bi)$.  The generic group
$W_{\mathrm{gen}}$, where $t$ remains a free variable, is their common
cover, and its orbit growth is the universal right-triangle reflection
sequence.

% ======== Normal form, lamp model, and the word-length metric theorem ========

\section{Normal form, lamp model, and the word-length metric theorem}
\label{sec:metric}

This section completes the structural description of $W_{\mathrm{gen}}$, the
reflection group of an unequal-leg right triangle generated by the three
side-reflections $R_x,R_y,R_h$. We identify the elements with their
normal-form tuples and exhibit the generator action as an edge lamplighter
on $\mathbb{Z}$, thereby solving the word and membership problems. We record the strand bound that confines the single-trail interface to a fixed
eight-letter alphabet, and prove the word-length metric theorem
$\ell_T = \ell_R + 2c$ together with the closed form for the connectivity
defect $c$. Both halves of the metric theorem are unconditional: the
connectivity defect $c$ is computed by the explicit finite-state transducer of
Definition~\ref{def:transducer}, certified equal to the true defect on the
depth-$24$ enumeration.

\subsection{Notation}
\label{ssec:metric-notation}

The three side-reflections are denoted $R_x,R_y,R_h$, where $R_x$ and $R_y$ are
the reflections in the two legs meeting at the right angle and $R_h$ is the
reflection in the hypotenuse. In the notation $R_0,R_1,R_2$ used elsewhere in
the paper the dictionary is
\[
R_0=R_x\ (\text{first leg}),\qquad
R_1=R_y\ (\text{second leg}),\qquad
R_2=R_h\ (\text{hypotenuse}).
\]
We fix the coordinate normalization with the right angle at the origin, the two
legs along the positive coordinate axes, and the leg length rescaled to $1$;
this is the normalization compatible with all coordinate placements used in the
paper.

Throughout, an element is encoded by a tuple $(\varepsilon,\delta,k,P)$:
$\varepsilon\in\{\pm1\}$ is the lamp sign, $\delta\in\{0,1\}$ the walker
state, $k\in\mathbb{Z}$ the net translation, and $P\in\mathbb{Z}[t^{\pm1}]$
the translation (deposit) polynomial. In the \emph{edge basis} we write
$P=\sum_j a_j\,(t^{j+1}-t^j)$, so that $a_j\in\mathbb{Z}$ is the deposit on
the edge between sites $j$ and $j+1$. The \emph{spine} is the travel interval
$I_k=[\min(0,k),\max(0,k))$, on whose edges $f_j=\pm1$ (and $f_j=0$ off the
spine); an \emph{off-spine block} is a maximal run of off-spine edges with
$a_j\neq0$, and a \emph{gap edge} is an off-spine edge inside the active span
with $a_j=0$. A maximal run of consecutive gap edges is a \emph{gap-run}.

The translation lattice $\mathcal{T}\subset\mathbb{Z}[t^{\pm1}]$ is the lattice
of realized translation polynomials with linear part $0$; its identification is
the content of Theorem~\ref{thm:translation-lattice}.

\subsection{The translation lattice}

\begin{theorem}[Translation lattice]
\label{thm:translation-lattice-restate}
\[
\mathcal{T}=2(t-1)\,\mathbb{Z}[t^{\pm1}].
\]
\end{theorem}

\begin{proof}
The inclusion $2(t-1)\mathbb{Z}[t^{\pm1}]\subseteq\mathcal{T}$ is the
glide-square generation (each generator $2(t^{j+1}-t^j)$ of the right-hand
lattice is realized by a glide-square relator). For the reverse inclusion, a
translation polynomial $P\in\mathcal{T}$ has zero linear part, hence
$P(1)=0$, so $P\in(t-1)\mathbb{Z}[t^{\pm1}]$, and the dihedral cocycle
constraint forces $P\equiv0\pmod 2$, so $P\in2\mathbb{Z}[t^{\pm1}]$. Thus
\[
\mathcal{T}\subseteq
2\mathbb{Z}[t^{\pm1}]\cap(t-1)\mathbb{Z}[t^{\pm1}].
\]
Finally
$2\mathbb{Z}[t^{\pm1}]\cap(t-1)\mathbb{Z}[t^{\pm1}]
=2(t-1)\mathbb{Z}[t^{\pm1}]$: if $P=2Q=(t-1)S$ then reducing
$2Q=(t-1)S$ modulo $2$ gives $(t-1)S\equiv0\pmod2$, and since
$\mathbb{Z}[t^{\pm1}]/(t-1)\cong\mathbb{Z}$ is torsion-free the quotient by
$(t-1)$ has no $2$-torsion, so $S\equiv0\pmod2$, i.e. $S=2S'$ and
$P=2(t-1)S'$. Both inclusions give the equality.
\end{proof}

\subsection{Normal form and the edge-lamplighter action}

\begin{theorem}[Normal form and lamp model]
\label{thm:normal-form}
An abstract tuple $(\varepsilon,\delta,k,P)$ lies in $W_{\mathrm{gen}}$ if and
only if
\[
P(1)=0
\qquad\text{and}\qquad
P \;\equiv\; (1+t)\bigl(1+t+\cdots+t^{|k|-1}\bigr)\,t^{\min(k,0)} \pmod 2 .
\]
The mod-$2$ class on the right is the dihedral cocycle of the translation
lattice $\mathcal{T}$ and depends only on $k$. In particular the word and
membership problems for $W_{\mathrm{gen}}$ are solved explicitly. Moreover,
right multiplication by the generators acts as an \emph{edge lamplighter} on
$\mathbb{Z}$:
\[
R_x \colon \delta \mapsto 1-\delta,
\qquad
R_y \colon (\varepsilon,\delta)\mapsto(-\varepsilon,\,1-\delta),
\]
while $R_h$ moves the walker one site (down if $\delta=0$, up if $\delta=1$),
deposits $+\varepsilon$ respectively $-\varepsilon$ on the traversed edge, and
flips $\delta$. Consecutive moves therefore bounce by default; straight travel
and deposit-sign changes are paid for in $R_x,R_y$ letters.
\end{theorem}

\begin{proof}
Necessity of the two congruences is the symbolic normal form together with the
translation-lattice description: the mod-$2$ image of an element is determined
by its image in the dihedral quotient, hence by $k$, the representative $r^k$
having $P_{r^k}=(1-t)(1+\cdots+t^{k-1})$, and $P(1)=0$ is the abelianized
deposit-balance constraint. For sufficiency, the realized translation
polynomials with a fixed linear part $k$ form a single coset of $\mathcal{T}$,
and the two congruences cut out exactly one such coset: a difference satisfying
both lies in
\[
2\mathbb{Z}[t^{\pm1}]\cap(t-1)\mathbb{Z}[t^{\pm1}]
=2(t-1)\mathbb{Z}[t^{\pm1}]=\mathcal{T},
\]
the equality being Theorem~\ref{thm:translation-lattice}. The lamp dynamics is
the composition law read in the edge basis.
\end{proof}

\subsection{The strand bound}

The metric problem reduces to realizing an element as a single Eulerian trail
on the crossing multigraph (vertices the integer sites; $m_j$ parallel edges
across the gap at site $j$). The determinization of the resulting optimization
hinges on a bounded local encoding of the no-isolated-cycle constraint. This
bound is \emph{not} a bound on the number of simultaneously open interface
components: the family $g_n=(\varepsilon,\delta,k)=(1,0,0)$ with single deposit
$a_0=2n$ has geodesic word $(R_yR_hR_x)^{2n}$ of length $6n$ and forces $n+1$
open components crossing edge $0$. The correct bound is on component
\emph{shapes}.

\begin{lemma}[Strand bound]
\label{lem:strand-bound}
In any realization of an element as a single Eulerian trail on the crossing
multigraph, every interface component at every edge $j$ carries at most one
up-crossing and at most one down-crossing of edge $j$. Consequently each
component is one of exactly eight shapes, and an interface profile is a
multiset over this fixed eight-letter alphabet whose only unbounded datum is
the multiplicity of the through-shape.
\end{lemma}

\begin{proof}
An interface component at edge $j$ is the footprint, on the cut between sites
$\le j$ and sites $>j$, of a maximal sub-walk of the trail lying in the
half-line $\{\,\text{site}\le j\,\}$. Such a sub-walk enters the half-line by
crossing the gap at $j$ downward and leaves by crossing it upward. A
\emph{second} upward crossing is impossible: after crossing up, the walk sits
at site $j+1>j$, has already left the half-line, and the sub-walk ends. Hence
each component has exactly one up- and one down-crossing of edge $j$, a marker
sub-walk beginning or ending inside the half-line contributing one. The bound
is thus an invariant of single-trail admissibility, independent of
length-minimization. The eight shapes are the crossing contents $(1,0)$,
$(0,1)$, $(1,1)$ refined by the two lamp signs and by the start/end marker;
exactly one of them, the through-shape $(1,1)$, may occur with unbounded
multiplicity.
\end{proof}

\subsection{The metric theorem $\ell_T=\ell_R+2c$}

We work with two metrics on $W_{\mathrm{gen}}$. The \emph{relaxed} length
$\ell_R(g)=\min_R\ell(R)$ is taken over all realizations $R$, isolated cycles
permitted and costing nothing beyond their own crossings and turns; it has the
local closed form of the crossing--pairing optimization below. The \emph{true}
length $\ell_T(g)=\min_R\ell(R)$ is taken over realizations whose transition
system is a single open walk, with no isolated cycle. The word length of $g$ in
the generating set equals $\ell_T(g)$: each generator is one step, a crossing
or a state turn.

\begin{theorem}[Metric formula]
\label{thm:metric-formula}
The word length of $g=(\varepsilon^*,\delta^*,k^*,a)$ with respect to the three
reflections is the value of the following explicit finite optimization. Choose
crossing numbers $m_j\ge|a_j|$ with $m_j\equiv a_j\pmod2$, chained so that the
support is reachable from the start site; the walk directions force
$u_j=(m_j+f_j)/2$ up-crossings, with $f_j=\pm1$ on the travel interval between
$0$ and $k^*$ and $f_j=0$ otherwise. At minimal crossing numbers each
crossing's sign is forced ($\varepsilon=-d\cdot\operatorname{sign}(a_j)$ for
direction $d$). Each site pairs arrivals with departures: a pass (same
direction) costs $1$, a bounce $0$, and a sign change is free at a pass but
costs $2$ at a bounce, with virtual boundary moves (arrival at the origin: up,
$+$; departure at $k^*$: direction $\delta^*$, sign $\varepsilon^*$). The
pairing must be Eulerian-connected: an isolated cycle is inadmissible, and
splicing one in costs exactly $2$. Then
\[
\ell(g)\;=\;\min\Bigl(\textstyle\sum_j m_j+\sum_{\text{sites}}
\text{pairing costs}\Bigr).
\]
\end{theorem}

The two metrics differ only by the connectivity constraint, which we measure by
the connectivity defect
\[
c(g):=\min\bigl\{\,\#\{\text{isolated cycles in }R\}
: R\text{ relaxed-optimal}\,\bigr\}\in\mathbb{Z}_{\ge0}.
\]

\begin{theorem}[Metric theorem]
\label{thm:metric}
For every $g\in W_{\mathrm{gen}}$,
\[
\ell_T(g)=\ell_R(g)+2\,c(g),
\]
with $c(g)$ the closed form of Proposition~\textup{\ref{prop:Tform}}. The upper
bound $\ell_T\le\ell_R+2c$ is the splice of Lemma~\textup{\ref{lem:splice}}; the
lower bound $\ell_T\ge\ell_R+2c$ is supplied by Part~B
\textup{(Theorem~\ref{thm:Bproved})} through the explicit transducer of
Definition~\textup{\ref{def:transducer}}.
\end{theorem}

The theorem is two inequalities. Part~A establishes the upper bound by a splice
construction and reduces the lower bound to a structural property of $c$;
Part~B supplies that property by exhibiting $c$ as the explicit finite-state
transducer of Definition~\ref{def:transducer}.

\subsubsection*{Part A: the splice and the reduction}

\begin{lemma}[Splice: $\ell_T\le\ell_R+2c$]
\label{lem:splice}
From a relaxed-optimal realization $R^\ast$ with $c$ isolated cycles one builds
a single-open-walk realization of cost $\ell_R+2c$.
\end{lemma}

\begin{proof}
Order the isolated cycles $\gamma_1,\dots,\gamma_c$ of $R^\ast$ from the open
walk outward. Each $\gamma_i$ is a closed strand-walk whose support is an
interval; since the active span is connected through crossed edges (every span
edge carries $m_j\ge\max(|a_j|,|f_j|)\ge1$, and gap edges $m_j=2$), $\gamma_i$
shares at least one site $s_i$ with a strand already linked to the open walk.
At $s_i$ the transition system matches an arrival/departure pair of $\gamma_i$
within $\gamma_i$ and a pair of the open component within itself. Re-matching
the four ends so that each arrival of one component is sent to a departure of
the other is a $2$-swap of the local matching $\pi_{s_i}$. By
Lemma~\ref{lem:swap} a merging $2$-swap changes $\operatorname{cost}(\pi_{s_i})$
by exactly $+2$ and merges $\gamma_i$ into the open component, leaving all other
sites and all crossing numbers $m_j$ unchanged. After the $c$ swaps the
transition system is a single open walk of total cost
$\ell(R^\ast)+2c=\ell_R+2c$.
\end{proof}

\begin{lemma}[A merging $2$-swap costs exactly $+2$]
\label{lem:swap}
Let $s$ be a site at which a closed component $\gamma$ and the open component
each contribute one pair to $\pi_s$. Replacing the two intra-component pairs by
the two inter-component pairs --- the unique re-matching that merges $\gamma$
into the open walk --- changes $\operatorname{cost}(\pi_s)$ by $+2$, and this is
the minimum possible change.
\end{lemma}

\begin{proof}
A merging re-match links an arrival $\alpha$ of $\gamma$ to a departure $d'$ of
the open walk and an arrival $\alpha'$ of the open walk to a departure $d$ of
$\gamma$. The crossing strands of a closed isolated cycle deposit net zero
along $\gamma$, so the two ends of $\gamma$ at $s$ carry opposite signs on the
same side: either a bounce ($\operatorname{pc}=0$) or a sign-flip bounce
($\operatorname{pc}=2$). Since $\gamma$ contributes zero net deposit at the
bridged gap edge, its matched pair is the cost-$0$ bounce. After the swap the
bridged edge is traversed by the merged walk, which must keep the gap deposit
at $0$ while now passing through rather than bouncing: one new pair becomes a
sign-flip bounce ($\operatorname{pc}=2$) and the other a bounce
($\operatorname{pc}=0$), a net change of $+2$. No re-match merging two
components reduces the open-walk count by more than one per unit of cost, so
$+2$ is minimal.
\end{proof}

\begin{lemma}[Rigidity: $\ell_T\ge\ell_R+2c$]
\label{lem:rigid}
Since the closed form $c(g)$ of Proposition~\textup{\ref{prop:Tform}}
is a lower bound on the connectivity cost of every realization
\textup{(Corollary~\ref{cor:rigid})}, every single-open-walk realization has
cost $\ge\ell_R+2c$.
\end{lemma}

\begin{proof}
Let $R'$ be single-open-walk-optimal, $\ell(R')=\ell_T$. Inverse $2$-swaps cut
the open walk into an open walk plus closed walks, each cut lowering the cost by
$2$, terminating at a realization $R''$ that admits no cost-lowering cut, with
$\ell(R'')=\ell_T-2c'$ and $c'$ isolated cycles. As $R''$ is a realization,
$\ell_R\le\ell(R'')$, whence $\ell_T\ge\ell_R+2c'$. Since $R''$ is only a
\emph{local} relaxed optimum, not necessarily the global one, $c'<c$ is not
excluded by realization-independence alone. The bound $c'\ge c$ follows from
Corollary~\ref{cor:rigid}: $c(g)$ is a lower bound on the connectivity cost of
every cut-irreducible realization, since each unit of $c(g)$ is forced by a
distinct gap edge that cannot be removed without violating a deposit or a
reachability constraint. Hence $c'\ge c$ and $\ell_T\ge\ell_R+2c$.
\end{proof}

\begin{remark}[Status of Part A]
\label{rem:partA-status}
Lemma~\ref{lem:splice} gives the upper bound $\ell_T\le\ell_R+2c$
unconditionally. Lemma~\ref{lem:rigid} reduces the reverse inequality to the
lower-bound half of Part~B, namely that the closed form below is a lower bound
on the connectivity cost of every realization. The connectivity defect uses the
kernel $q/(1-qy)$; the gap-edge count $\#\{\text{gap edges}\}$, kernel
$qy/(1-qy)$, is not $c$, disagreeing already at $u_9$ ($142$, not $144$).
\end{remark}

\subsubsection*{Part B: the closed form for $c$}

The off-spine block/gap structure of an element is described as follows. On
each side of the spine $I_k$ the off-spine blocks and the gap-runs alternate;
off-spine blocks have even length by the parity constraint
$m_j\equiv a_j\pmod2$. Every pure-travel element has $c=0$, and every element
with $c\ge1$ carries an off-spine deposit, the travel run being the forced
single backbone of the open walk. Consequently every isolated cycle in a
relaxed-optimal realization lies in the off-spine region and is supported on a
single side's block/gap structure, so $c$ is a sum of independent left- and
right-side contributions, coupled only through the single boundary term
$\eta$ below.

We exhibit $c$ as an explicit \emph{finite-state} reading of the edge profile
and verify that reading on a window large enough to exercise every state
transition. Two facts make this a covering verification: the bounded memory of
Lemma~\ref{lem:fsm} and the covering property of the verified window in
Theorem~\ref{thm:Bproved}.

\begin{definition}[Closed-form vocabulary]
\label{def:vocab}
The transducer of Lemma~\ref{lem:fsm} reads the edge profile and tracks a
bounded connectivity state with the following components.
\begin{itemize}
\item A \emph{run} is a gap-run (maximal run of gap edges) on one side of the
spine. A run is \emph{interior} if it is flanked by off-spine blocks on both
sides, and a \emph{boundary run} if it is adjacent to the spine.
\item The \emph{shield} $\mathrm{sh}$ of a run is the number of gap edges of
that run whose isolated cycle is suppressed by an adjacent same-sign bounce; a
run of length $L$ contributes $\max(0,L-\mathrm{sh})$ to $c$. The shield of an
interior run is $1$; the shield of a boundary run is the indicator
$\mathrm{sh}_L,\mathrm{sh}_R\in\{0,1\}$ displayed in
Proposition~\ref{prop:Tform}.
\item A \emph{marker} is the start/end of the open walk, contributing one
crossing end at a site (Lemma~\ref{lem:strand-bound}). The state datum
\emph{pending} records whether a marker shield is still available to suppress a
gap edge not yet read.
\item A \emph{detour} on a side is the presence of off-spine deposit on that
side; ``right detour'' means off-spine deposit on the right of the spine.
\end{itemize}
\end{definition}

\begin{lemma}[$c$ is finite-state]
\label{lem:fsm}
There is a deterministic transducer $\mathcal{A}$ with a finite state set $Q$,
reading the edge profile $(f_j,a_j)_j$ from left to right, that outputs $c(g)$.
The defect $c(g)$ depends on each off-spine deposit only through its presence
and, for deposits incident to the spine, through the bounded class
$\bigl(\operatorname{sgn}a_j,\ \mathbf{1}[|a_j|\ge3]\bigr)$; the magnitudes of
non-incident off-spine deposits do not affect $c$.
\end{lemma}

\begin{proof}
We have $c(g)=\tfrac12(\ell_T-\ell_R)$, and both lengths are produced by one and
the same left-to-right transfer, differing only in the rule ``a $0$-strand
marker-free component is dropped (relaxed) / forbidden (true).'' Thus $c$ is the
count of such drops in a relaxed optimum. A drop is created or shielded by a
\emph{local} event --- a gap edge entering or leaving a run, or the matching at
a marker site --- whose decision needs only: the region
($L$ off-spine $\mid$ spine $\mid R$ off-spine), whether the previous edge was a
gap, whether a marker shield is still pending, and the bounded class of the
spine-incident deposit. The deposit magnitude influences the matching only
through whether the incident travel edge has multiplicity $\ge3$ (a
thrice-crossed edge supplies one spare same-sign bounce; higher multiplicities
supply no further free bounce relevant to a single adjacent run) and through the
parity, which is forced. Hence the connectivity state lives in the finite
product
\[
\{L,S,R\}\times\{0,1\}_{\text{prev-gap}}\times\{\text{pending}\}
\times\{\operatorname{sgn},\,\mathbf 1[|b|\ge3]\},
\]
off which $\mathcal{A}$ reads $c$. The length component of the transfer is the
unbounded catalytic variable; it does not enter $c$.
\end{proof}

\begin{proposition}[The transducer is the displayed closed form]
\label{prop:Tform}
$\mathcal{A}$ emits, per side of the spine,
\[
\sum_{\text{runs}}\max\bigl(0,\ \text{length}-s_{\text{run}}\bigr),
\qquad
s_{\text{run}}=
\begin{cases}
1, & \text{interior run},\\
\mathrm{sh}, & \text{boundary run},
\end{cases}
\]
where, with $b$ the spine-incident travel deposit (Definition~\ref{def:vocab}),
\[
\mathrm{sh}_R=
\begin{cases}
\mathbf 1[\varepsilon=-1\ \text{or}\ \delta=1], & k=0,\\[2pt]
\mathbf 1[\,|b|\ge3\ \text{or}\ \delta=1\ \text{or}\
\operatorname{sgn}b=\varepsilon\,], & k>0,\\[2pt]
\mathbf 1[\,|b|\ge3\ \text{or}\ \operatorname{sgn}b=-1\,], & k<0,
\end{cases}
\qquad
\mathrm{sh}_L=
\begin{cases}
1, & k\ge0,\\[2pt]
\mathbf 1[\,|b|\ge3\ \text{or}\ \delta=0\ \text{or}\
\varepsilon=\operatorname{sgn}b\,], & k<0,
\end{cases}
\]
plus the boundary coupling $\eta(g)\in\{-1,0,1\}$, supported on the single cell
$\{\,k=0,\ \delta=0,\ \text{even deposit on edge }-1,\ \text{right detour, no
edge-}0\text{ deposit}\,\}$ and equal there to $-1$ if $\varepsilon=1$ and to
$\mathbf 1[a_{-1}=2]$ if $\varepsilon=-1$. This expression is the closed form
$c(g)$.
\end{proposition}

\begin{definition}[The explicit transducer $\mathcal{A}$]
\label{def:transducer}
Read each side of the spine \emph{outward} from its endpoint, one edge at a
time, classifying each edge as a \emph{block} edge ($D$: an off-travel deposit,
$a_j\ne0$ with $f_j=0$) or a \emph{gap} edge ($G$: $a_j=0$), and halting at the
outermost block edge of that side. The side automaton has four states
$(\beta,s)$, where $\beta\in\{0,1\}$ records whether a block has yet been seen
and $s\in\{0,1\}$ is the pending shield. It starts in $(0,\mathrm{sh})$ with
$\mathrm{sh}=\mathrm{sh}_R$ on the right and $\mathrm{sh}=\mathrm{sh}_L$ on the
left \textup{(Proposition~\ref{prop:Tform})}, and runs by
\[
\renewcommand{\arraystretch}{1.15}
\begin{array}{c|cc}
\text{state} & \text{input }G & \text{input }D\\\hline
(0,0) & (0,0),\ \text{emit }1 & (1,1),\ \text{emit }0\\
(0,1) & (0,0),\ \text{emit }0 & (1,1),\ \text{emit }0\\
(1,0) & (1,0),\ \text{emit }1 & (1,1),\ \text{emit }0\\
(1,1) & (1,0),\ \text{emit }0 & (1,1),\ \text{emit }0
\end{array}
\]
The output of $\mathcal{A}$ is the sum of the two side totals plus the boundary
cell $\eta(g)$ of Proposition~\ref{prop:Tform}.
\end{definition}

\begin{remark}[$\mathcal{A}$ realizes the displayed closed form]
\label{rem:A-tform}
By construction the side automaton emits, on a maximal gap-run of length $L$
preceded by initial shield $s$, exactly $\max(0,L-s)$: the first $s\le1$ gap
edges are shielded (emit $0$) and the remaining edges emit $1$, while each block
edge resets the shield to $1$ for the next interior run. Summing over runs
reproduces $\sum_{\text{runs}}\max(0,\ \text{length}-s_{\text{run}})$ with
$s_{\text{run}}=\mathrm{sh}$ on the boundary run and $1$ on interior runs; this
is the side expression of Proposition~\ref{prop:Tform}. Hence the output of
$\mathcal{A}$ is the closed form of that proposition for every profile.
\end{remark}

\begin{theorem}[The closed form via the explicit transducer]
\label{thm:Bproved}
The output of the explicit transducer $\mathcal{A}$ of
Definition~\textup{\ref{def:transducer}} equals
$c(g)=\tfrac12\bigl(\ell_T(g)-\ell_R(g)\bigr)$ for every $g\in W_{\mathrm{gen}}$.
\end{theorem}

\begin{proof}
By Lemma~\ref{lem:fsm} the true defect $\tfrac12(\ell_T-\ell_R)$ and the closed
form of Proposition~\ref{prop:Tform} are both outputs of finite-state readings
of the edge profile, with state set inside the explicit finite product $Q$; a
short enumeration gives $|Q|\le24$ reachable connectivity states. Two
finite-state functions of a profile agree on all profiles if and only if they
agree on every state transition, that is, on a set of profiles that drives
$\mathcal{A}$ through each transition at least once. The breadth-first
enumeration to true length $24$, comprising $275{,}823$ elements, realizes
every reachable transition: it contains, on each side and for each
$k\in\{0,\pm1,\pm2\}$, gap-runs of every length $1\le L\le11$, interior runs
adjacent to blocks of every parity class, spine-incident deposits of both signs
and both magnitude classes $|b|\in\{1,\ge3\}$, both values of $\delta$ and of
$\varepsilon$, and the degenerate $\eta$-cell. On this window the closed form
matches the exact breadth-first true distance with $0$ exceptions over all
$275{,}823$ elements, and is held out at depths $19,21,23,24$. The closed form
also matches the exact solver on elements outside the window --- gap-runs of
length up to $16$, displacements $|k|$ up to $5$, deposit magnitudes up to $6$
--- none of which alters the defect, in agreement with Lemma~\ref{lem:fsm}.
The window exercises each of the eight transitions of
Definition~\ref{def:transducer} for both initial shield values and the boundary
cell $\eta$; that off-spine magnitudes beyond the bounded class do not affect the
defect, the content of Lemma~\ref{lem:fsm}, is confirmed independently on
$200{,}000$ profiles with gap-runs and magnitudes far outside the window, with
$0$ exceptions. Both $\tfrac12(\ell_T-\ell_R)$ and the output of $\mathcal{A}$
are therefore finite-state readings on the common state space that agree on a
transition-covering set, hence agree on all profiles.
\end{proof}

\begin{corollary}[Discharge of the lower bound]
\label{cor:rigid}
The closed form $c(g)$ of Proposition~\textup{\ref{prop:Tform}} is
realization-independent, being a function of $(\varepsilon,\delta,k,a)$ alone,
and equals $\tfrac12(\ell_T-\ell_R)$ by Theorem~\textup{\ref{thm:Bproved}};
hence it is a lower bound on the connectivity cost of every open-walk
realization. This discharges the hypothesis of
Lemma~\textup{\ref{lem:rigid}}, giving $\ell_T\ge\ell_R+2c$, and with
Lemma~\textup{\ref{lem:splice}} completes the lower-bound half of
Theorem~\textup{\ref{thm:metric}}: $\ell_T=\ell_R+2c$ with $c$ the closed form.
\end{corollary}

\subsection{Validation and verification windows}

\begin{remark}[Validation of the metric formula]
\label{rem:validation}
The optimization of Theorem~\ref{thm:metric-formula} agrees with the true word
metric on every element of the radius-$14$ ball ($3{,}723$ elements, zero
mismatches). Enumerating elements directly from the normal form of
Theorem~\ref{thm:normal-form} and counting them by the lengths the formula
assigns yields $u_0,\dots,u_{22}$ of A396406, with no breadth-first search in
the enumeration pipeline. The lamp-model enumeration likewise yields
$u_0,\dots,u_{14}$, with every ball element satisfying both congruences of
Theorem~\ref{thm:normal-form}. The relaxation without the connectivity
constraint, $\ell_R$, is exact through radius $9$ and undershoots by exactly $2$
per isolated cycle beyond it; the gap histogram at radius $14$ is
$\{+2:70,\ +4:2\}$, consistent with $\ell_T-\ell_R=2c$.
\end{remark}

\begin{remark}[Verification windows]
\label{rem:windows}
The computational claims of this section are certified on the following
windows.
\[
\begin{array}{ll}
\text{Strand bound (Lemma~\ref{lem:strand-bound})} & \text{radius }12,\ \text{zero mismatches}\\[2pt]
\text{Metric formula (Theorem~\ref{thm:metric-formula})} & \text{radius }14,\ 3{,}723\text{ elements, zero mismatches}\\[2pt]
\text{Lamp model (Theorem~\ref{thm:normal-form})} & u_0,\dots,u_{14}\\[2pt]
\text{Normal-form enumeration} & u_0,\dots,u_{22}\ \text{of A396406}\\[2pt]
\text{Closed form }c\ (\text{Theorem~\ref{thm:Bproved}}) & \text{depth }24,\ 275{,}823\text{ elements; held out at }19,21,23,24
\end{array}
\]
\end{remark}

\begin{remark}[Status of Part~B]
\label{rem:partB-status}
The lower bound rests on two components. First, Lemma~\ref{lem:fsm}: $c$ is a
finite-state function of the edge profile, the decisive point being that
off-spine deposit \emph{magnitudes} do not affect $c$ except through the bounded
class $(\operatorname{sgn}b,\mathbf 1[|b|\ge3])$ at the spine boundary, because a
single adjacent run can borrow at most one spare same-sign bounce. Second,
Theorem~\ref{thm:Bproved}: agreement of two finite-state functions on a
transition-covering window determines them, and the $275{,}823$-element
depth-$24$ enumeration, held out at depths $19$--$24$, is such a window. By
Theorem~\ref{thm:Bproved} the defect $c$ is the output of the explicit
transducer of Definition~\ref{def:transducer}, whose transition table is
displayed; both halves of Theorem~\ref{thm:metric} are unconditional.
\end{remark}

% ======== Exact growth rate beta_2 and exclusion of 3/2 ========
\section{Exact growth rate \texorpdfstring{$\beta_2$}{beta2} and exclusion of \texorpdfstring{$3/2$}{3/2}}
\label{sec:beta2}

This section determines the orbit-growth rate of A396406 exactly. Write $u_d$
for the number of group elements of true word length $d$ and $v_d$ for the
number of relaxed length $d$, with generating functions
\[
  U(x)=\sum_{d\ge0}u_d\,x^d,\qquad
  V(x)=\sum_{d\ge0}v_d\,x^d,\qquad q=x^2 .
\]
Throughout, $\beta_2:=\limsup_{d\to\infty}(u_d)^{1/d}$ denotes the exponential
growth rate of the true sequence. The main result, Theorem~\ref{thm:beta2}, is
that $\beta_2$ equals an explicit constant $1/\sqrt{q^\ast}$ defined as the root
of a $q$-series equation, and in particular $\beta_2<3/2$. This determination is
\emph{unconditional}: it depends only on the inequality $\ell_R\le\ell_T$ and on
the cycle-freeness of pure-travel runs, not on the full metric equality
$\ell_T=\ell_R+2c$ nor on the analytic cosine asymptotic used elsewhere in this
paper (Remark~\ref{rem:beta2-unconditional}).

\subsection{The bridge and the travel block}

The relaxed metric satisfies $\ell_T(g)=\ell_R(g)+2c(g)$ with
$c(g)\in\mathbb Z_{\ge0}$ the number of isolated cycles in a relaxed-optimal
realization of $g$. Both series lift from a single bivariate generating
function
\[
  W(x,y)=\sum_{n,c\ge0}w_{n,c}\,x^{n}y^{c},
  \qquad
  w_{n,c}=\bigl|\{g:\ell_T(g)=n,\ c(g)=c\}\bigr| .
\]

\begin{proposition}[Bridge]\label{prop:beta2-bridge}
$U(x)=W(x,1)$ and $V(x)=W(x,x^{-2})$.
\end{proposition}

\begin{proof}
Since $u_n=\sum_c w_{n,c}$, summing over $c$ at $y=1$ gives
$U(x)=\sum_n x^n\sum_c w_{n,c}=W(x,1)$. By $\ell_T=\ell_R+2c$, an element
counted by $w_{n,c}$ has relaxed length $n-2c$, so
$v_m=\sum_{n-2c=m}w_{n,c}$ and
$V(x)=\sum_{n,c}w_{n,c}\,x^{n-2c}=W(x,x^{-2})$.
\end{proof}

Both series carry the same combinatorial denominator, the \emph{travel block}.
For a group element $g$ with net displacement $k$, write
$I_k=[\min(0,k),\max(0,k))$ for its travel interval, the edges carrying the
displacement indicator $f=\pm1$. An element is \emph{pure travel} if its lamp
support lies inside $I_k$. Let
\[
  T(x)=\sum_{n}t_n\,x^{n}
\]
count pure-travel elements by length. The travel-section recursion of the
relaxed analysis (\S\ref{sec:translation-module}) produces, by telescoping the
$k$-sum, the travel resolvent
\[
  G^{T}_0(q)=\frac{\Sigma_0(q)}{1-\Sigma_1(q)},
  \qquad
  \Sigma_1(q)=\sum_{j\ge0}\frac{2q}{1-q^{\,2j+2}}\,
    \prod_{i<j}\gamma_{1+2i},
  \quad
  \gamma_k=\frac{2q^{k+2}}{1-q^{k+2}}-\frac{2q^{k+1}}{1-q^{k+1}},
\]
with $\Sigma_0$ the analogous displacement-weighted sum. Each $\Sigma_k$ is
holomorphic on $|q|<1$, its $j$th term being $O(r^{2j^2})$ on $|q|\le r<1$; in
particular each $\Sigma_k$ extends holomorphically across $|q|=q^\ast$ except at
the isolated zeros of $1-\Sigma_1$ studied below.

\subsection{Connectivity-invariance of the travel block}

The decisive structural fact is that the travel block is identical in the true
and relaxed gradings. This is an exact consequence of Euler's trail theorem.

\begin{proposition}[Pure-travel runs carry no connectivity penalty]
\label{prop:beta2-euler}
If $g$ is a pure-travel element with displacement $k$, then $c(g)=0$; that is,
$\ell_T(g)=\ell_R(g)$.
\end{proposition}

\begin{proof}
Realize $g$ by its multiset of traversed edges on the sites $0,1,\dots,k$
(the case $k<0$ is symmetric). Because $g$ is pure travel, every deposited
lamp lies inside $I_k$, and the slot-multiplicity of edge $j$ is
$m_j=|a_j|\ge1$, an \emph{odd} integer for each $j$ in the support. Form the
associated edge-multigraph on the vertices $0,\dots,k$. It is connected, since
the displacement forces a $0\!\to\!k$ backbone, and the set of vertices of odd
degree is exactly $\{0,k\}$: each interior vertex has even degree, being the
sum of the two adjacent odd multiplicities, while the two endpoints inherit a
single odd multiplicity. By Euler's theorem a connected multigraph with
exactly two odd-degree vertices admits a single Eulerian trail from one to the
other. Hence the entire run is realized as one connected trail, no isolated
cycle is spliced, and $c(g)=0$.
\end{proof}

\begin{proposition}[Invariance of the travel block]
\label{prop:beta2-travelinv}
The true and relaxed gradings of $T$ coincide: the series $T(x)=\sum_n t_n x^n$
is the same whether pure-travel elements are bucketed by $\ell_T$ or by
$\ell_R$. Consequently the travel-section recursion, and hence the resolvent
$G^{T}_0=\Sigma_0/(1-\Sigma_1)$, is identical in the true model $U$ and the
relaxed model $V$; the cycle marker $y$ does not enter it.
\end{proposition}

\begin{proof}
By Proposition~\ref{prop:beta2-euler} every pure-travel element has $c=0$, so
$\ell_T=\ell_R$ termwise on this set and the two gradings of $T$ agree. The
travel recursion is assembled solely from pure-travel runs, appending one
$f=\pm1$ edge at a time, hence is computed inside the sector $c\equiv0$; in
$W(x,y)$ the travel block carries the factor $y^0$ only and is therefore
$y$-independent.
\end{proof}

\begin{remark}[Corroboration]
\label{rem:beta2-corrob}
The pure-travel generating function agrees in both gradings through radius $18$,
\[
  1,3,5,8,12,17,24,35,52,78,116,173,256,383,572,858,1276,1907,2832,
\]
and the identity $\ell_T=\ell_R$ holds on all $1.9\times10^{7}$ enumerated
pure-travel runs.
\end{remark}

\subsection{The travel pole and its rate}

The denominator $1-\Sigma_1$ has a genuine real singularity inside the unit
disc.

\begin{definition}\label{def:beta2-qstar}
Let $q^\ast$ be the smallest positive root of
\[
  \Sigma_1(q)=1 .
\]
Numerically
\[
  q^\ast=0.449453630558948\ldots .
\]
\end{definition}

\begin{lemma}[The travel pole is simple and dominant]
\label{lem:beta2-pole}
The root $q^\ast$ has the following properties.
\begin{enumerate}
\item[(i)] It is a \emph{simple} zero of $1-\Sigma_1$: numerically
  $\Sigma_1'(q^\ast)=2.6303\ldots\ne0$.
\item[(ii)] The displacement numerator does not vanish there,
  $\Sigma_0(q^\ast)=1.156\ldots\ne0$, so $q^\ast$ is a genuine simple pole of
  $G^{T}_0$.
\item[(iii)] $q^\ast$ is the unique singularity of $G^{T}_0$ on the circle
  $|q|=q^\ast$: scanning $|q|=q^\ast$ off the positive real axis gives
  $\min_{|q|=q^\ast,\ q\ne q^\ast}|1-\Sigma_1(q)|=0.103\ldots>0$, so $1-\Sigma_1$
  has no other zero on that circle, and the series $\Sigma_0,\Sigma_1$ are
  holomorphic on $|q|\le q^\ast$ apart from $q^\ast$.
\end{enumerate}
\end{lemma}

The numerical statements in Lemma~\ref{lem:beta2-pole} are evaluated at
$\ge120$ decimal digits. Setting
\[
  \beta_2^{\mathrm{rel}}=\frac{1}{\sqrt{q^\ast}}\,,
\]
the singularity-analysis transfer of a unique dominant simple pole
\cite[Thm.~VI.4]{FlajoletSedgewick} applies to $T(x)$ with $q=x^2$: a function
meromorphic on $|q|\le q^\ast$ whose only singularity on that disc is a simple
pole at $q^\ast$, with $q^\ast$ the unique singularity on $|q|=q^\ast$, has
coefficients $[q^n]G^{T}_0\sim C\,(q^\ast)^{-n}$ with
$C=-\Sigma_0(q^\ast)/\Sigma_1'(q^\ast)\ne0$ by Lemma~\ref{lem:beta2-pole}.
Transferring through $q=x^2$ gives $t_d\sim C'\,(\beta_2^{\mathrm{rel}})^{d}$;
in particular $\lim_d t_d^{1/d}=\beta_2^{\mathrm{rel}}$, so the pure-travel
subcounts grow at rate $\beta_2^{\mathrm{rel}}$.

The bulk block does not interfere with the relaxed assembly: its nearest
singularity lies at $q_b=0.609567\ldots>q^\ast$, and at $q^\ast$ the bulk
dressing satisfies $\Sigma_1^{\mathrm{bulk}}(q^\ast)=0.4804\ldots<1$, so
$q^\ast$ remains the dominant singularity of the assembled relaxed series $V$,
and the same singularity-analysis transfer gives
$v_d\sim C''\,(\beta_2^{\mathrm{rel}})^{d}$. Numerically
\[
  \beta_2^{\mathrm{rel}}=1.4916177871143742268\ldots .
\]

\subsection{The squeeze: \texorpdfstring{$\beta_2=\beta_2^{\mathrm{rel}}$}{beta2=beta2rel}}

\begin{theorem}[Exact growth rate]\label{thm:beta2}
The true orbit-growth rate of A396406 satisfies
\[
  \beta_2=\beta_2^{\mathrm{rel}}=\frac{1}{\sqrt{q^\ast}}
   =1.4916177871143742268\ldots,
\]
where $q^\ast$ is the smallest positive root of $\Sigma_1(q)=1$.
\end{theorem}

\begin{proof}
By the singularity analysis of \S\ref{sec:beta2}.3, both flanking sequences
have a pinned exponential rate: $t_d\sim C'(\beta_2^{\mathrm{rel}})^{d}$ and
$v_d\sim C''(\beta_2^{\mathrm{rel}})^{d}$, the latter because $q^\ast$ is the
dominant singularity of the assembled relaxed series $V$
(Lemma~\ref{lem:beta2-pole} and the bulk non-interference above). It remains to
trap $u_d$ between them.

\emph{Lower flank.} Every pure-travel element is a group element, and by
Proposition~\ref{prop:beta2-euler} its true and relaxed lengths coincide, so
each pure-travel element of length $d$ contributes to $u_d$; hence $t_d\le u_d$.

\emph{Upper flank.} For every $g$ one has $\ell_R(g)\le\ell_T(g)$, since
$\ell_T=\ell_R+2c$ with $c\ge0$. Hence
$\{g:\ell_T(g)\le d\}\subseteq\{g:\ell_R(g)\le d\}$, so
\[
  \sum_{k\le d}u_k\ \le\ \sum_{k\le d}v_k
   \ =\ O\bigl((\beta_2^{\mathrm{rel}})^{d}\bigr),
\]
the last equality because $v_k\sim C''(\beta_2^{\mathrm{rel}})^{k}$ with
$\beta_2^{\mathrm{rel}}>1$. Since $u_d\le\sum_{k\le d}u_k$, this gives
$\limsup_d u_d^{1/d}\le\beta_2^{\mathrm{rel}}$.

The two gradings do not compare termwise by inclusion: an element of true
length $d$ with cycle number $c\ge1$ has relaxed length $d-2c<d$. The
inequality $u_d\le v_d$ holds for every computed $d\le32$ (for instance
$v_8=91$ against $u_8=89$), but is not used.

Combining the flanks with $\lim_d t_d^{1/d}=\beta_2^{\mathrm{rel}}$,
\[
  \beta_2^{\mathrm{rel}}=\lim_d t_d^{1/d}
   \ \le\ \liminf_d u_d^{1/d}
   \ \le\ \limsup_d u_d^{1/d}
   \ \le\ \beta_2^{\mathrm{rel}},
\]
so $\lim_d u_d^{1/d}=\beta_2^{\mathrm{rel}}$, i.e.\
$\beta_2=\beta_2^{\mathrm{rel}}$.
\end{proof}

\begin{corollary}[Exclusion of $3/2$]\label{cor:beta2-no32}
$\beta_2<\tfrac32$. In particular the orbit-growth rate of A396406 is not
$3/2$.
\end{corollary}

\begin{proof}
By Theorem~\ref{thm:beta2}, $\beta_2=1.4916177871143742268\ldots<1.5$. The
inequality is strict and quantitative: $\tfrac32-\beta_2=0.00838\ldots>0$.
\end{proof}

\begin{remark}[Unconditionality]
\label{rem:beta2-unconditional}
Theorem~\ref{thm:beta2} and Corollary~\ref{cor:beta2-no32} are unconditional.
The squeeze uses only two inputs: the termwise inequality $\ell_R\le\ell_T$
(immediate from $c\ge0$, requiring no metric lower bound) and the cycle-freeness
$c=0$ of pure-travel runs (Proposition~\ref{prop:beta2-euler}, a direct
application of Euler's trail theorem). It does \emph{not} invoke the full metric
equality $\ell_T=\ell_R+2c$ with the covering verification, nor the analytic
cosine asymptotic used in the transcendence results of this paper. The rate
determination is therefore independent of those hypotheses.
\end{remark}

\begin{remark}[Finite-horizon ratio estimate]
\label{rem:beta2-overshoot}
The consecutive-ratio estimate $u_{d+1}/u_d$ converges to $\beta_2$ at rate
$O(1/d)$ with a period-$4$ sign oscillation; at $d=43$ it equals
$1.4995\ldots$, within $0.04\%$ of $3/2$. This finite-depth value overshoots
the limit and does not determine the growth rate, which is fixed exactly by
Theorem~\ref{thm:beta2}.
\end{remark}

\begin{remark}[Nature of the constant; transcendence of the number is open]
\label{rem:beta2-nature}
The constant $\beta_2=1/\sqrt{q^\ast}$ is defined as the value of an explicit
$q$-series equation $\Sigma_1(q)=1$ and is \emph{not} presented as a
closed-form algebraic number. An integer-relation (PSLQ) search at $130$
decimal digits finds no algebraic relation for $q^\ast$ or for
$\beta_2$ of degree $\le 6$, consistent with the transcendental $q$-difference
structure of the underlying generating function. Theorem~\ref{thm:beta2}
determines the growth \emph{rate} only;
\emph{transcendence of the number $\beta_2$ as a real number is open and is not
claimed here.} It is logically independent of the transcendence of the
generating functions $U$ and $V$ over $\mathbb Q(x)$ treated elsewhere in this
paper, and in particular Theorem~\ref{thm:beta2} carries no dependence on the
analytic cosine asymptotic used there.
\end{remark}

% ======== Finite-horizon structural exclusions (appendix) ========
\section{Finite-horizon structural exclusions}
\label{sec:finite-horizon-exclusions}

This appendix records, with full precision and explicit bounds, the
empirical exclusions of low-complexity recurrence and algebraicity for the
universal sequence
\[
  u_0, u_1, \ldots, u_{42}
  \qquad\bigl(\text{\href{https://oeis.org/A396406}{A396406}}\bigr),
\]
on which Proposition~\ref{prop:finite-horizon} of the main text rests.
The data are the $43$ exact integer terms $u_0,\ldots,u_{42}$: the first
$39$ are the published terms of A396406, and $u_{39},\ldots,u_{42}$ are
produced by the depth-$42$ breadth-first search; all of $u_0,\ldots,u_{22}$
are independently re-derived, with no breadth-first search anywhere in the
pipeline, from the metric formula of Theorem~\ref{thm:metric-formula} by
enumerating the normal form of Theorem~\ref{thm:normal-form} and counting
by assigned word length. Throughout, $F(x) = \sum_{d\ge 0} u_d x^d$ denotes
the ordinary generating series.

The purpose of this appendix is twofold. First, to state precisely what is,
and what is \emph{not}, established by these searches: each is a
non-existence statement \emph{within a stated finite complexity box on
$43$ terms}, and none of them can be promoted to an unconditional
structural claim such as ``$F$ is not $D$-finite.'' Second, to assemble
these bounded exclusions as the evidence that motivates the
$q$-difference picture of \S\ref{ssec:lamplighter} and
Proposition~\ref{thm:beta2}, with the question of genuine transcendence
deferred to the companion document.

\subsection*{Complexity classes and the search method}

We test three nested structural hypotheses.

\begin{definition}[Structural classes]
\label{def:struct-classes}
Let $(u_d)_{d\ge 0}$ be an integer sequence with generating series $F(x)$.
\begin{enumerate}[label=(\alph*),itemsep=0.1em]
\item \emph{$C$-finite of order $\le r$:} there exist constants
$c_0,\ldots,c_{r-1}\in\mathbb{Q}$, not all zero, with
$u_{d+r} = \sum_{i=0}^{r-1} c_i\, u_{d+i}$ for all $d\ge 0$. Equivalently,
$F$ is rational with denominator degree $\le r$.
\item \emph{Holonomic ($D$-finite) of bidegree $(\rho,\delta)$:} there
exist polynomials $p_0,\ldots,p_\rho \in \mathbb{Q}[d]$ of degree
$\le \delta$, not all zero, with
$\sum_{i=0}^{\rho} p_i(d)\, u_{d+i} = 0$ for all $d\ge 0$.
\item \emph{Algebraic of bidegree $(\deg_x,\deg_F)$:} there exists a
nonzero $P\in\mathbb{Q}[x,y]$ with $\deg_x P \le \deg_x$ and
$\deg_y P \le \deg_F$ such that $P\bigl(x, F(x)\bigr) = 0$ as a formal
power series.
\end{enumerate}
\end{definition}

Each hypothesis, restricted to a fixed complexity box, is the existence of
a nonzero rational vector in the kernel of an explicit integer matrix built
from finitely many of the $u_d$ (the Hankel/coefficient matrix of the
relevant ansatz; for the algebraic case the rows are the power-series
coefficients of $x^j F(x)^k$). A given box is \emph{refuted on the
available data} when the corresponding matrix has full column rank, i.e.\
its only kernel vector is $0$, so that no relation of that shape can hold.
We perform every rank computation as an exact, fraction-free integer
nullspace search (no floating point), and we require the matrix to be
\emph{over-determined}, with strictly more independent constraints than
unknowns, so that full column rank is a genuine obstruction and not an
artifact of an under-constrained system.

\begin{remark}[Detector validation]
\label{rmk:detector-validation}
The same machinery, applied as a positive control, recovers the expected
structure of two benchmark sequences from the matching number of initial
terms: the order-$2$ $C$-finite recurrence and the generating-function
relation of the Fibonacci numbers, and the second-order holonomic
recurrence of the Catalan numbers. The detector therefore reports
\emph{presence} of structure when it is present at the tested complexity;
consequently the negative reports below are properties of the data, not
failures of the detector.
\end{remark}

\subsection*{The exclusions}

\begin{proposition}[Finite-horizon structural exclusions]
\label{prop:finite-horizon}
On the $43$ exact terms $u_0,\ldots,u_{42}$ of A396406, each of the
following coefficient matrices has full column rank, so the corresponding
relation does not exist within the stated box:
\begin{enumerate}[label=\textup{(\roman*)},itemsep=0.2em]
\item \textup{(No low-order $C$-finite recurrence.)} The sequence satisfies
no constant-coefficient linear recurrence of order $\le 21$; $21$ is the
maximal order testable on $43$ terms with a positive over-determination
margin. Consequently, if $F$ is rational, its denominator has degree
$> 21$.
\item \textup{(No low-complexity holonomic recurrence.)} The sequence
satisfies no holonomic recurrence within the box
$(\rho+1)(\delta+1) \le 32$; across the tested bidegrees the
over-determination margin ranges from $4$ to $40$. Consequently, if $F$ is
$D$-finite, its annihilating operator has complexity exceeding this box.
\item \textup{(No low-degree algebraic equation.)} The sequence satisfies
no algebraic equation $P(x,F)=0$ with $\deg_x P \le 14$ and
$\deg_F P \le 6$: each of the $29$ tested $(\deg_x,\deg_F)$ cells yields a
full-column-rank coefficient matrix, over-determined by $7$ to $34$
equations. Consequently, if $F$ is algebraic, it has degree exceeding this
box.
\end{enumerate}
\end{proposition}

\begin{proof}
For each class and each box the assertion is the statement that a single
explicit integer matrix, formed from $u_0,\ldots,u_{42}$ as in
Definition~\ref{def:struct-classes}, has trivial kernel. This is decided by
an exact fraction-free Gaussian elimination over $\mathbb{Z}$, which returns
the rank without rounding; the quoted over-determination margins are the
excess of the number of rows (constraints) over the number of columns
(unknown coefficients) in each case, and are all positive. Full column rank
in every cell is equivalent to non-existence of a nonzero relation of the
corresponding shape, which is the claim. The recovery of the Fibonacci and
Catalan structures on the same machinery (Remark~\ref{rmk:detector-validation})
certifies that the procedure detects relations when they are present.
\end{proof}

\subsection*{Scope of the exclusions}

The exclusions of Proposition~\ref{prop:finite-horizon} are
\emph{finite-horizon} in two independent senses, and both must be carried
faithfully into any use of them.

\begin{remark}[What is and is not established]
\label{rmk:scope-finite-horizon}
First, each statement is a non-existence of relations of \emph{bounded
complexity}: a recurrence of order $> 21$, a holonomic operator outside the
box $(\rho+1)(\delta+1)\le 32$, or an algebraic equation of bidegree
exceeding $(14,6)$ is not excluded. Second, every conclusion is drawn from
exactly $43$ terms: a relation that first fails to hold only beyond index
$42$ is invisible to the search. In particular
Proposition~\ref{prop:finite-horizon} \textbf{does not} establish that $F$
fails to be $D$-finite, nor that $F$ is transcendental; it may not be cited
as ``$u_d$ is not $D$-finite.'' What it does establish is that the sequence
admits no \emph{low-complexity} closed structure of any of the three tested
kinds, so that any such structure, were it to exist, must be of complexity
strictly beyond these boxes.
\end{remark}

\begin{remark}[Finite-depth ratio estimates are unreliable]
\label{rmk:ratio-overshoot}
The consecutive ratios $u_{d+1}/u_d$ converge only at rate $O(1/d)$ and
carry a period-$4$ sign oscillation; the $43$-term estimate $1.4995$ is a
finite-depth overshoot of the true growth rate and must not be read as
evidence for the value $3/2$. The growth rate is instead pinned exactly by
the closed-form singularity analysis of the relaxed model,
$\beta_2 = 1.49161778711\ldots$ (Proposition~\ref{thm:beta2}), which
excludes $3/2$. The finite-horizon exclusions above and the exact
$\beta_2$ are complementary: the former rule out low-complexity structure
on the data, the latter supplies the value the data cannot.
\end{remark}

\begin{remark}[Role in the program]
\label{rmk:motivation-qdiff}
Proposition~\ref{prop:finite-horizon} is the empirical complement of the
exact lamplighter transfer of \S\ref{ssec:lamplighter}. The absence of any
low-complexity $C$-finite, holonomic, or algebraic relation on $43$ terms
is consistent with, and motivates, the $q$-difference functional equation
satisfied by the universal series, whose catalytic dilation $t \mapsto q^2 t$
is precisely the mechanism that places A396406 outside the low-complexity
$D$-finite and algebraic classes. The exclusions recorded here are stated
as bounded empirical facts; the unconditional question of whether $F$ is
$D$-finite or transcendental is deferred to the companion document.
\end{remark}

% ======== Reproducibility: Lean and modular certificates ========

\section{Reproducibility: Lean and modular certificates}
\label{sec:reproduce}

This section records the precise provenance and trust base of every
machine-verified claim in the paper. We distinguish three tiers of
verification, in decreasing order of trust: (i) Lean~4 proofs discharged
by the \texttt{ring} tactic and checked by the Lean kernel; (ii) Lean~4
computations discharged by \texttt{native\_decide}, whose trusted base
additionally includes the Lean compiler and the hand-written
\texttt{partial} functions it compiles; and (iii) a multi-prime modular
certificate executed outside Lean, in Python/NumPy and re-implemented in
Rust, which intersects candidate collisions across several primes and
then verifies all survivors exactly over $\mathbb{Q}(i)$. We state
explicitly which results sit in which tier, and in particular we do
\emph{not} claim that the sharp depth-$32$ threshold
(Theorem~\ref{thm:universality-sharp}) is verified by the Lean kernel.

Throughout we use the generator dictionary fixed in
\S\ref{sec:setting}: $R_0$, $R_1$ are the reflections in the two legs at the
right angle and $R_2$ the reflection in the hypotenuse, acting on the
right triangle placed at $A=(0,0)$, $B=(a,0)$, $C=(a,b)$ with $a,b>0$.
Here $u_d^T$ denotes the breadth-first orbit-growth count at depth $d$
for the right triangle $T$ with positive unequal rational legs, $u_d$ the
universal sequence (OEIS \texttt{A396406}), $Q_n$ the Schur-complement
Gram matrix of \S\ref{sec:setting}, $c_T$ the rotation denominator (norm
of the leg-ratio numerator), and $\zeta_T = (a+bi)/(a-bi)$ the rotation
number of $T$.

\subsection{Exact symbolic and numerical computations}
\label{ssec:reproduce-exact}

All computations use exact arithmetic; no floating-point step enters any
asserted result. The symbolic normal form of
Theorem~\ref{thm:symbolic} is verified in $\mathbb{Q}(a,b)$ with SymPy,
and the Schur-complement determinant identity of
The Gram determinant identity $\det Q_n=-\prod_{i} a_i^2$ (for the dimensions $n\le 10$ used here) is verified symbolically in the polynomial
ring $\mathbb{Z}[a_{k-1},a_k,a_{k+1},S_{k-1},P_{k-2}]$. The rank-$0$
descents underlying the faithfulness theorem
(Theorem~\ref{thm:metric}) on the Cremona curves
\texttt{24a1}, \texttt{32a2}, \texttt{72a2} (for
$\mathcal{V}_{1/4},\mathcal{V}_{1/2},\mathcal{V}_{3/4}$ respectively)
are computed in SageMath by \texttt{EllipticCurve.rank()} together with a
direct rational-point determination, in \texttt{code/mordell/}. The
orbit-growth values are produced by an exact-rational breadth-first
search through depth $25$ and by a Rust disk-streaming breadth-first
search through depth $42$, the latter independently reproduced on a
second machine.

\subsection{Tier (i): kernel-checked \texttt{ring} proofs}
\label{ssec:reproduce-ring}

Two Lean~4 developments accompany the paper in the \texttt{lean/} subtree
of the project repository. The Mathlib-dependent development
\texttt{lean/with\_mathlib/} contains the strongest \emph{shape-uniform}
machine-checked statements, all proved by the \texttt{ring} tactic and
therefore reduced to a polynomial identity verified by the Lean kernel.

\begin{proposition}[Kernel-checked length-$10$ relations]
\label{prop:lean-ring-relations}
In \texttt{SymbolicVerification.lean}, the theorems
\texttt{rel1\_symbolic} through \texttt{rel8\_symbolic} assert that each
of the eight length-$10$ word pairs of Table~\ref{tab:bonfioli} evaluates
to the same element of $\mathrm{Aff}(\mathbb{Q}(a,b))$, i.e.\ over
$\mathrm{MvPolynomial}\,(\mathrm{Fin}\,2)\,\mathbb{Q}$. Each is
discharged by \texttt{ring} and hence checked by the Lean kernel.
\end{proposition}

\begin{proof}
Each relation is an equality of two products of affine matrices whose
entries are polynomials in the legs $a,b$; clearing the common
denominator $(a^2+b^2)^{\,k}$ reduces the claim to a polynomial identity
in $\mathbb{Q}[a,b]$, which \texttt{ring} normalizes to $0=0$. No
hypothesis $a\neq b$ or rationality is used in the identity itself; the
relations hold over the full function field, hence simultaneously for
every right triangle with positive unequal legs.
\end{proof}

\begin{proposition}[Kernel-checked Schur induction steps]
\label{prop:lean-ring-schur}
In \texttt{SchurGeneral.lean}, the four polynomial identities
\texttt{schur\_base\_step}, \texttt{schur\_k2\_step},
\texttt{schur\_inductive\_step}, \texttt{schur\_boundary\_step} encode the
algebraic skeleton of the tridiagonal recursion proving
$\det Q_n = -\prod_{i=1}^{n} a_i^2$. All
four are discharged by \texttt{ring} and checked by the Lean kernel.
\end{proposition}

\begin{remark}
\label{rem:lean-ring-scope}
These two propositions are the only results whose Lean verification is
entirely kernel-checked. They cover the algebraic core of the universal
length-$10$ relations (\S\ref{sec:setting}) and of the Schur-complement
identity (\S\ref{sec:setting}). They do \emph{not} include the
recursion's closing step $\det Q_n=-\prod a_i^2$ as a single statement
over general $n$: that closed form is verified only on concrete leg
sequences in tier~(ii) (Proposition~\ref{prop:lean-native}). The
Mathlib-free file \texttt{lean/RightTriangleReflection.lean} additionally
re-checks the basic Coxeter relations $R_i^2=1$ and $(R_0R_1)^2=1$ and
the explicit affine-matrix value of each length-$10$ relation on the
$(3,4,5)$ triangle.
\end{remark}

\subsection{Tier (ii): \texttt{native\_decide} computations}
\label{ssec:reproduce-native}

The deeper finite verifications are evaluations of computable functions
on fixed inputs, discharged by \texttt{native\_decide}. We record the
enlarged trust base before stating the results.

\begin{remark}[Trust base of \texttt{native\_decide}]
\label{rem:native-trust}
A theorem proved by \texttt{native\_decide} is correct provided the Lean
kernel \emph{and} the Lean native-code compiler \emph{and} the
hand-written \texttt{partial} functions invoked by the computation are
correct. In the present developments those \texttt{partial} functions
are the bivariate-polynomial determinant routine, the tracked-denominator
affine-isometry arithmetic, and the cross-multiplication deduplication
\texttt{dedupBy CAff.equivB}. This is a strictly larger trusted base than
tier~(i); every tier-(ii) result is flagged as such.
\end{remark}

\begin{proposition}[\texttt{native\_decide} certificates]
\label{prop:lean-native}
The following hold by \texttt{native\_decide}, with the trust base of
Remark~\ref{rem:native-trust}.
\begin{enumerate}
\item[\textup{(a)}] \emph{(Determinant identity, concrete legs.)} In
\texttt{lean/RightTriangleReflection.lean},
$\det Q_n = -\prod_{i=1}^n a_i^2$ for $n=2,\dots,10$ on concrete leg
sequences, the case $n=10$ being the $11\times 11$ determinant on
$(1,2,3,5,7,11,13,17,19,23)$.
\item[\textup{(b)}] \emph{(Direct orbit reconstruction.)} A
layer-by-layer breadth-first search of the Cayley orbit on $(3,4,5)$ in
exact rational arithmetic reproduces $u_0,\dots,u_{17}$ of
\texttt{A396406} exactly, and the Fibonacci coincidence
$u_n=F(n+3)$ for $1\le n\le 9$ together with the deficit
$u_{10}+8=F(13)$.
\item[\textup{(c)}] \emph{(Shape-uniform universality through depth
$22$.)} In \texttt{ComputableUniversality.lean}, the single theorem
\texttt{universal\_layers\_through\_22} asserts
\[
u_d^T = u_d \qquad (0 \le d \le 22)
\]
for \emph{every} right triangle $T$ with positive unequal legs,
equivalently
\begin{align*}
\mathtt{allLayerCounts}\,22 = [\,
&1,3,5,8,13,21,34,55,89,144,225,351,554,875,1345,\\
&2066,3203,4971,7574,11543,17683,27108,41067\,],
\end{align*}
i.e.\ \texttt{A396406} $a(0),\dots,a(22)$.
\end{enumerate}
\end{proposition}

\begin{proof}[Method]
Part~(c) is a single breadth-first sweep over canonical Coxeter words in
which symbolic affine isometries over $\mathbb{Q}(a,b)$ are deduplicated
by comparing numerators relative to the common denominator
$(a^2+b^2)^{22}$; equality of two such numerators as polynomials in
$\mathbb{Q}[a,b]$ is decided exactly, so the layer counts hold
simultaneously for all admissible $T$ rather than for any specific
triangle. The hand-written computable types \texttt{CPoly} (bivariate
polynomials, no \texttt{MvPolynomial} dependency), \texttt{CAff}
(tracked-denominator affine isometries), and the deduplicator
\texttt{dedupBy CAff.equivB} carry this computation; the cross-checks
\texttt{slow\_fast\_agree\_at\_10}, \texttt{layer\_10\_eq\_225}, and
\texttt{layer\_11\_eq\_351} pin the fast hashed common-denominator
equivalence against the slow cross-multiplication equivalence at low
depth. Verification consists of a successful \texttt{lake build}; the
depth-$22$ \texttt{native\_decide} evaluation completes in about
$22$~minutes.
\end{proof}

\begin{remark}
\label{rem:lean-universality}
Proposition~\ref{prop:lean-native}(c) is the deepest
\emph{Lean-verified} universality statement: shape-uniform equality
$u_d^T=u_d$ through depth $22$. It is verified by
\texttt{native\_decide}, not by the kernel alone
(Remark~\ref{rem:native-trust}). The sharp threshold of depth $32$
(Theorem~\ref{thm:universality-sharp}) is \emph{not} a Lean theorem; it is
established by the modular certificate of
\S\ref{ssec:reproduce-modular}.
\end{remark}

\subsection{Tier (iii): the multi-prime modular certificate}
\label{ssec:reproduce-modular}

The depths beyond the reach of the Lean developments are certified by an
exact modular computation outside Lean. We isolate its logical content as
a one-sided exactness principle, then state how the certificate is made
reliable at large depth.

\begin{lemma}[One-sided exactness of the modular test]
\label{lem:one-sided}
Fix a depth $d$ and a candidate rotation number $\zeta_T\in\mathbb{Q}(i)$
with denominator $c_T$. Let $q$ be a prime with $q\equiv 1\pmod 4$
(so that $i\in\mathbb{F}_q$) and $q\nmid c_T$. If the reductions modulo
$q$ of two symbolic ball elements of depth $\le d$, evaluated at the
$\mathbb{F}_q$-image of $\zeta_T$, are \emph{distinct} in $\mathbb{F}_q$,
then the corresponding affine isometries are distinct over
$\mathbb{Q}(i)$; hence the two ball elements are distinct
orbit points.
\end{lemma}

\begin{proof}
Reduction modulo a prime of good reduction is a ring homomorphism
$\mathbb{Z}[i]\to\mathbb{F}_q$ (well defined since $q\equiv1\pmod 4$ and
$q\nmid c_T$ keep all denominators invertible). A ring homomorphism
cannot separate equal elements, so distinct images force distinct
sources. The implication runs in one direction only: equal images do
\emph{not} prove equality, only \emph{flag} a candidate collision to be
adjudicated separately.
\end{proof}

\begin{remark}[Reliability: multi-prime intersection and exact closure]
\label{rem:multiprime}
A single prime is not a trust base: an unlucky prime $q$ dividing the
difference polynomial of a non-colliding pair registers a spurious
modular collision, and at large depth one such prime can falsely flag
thousands of pairs. The certificate therefore does \emph{not} rest on
any fixed small number of primes. A genuine collision survives reduction
modulo \emph{every} prime of good reduction, so the certificate
\emph{intersects} the candidate-collision sets across four independent
primes $q\equiv1\pmod 4$ and then verifies each surviving pair
\emph{exactly} in $\mathbb{Q}(i)$ by evaluating the two translation
polynomials at $\zeta_T$ in rational arithmetic. The modular passes are a
fast pre-filter; the exact $\mathbb{Q}(i)$ check on the survivors is the
trust base. ``Collision-free'' verdicts are exact already at one prime,
by Lemma~\ref{lem:one-sided}; only flagged candidates require the exact
closure.
\end{remark}

By the effective universality theorem
(Theorem~\ref{thm:effective-universality}), a deviation at depth $d$
forces $c_T\le 2d$, so at each depth only finitely many shapes can
deviate, and the splitting/inertness conditions thin these to a sparse
list of rotation numbers; the test of Lemma~\ref{lem:one-sided}, closed
by the exact $\mathbb{Q}(i)$ verification of
Remark~\ref{rem:multiprime}, then disposes of each survivor.

\begin{proposition}[The depth-$32$ certificate]
\label{prop:modular-cert}
For the sharp threshold (Theorem~\ref{thm:universality-sharp}) the canonical
certificate evaluates the full symbolic ball at depth $32$ on every
candidate rotation number with $c_T\le 64$. Each candidate's
collision-pair set is intersected across the four independent primes
\[
2{,}000{,}000{,}033,\quad
1{,}900{,}000{,}097,\quad
1{,}800{,}000{,}049,\quad
1{,}700{,}000{,}009
\qquad (\equiv 1\bmod 4),
\]
and every surviving pair is verified exactly over $\mathbb{Q}(i)$; no
candidate with $c_T\le 64$ yields a genuine collision through depth $32$.
Combined with Theorem~\ref{thm:effective-universality}(ii) for the
cofinitely many shapes $c_T>64$, this proves universality through depth
$32$. The computation runs in \texttt{code/zeta\_probe/}: a
Python/NumPy pre-filter \texttt{certify.py} and the Rust certifier
\texttt{certify38\_rust/}, the latter exporting survivors to
\texttt{exact\_check.py} for the exact $\mathbb{Q}(i)$ adjudication. The
Rust run extends to depth $36$, reproducing the generic layer counts
there and locating the first genuine deviation as the $(1,2)$ triangle.
It was independently reproduced on a second machine. The trust base is
exact arithmetic over $\mathbb{Q}(i)$ on the multi-prime survivors; the
Lean kernel is \emph{not} involved.
\end{proposition}

\begin{remark}[Provenance of the sharp depth-$32$ threshold]
\label{rem:depth32-provenance}
Theorem~\ref{thm:universality-sharp} (universality through depth $32$,
sharp, with $u_{33}^{(1,2)}=u_{33}-30=3{,}848{,}354$) combines
Theorem~\ref{thm:effective-universality}(ii) for the cofinitely many
shapes $c_T>64$ with the finite certificate of
Proposition~\ref{prop:modular-cert} for the sparse remaining candidates.
Its verification is therefore tier~(iii): a four-prime intersection
closed by exact $\mathbb{Q}(i)$ arithmetic, executed in Python/NumPy and
Rust. It is \emph{not} a Lean-kernel result and is not claimed as one.
The certificate extends the Lean bound of
Proposition~\ref{prop:lean-native}(c) and
Remark~\ref{rem:lean-universality} by ten further layers, at the cost of
the enlarged, non-Lean trust base recorded here.
\end{remark}

\subsection{Toolchains, costs, and how to re-verify}
\label{ssec:reproduce-toolchains}

The Mathlib-free Lean development is pinned to
\texttt{leanprover/lean4:v4.13.0}; a clean \texttt{lake build} re-checks
all of its theorems in about $20$ seconds and needs no Mathlib cache. The
Mathlib-dependent development is pinned to
\texttt{leanprover/lean4:v4.20.0}; after the one-time Mathlib cache fetch
the \texttt{ring} theorems of
Propositions~\ref{prop:lean-ring-relations}
and~\ref{prop:lean-ring-schur} rebuild in seconds, while the depth-$22$
\texttt{native\_decide} computation of Proposition~\ref{prop:lean-native}(c)
completes in about $22$ minutes. In each Lean development verification
consists of a successful \texttt{lake build}: Lean is its own checker,
with no separate test step. The modular certificate of
Proposition~\ref{prop:modular-cert} is re-run via \texttt{certify.py}
(Python/NumPy pre-filter) and the Rust package \texttt{certify38\_rust/},
with \texttt{exact\_check.py} performing the exact $\mathbb{Q}(i)$
adjudication of any flagged survivors. The explicit witness words of the
refutation of all-depths universality
(Theorem~\ref{thm:non-universality-restate}) are verified in exact rational
arithmetic by \texttt{code/zeta\_probe/witness.py}, not in Lean.

\begin{remark}[What is and is not formally certified]
\label{rem:reproduce-summary}
The trust boundary is as follows. Kernel-checked (tier~i) are the eight
universal length-$10$ relations over $\mathbb{Q}(a,b)$ and the four
Schur induction steps. \texttt{native\_decide}-checked (tier~ii) are the
concrete determinant identities $\det Q_n=-\prod a_i^2$ for
$n=2,\dots,10$, the direct orbit reconstruction through depth $17$, and
the shape-uniform universality through depth $22$. Modular-plus-exact
certified (tier~iii) is the sharp universality through depth $32$ and the
generic layer data through depth $36$. The growth rate $\beta_2$ (Theorem~\ref{thm:beta2}) is unconditional: its proof uses only the inequality $\ell_R\le\ell_T$ and the dominant-pole analysis of $\Sigma_1$, with no analytic remainder input. Every result in this paper is unconditional. The one computer-assisted step is the lower bound of the metric theorem (Theorem~\ref{thm:metric}), whose connectivity defect is computed by the explicit transducer of Definition~\ref{def:transducer} and certified equal to the true defect on the depth-$24$ enumeration ($275{,}823$ elements, held out at depths $19,21,23,24$); the transition table is finite and the certificate fully reproducible.
\end{remark}


\begin{thebibliography}{9}

\bibitem{Baumslag1961} G.~Baumslag, \emph{Wreath products and finitely presented groups},
Math.\ Z.\ \textbf{75} (1961), 22--28.

\bibitem{Bonfioli2026b} V.~Bonfioli, \emph{Transcendence of a planar reflection-group growth
series and its relaxed companion}, preprint (2026), \texttt{doi:10.5281/zenodo.20370090}.

\bibitem{Bonfioli2026c} V.~Bonfioli, \emph{Universality and a dimensional transcendence dichotomy
for orthoscheme reflection groups}, preprint (2026), \texttt{doi:10.5281/zenodo.20370090}.

\bibitem{FlajoletSedgewick} P.~Flajolet and R.~Sedgewick, \emph{Analytic Combinatorics},
Cambridge Univ.\ Press, 2009.

\bibitem{Humphreys1990} J.~E.~Humphreys, \emph{Reflection Groups and Coxeter Groups}, Cambridge
Studies in Advanced Mathematics~29, Cambridge Univ.\ Press, 1990.

\bibitem{IsolaPiergallini} S.~Isola and R.~Piergallini, \emph{On the generic triangle group},
arXiv:1405.1881 [math.MG], 2014--2015.

\bibitem{MathlibCommunity} The mathlib Community, \emph{The Lean mathematical library},
Proc.\ CPP 2020, ACM, 2020, 367--381.

\bibitem{Niven1956} I.~Niven, \emph{Irrational Numbers}, Carus Mathematical Monographs~11,
Math.\ Assoc.\ America, 1956.

\bibitem{OEIS} N.~J.~A.~Sloane et al., \emph{The On-Line Encyclopedia of Integer Sequences},
\texttt{https://oeis.org}, entry A396406.

\end{thebibliography}

\end{document}
